%===========================================================
%  pandoc-template.tex — DollarHua PDF Toolbook Template
%  Version: v2.0 · Updated: 2026-09-07
%
%  Used with: pandoc merged.md --template=this -o output.tex
%  Compile with: xelatex (2 passes: TOC + cross-refs; optional 3rd if "Rerun")
%
%  Optimizations over v1.3:
%    • Cover rendered as literal titlepage (v1.3 \@maketitle was shadowed by ctex)
%    • Chapter: display style with 0.6×\textwidth decorative band + 章号居中
%    • Header: chapter name (left) + book title (right), Bond Navy on blossom rule
%    • Footer: page number centered
%    • Part divider macro: \partdivider{...} for "第N部分" inline banners
%    • Tighter line spacing for body, generous for chapter intros
%    • First-line indent 2em CJK, parskip 0.4em
%===========================================================

% ── Document class (ctexbook = Chinese-enabled book) ──
% fontset=none → ctex doesn't preload any CJK font; we set our own below.
\documentclass[a4paper,11pt,oneside,fontset=none,UTF8,zihao=-4]{ctexbook}

% ── Page geometry ──
% headheight set to 24pt to satisfy fancyhdr and prevent header clipping
\usepackage[top=2.6cm, bottom=2.6cm, left=2.8cm, right=2.4cm,
            headheight=24pt, headsep=8mm, footskip=12mm,
            includeheadfoot]{geometry}

% ── DollarHua brand colors ──
\usepackage{xcolor}
\definecolor{dh-pink}{HTML}{FEC6CD}
\definecolor{dh-blossom}{HTML}{F3A6AF}
\definecolor{dh-gold}{HTML}{E3A04B}
\definecolor{dh-navy}{HTML}{1D1E50}
\definecolor{dh-bronze}{HTML}{B86C40}
\definecolor{dh-skin}{HTML}{F9DECE}
\definecolor{dh-white}{HTML}{FCFDFD}
\definecolor{dh-red}{HTML}{E85D61}
\definecolor{dh-navy-muted}{HTML}{8E8FA2}

% ── CJK typography (cross-platform auto-detect) ──
\usepackage{fontspec}
\usepackage{xeCJK}

\IfFontExistsTF{PingFang SC}{
  \setCJKmainfont{PingFang SC}[
    BoldFont={PingFang SC Semibold},
    ItalicFont={PingFang SC Light}
  ]
  \setCJKsansfont{PingFang SC}
  \setCJKmonofont{PingFang SC}
}{
  \IfFontExistsTF{Source Han Serif CN}{
    \setCJKmainfont{Source Han Serif CN}[
      BoldFont={Source Han Serif CN Bold},
      AutoFakeSlant=0.15
    ]
    \setCJKsansfont{Source Han Sans CN}[
      BoldFont={Source Han Sans CN Bold}
    ]
    \setCJKmonofont{Source Han Sans CN}
  }{
    \IfFontExistsTF{Noto Serif CJK SC}{
      \setCJKmainfont{Noto Serif CJK SC}[
        BoldFont={Noto Serif CJK SC Bold},
        AutoFakeSlant=0.15
      ]
      \setCJKsansfont{Noto Sans CJK SC}
      \setCJKmonofont{Noto Sans Mono CJK SC}
    }{
      \IfFontExistsTF{Source Han Serif SC}{
        \setCJKmainfont{Source Han Serif SC}[
          BoldFont={Source Han Serif SC Bold},
          AutoFakeSlant=0.15
        ]
        \setCJKsansfont{Source Han Sans SC}
        \setCJKmonofont{Source Han Sans SC}
      }{
        % Ultimate fallback: SimSun on Windows
        \setCJKmainfont{SimSun}[AutoFakeBold=2.5]
        \setCJKsansfont{SimHei}
        \setCJKmonofont{KaiTi}
      }
    }
  }
}

% CJK line breaking + spacing
\XeTeXlinebreaklocale "zh"
\XeTeXlinebreakskip = 0pt plus 1pt
\xeCJKsetup{CJKecglue=\hskip 0.15em plus 0.05em minus 0.05em}

% Line spreading for CJK body text readability
\linespread{1.45}\selectfont

% ── Unicode character support ──
\usepackage{newunicodechar}
\newunicodechar{≈}{\ensuremath{\approx}}
\newunicodechar{≥}{\ensuremath{\geq}}
\newunicodechar{≤}{\ensuremath{\leq}}
\newunicodechar{×}{\ensuremath{\times}}
\newunicodechar{—}{\textemdash}
\newunicodechar{–}{\textendash}
\newunicodechar{…}{\dots}
\newunicodechar{°}{\ensuremath{^{\circ}}}
\newunicodechar{⚠}{\textcolor{dh-bronze}{\textbf{[注意]}}}
\newunicodechar{️}{} % Ignore variation selector-16

% ── Page style (headers / footers) ──
\usepackage{fancyhdr}
\pagestyle{fancy}
\fancyhf{}
\renewcommand{\headrulewidth}{0.5pt}
\renewcommand{\headrule}{\hbox to\headwidth{%
  \color{dh-blossom}\leaders\hrule height \headrulewidth\hfill}}
% In oneside mode: left header shows current chapter, right header shows book title
\fancyhead[L]{\footnotesize\color{dh-navy}\leftmark}
\fancyhead[R]{\footnotesize\color{dh-navy}\mytitle}
\newcommand{\footerdisclaimer}{{\tiny\color{dh-navy-muted}投资有风险，入市需谨慎}}
\fancyfoot[C]{\thepage\\[1pt]\footerdisclaimer}
\fancypagestyle{plain}{
  \fancyhf{}
  \fancyfoot[C]{\thepage\\[1pt]\footerdisclaimer}
  \renewcommand{\headrulewidth}{0pt}
  \renewcommand{\headrule}{}
}
% Chapter starts: show only page number and disclaimer, no header
\fancypagestyle{chapterstart}{
  \fancyhf{}
  \fancyfoot[C]{\thepage\\[1pt]\footerdisclaimer}
  \renewcommand{\headrulewidth}{0pt}
  \renewcommand{\headrule}{}
}

% ── Paragraph: first-line indent 2 CJK chars, generous inter-paragraph ──
\setlength{\parindent}{2em}
\setlength{\parskip}{0.4em plus 0.1em minus 0.1em}
\setlength{\emergencystretch}{1.5em}
\clubpenalty=10000
\widowpenalty=10000
\displaywidowpenalty=10000

% ── Chapter/section title styling ──
\usepackage{titlesec}

% Mixed-color helpers (avoid titlesec choking on `!NN` mixing syntax)
\definecolor{dh-navy-soft}{HTML}{3A3B72}    % 80% of dh-navy
\definecolor{dh-navy-faint}{HTML}{8485A0}   % 50% of dh-navy

% Chapter: display style with a band on top, the "第N章" label, the
% chapter title, and a trailing rule.
%
% titlesec's display-style "label" slot is sensitive: it must be a
% single, balanced group (no bare \par, no unprotected \vspace).
% We therefore wrap the entire label in a single {…} and use
% \vskip / \hrulefill instead of \par / \rule.
\titleformat{\chapter}[display]
  {\centering\Huge\bfseries\sffamily\color{dh-navy}}
  {{\color{dh-blossom}\hrule height 1.5pt}\vskip 0.7cm
   {\color{dh-navy-soft}\Large 第\ \thechapter\ 章}}
  {1.0em}
  {}
  [\vskip 0.6cm{\color{dh-blossom}\hrule height 1pt}]

\titlespacing*{\chapter}{0pt}{50pt}{40pt}

% Section: large bold with blossom number prefix (each section starts on a new page)
\newcommand{\sectionbreak}{\clearpage}
\titleformat{\section}
  {\Large\bfseries\sffamily\color{dh-navy}}
  {\color{dh-blossom}\thesection}
  {0.8em}
  {}
  [\vspace{-0.2em}{\color{dh-blossom}\titlerule[0.8pt]}]

\titlespacing*{\section}{0pt}{2.0em}{0.8em}

% Subsection: medium bold
\titleformat{\subsection}
  {\large\bfseries\sffamily\color{dh-navy}}
  {\color{dh-blossom}\thesubsection}
  {0.6em}
  {}

\titlespacing*{\subsection}{0pt}{1.6em}{0.6em}

% Subsubsection: small bold
\titleformat{\subsubsection}
  {\normalsize\bfseries\sffamily\color{dh-navy}}
  {\color{dh-blossom}\thesubsubsection}
  {0.5em}
  {}

% Subsection / subsubsection numbering: drop the \thesection component
% so a chapter with no parent \section does not render as "1.0.1, 1.0.2".
% We use plain chapter.X (and chapter.X.Y) instead.
\renewcommand{\thesubsection}{\thechapter.\arabic{subsection}}
\renewcommand{\thesubsubsection}{\thechapter.\arabic{subsection}.\arabic{subsubsection}}

% Make sure chapter shows in header (with short title)
\renewcommand{\chaptermark}[1]{\markboth{#1}{}}

% ── TOC styling ──
\usepackage{tocloft}
\renewcommand{\cftchapfont}{\bfseries\color{dh-navy}}
\renewcommand{\cftchappagefont}{\bfseries\color{dh-navy}}
\renewcommand{\cftsecfont}{\color{dh-navy}}
\renewcommand{\cftsubsecfont}{\color{dh-navy!85}}
\renewcommand{\cftdotsep}{1.5}
\renewcommand{\cftchapleader}{\color{dh-blossom}\cftdotfill{\cftdotsep}}
\setlength{\cftbeforechapskip}{6pt}
\setlength{\cftbeforesecskip}{3pt}

% ── Hyperref (must be near end of preamble) ──
\usepackage[unicode,bookmarks=true,bookmarksnumbered=true]{hyperref}
\hypersetup{
  colorlinks=true,
  linkcolor=dh-navy,
  urlcolor=dh-gold,
  citecolor=dh-bronze,
  pdfborder={0 0 0},
  pdftitle={四倍做多认知，长期做多人生},
  pdfauthor={花花 \textbar{} @BerryUIKI}
}

% ── Image support ──
\usepackage{graphicx}
\usepackage{float}
% Multiple graphicspath entries so image resolution works regardless of
% whether xelatex is invoked from the workspace root, the output/ dir,
% or the output/tex/ split dir.  Each entry is searched in order.
\graphicspath{%
  {./}%
  {../}%
  {./articles/2026-quadruple-long-life/analysis/output/}%
  {../articles/2026-quadruple-long-life/analysis/output/}%
  {./assets/}%
  {../assets/}%
  {./images/}%
  {../images/}%
}
\usepackage[export]{adjustbox}

% Pandoc image wrapper: ensure images fit within page width and height
\makeatletter
\newsavebox\pandoc@box
\providecommand*\pandocbounded[1]{%
  \sbox\pandoc@box{#1}%
  \Gscale@div\@tempa{\textheight}{\dimexpr\ht\pandoc@box+\dp\pandoc@box\relax}%
  \Gscale@div\@tempb{\linewidth}{\wd\pandoc@box}%
  \ifdim\@tempb\p@<\@tempa\p@\let\@tempa\@tempb\fi%
  \ifdim\@tempa\p@<1\p@\scalebox{\@tempa}{\usebox\pandoc@box}%
  \else\usebox{\pandoc@box}%
  \fi%
}
\makeatother

% ── Code highlighting ──
\usepackage{listings}
\lstset{
  basicstyle=\ttfamily\footnotesize\color{dh-navy},
  backgroundcolor=\color{dh-white},
  frame=single,
  framerule=0.5pt,
  rulecolor=\color{dh-blossom},
  breaklines=true,
  numbers=left,
  numberstyle=\tiny\color{dh-blossom},
  keywordstyle=\color{dh-navy}\bfseries,
  commentstyle=\color{dh-bronze}\itshape,
  stringstyle=\color{dh-red},
  tabsize=2,
  showstringspaces=false
}

% ── Pandoc compatibility ──
\providecommand{\tightlist}{\setlength{\itemsep}{2pt}\setlength{\parskip}{1pt}}
\usepackage{longtable}
\usepackage{calc}
\setlength{\LTleft}{0pt}
\setlength{\LTright}{0pt}

% ── Footnote styling ──
\renewcommand{\footnoterule}{%
  {\color{dh-blossom}\hrule width 0.3\textwidth height 0.4pt}%
  \vspace{4pt}}

% ── Tables ──
\usepackage{booktabs}
\usepackage{array}
\usepackage{colortbl}
\usepackage{tabularx}
\arrayrulecolor{dh-blossom}
\setlength{\tabcolsep}{6pt}
\renewcommand{\arraystretch}{1.25}

% ── Block quotes (brand callout box) ──
\usepackage[most]{tcolorbox}
\newtcolorbox{dhquote}{
  blanker,
  breakable,
  borderline west={3pt}{0pt}{dh-blossom},
  left=12pt, right=8pt,
  top=4pt, bottom=4pt,
  before skip=0.6em, after skip=0.6em,
  fontupper=\small\color{dh-bronze}\itshape
}
\renewenvironment{quote}{\begin{dhquote}}{\end{dhquote}}

% ── Part divider macro (standalone full-page banner for "第N部分") ──
\newcommand{\partdivider}[1]{%
  \clearpage
  \thispagestyle{empty}%
  \vspace*{5.0cm}%
  \begin{center}%
    {\color{dh-blossom}\rule{0.35\textwidth}{1.5pt}\par}%
    \vspace{1.2cm}%
    {\Huge\bfseries\sffamily\color{dh-navy} #1\par}%
    \vspace{1.2cm}%
    {\color{dh-blossom}\rule{0.35\textwidth}{1.5pt}\par}%
  \end{center}%
  \vfill
  \clearpage
}

% ── Pandoc variables (filled in by metadata flags) ──
\newcommand{\mytitle}{四倍做多认知，长期做多人生}
\title{\mytitle}
\author{花花 \textbar{} @BerryUIKI}
\date{2026-09-07}

% ────────────────────────────────────────────────────────────
%  Body
% ────────────────────────────────────────────────────────────
\begin{document}

% ---- Custom Title Page (literal, bypasses ctex's maketitle) ----
\begin{titlepage}
\thispagestyle{empty}
\centering
\vspace*{3.5cm}

% Top decorative band
{\color{dh-blossom}\rule{0.55\textwidth}{2pt}\par}
\vspace{2.2cm}

% Main title (symmetrical 6+6 two-line title for elegant cover typography)
{\fontsize{34pt}{48pt}\selectfont\bfseries\color{dh-navy}%
四倍做多认知\\[0.35em]%
长期做多人生\par}

\vspace{0.6cm}

% Decorative divider
{\color{dh-gold}\rule{0.25\textwidth}{1.2pt}\par}
\vspace{0.4cm}

% Subtitle (English)
{\Large\color{dh-navy!85}Long-Term Investing Handbook\par}
\vspace{1.5cm}

% Author
{\large\color{dh-bronze}花花 \textbar{} @BerryUIKI\par}
\vspace{1.2cm}

% Mid decorative band
{\color{dh-blossom}\rule{0.35\textwidth}{1pt}\par}
\vfill

% Bottom — edition / date
{\large\color{dh-navy}第 1 版 · 2026-09-07\par}
\vspace{0.4cm}
{\color{dh-gold}\rule{0.55\textwidth}{1pt}\par}
\vspace{2cm}
\end{titlepage}

% ---- Frontmatter & Table of Contents ----
\frontmatter

\tableofcontents

\newpage

\chapter{扉页 · Opening Statement}\label{ux6249ux9875-opening-statement}

\begin{dhquote}
财富，不是为了炫耀，而是为了拥有选择人生的自由。
投资，不是为了赢过别人。而是为了让未来的自己，拥有更多选择。 Wealth is
not for showing off --- it is for the freedom to choose your own life.
Investing is not about beating others; it is about giving your future
self more choices.
\end{dhquote}

\newpage

\chapter{阅读须知 · Reading
Notice}\label{ux9605ux8bfbux987bux77e5-reading-notice}

\begin{dhquote}
欢迎翻开这本书。在开始阅读之前，希望你先了解几点。
\end{dhquote}

\textbf{第一，} 本书旨在分享长期投资理念、资产配置方法（asset
allocation）及个人研究观点，仅用于知识交流与学习参考，不构成任何形式的投资建议（investment
advice），也不应被视为买卖任何金融产品的依据。

\textbf{第二，}
投资有风险，市场存在波动，任何投资都可能产生亏损。书中提及的案例、市场、产品或历史数据，仅用于帮助读者理解相关概念，不代表未来表现，也不构成收益承诺。

\textbf{第三，}
本书主要面向中国大陆境外投资者编写。不同国家和地区对金融产品、证券投资及税务等相关规定存在差异，请读者在阅读和实践过程中，自觉遵守所在地法律法规，并根据自身情况独立作出投资决策。

\textbf{最后，}
希望你带着思考阅读这本书。我更希望你理解长期投资（long-term
investing）背后的逻辑，而不是记住某一只股票、某一只基金，或某一种策略。真正能够陪伴你很多年的，不是某一次正确的判断，而是一套持续受益的财富认知。愿这本书，能成为你长期投资旅程的一个起点。

\begin{dhquote}
\emph{Reading Notice (English):} Welcome to this book. Before you begin,
I hope you will keep a few things in mind. First, this book shares
long-term investing philosophy, asset allocation methods, and personal
research views for knowledge exchange and learning only. It does not
constitute investment advice of any kind, nor should it be taken as a
basis for buying or selling any financial product. Second, investing
carries risk; markets fluctuate, and any investment may incur losses.
Cases, markets, products, or historical data mentioned here serve only
to illustrate concepts --- they do not represent future performance and
are not a promise of returns. Third, this book is written primarily for
investors outside mainland China. Rules on financial products,
securities, and taxation differ by jurisdiction; please comply with the
laws where you live and make independent decisions based on your own
situation. Finally, I hope you read with a critical mind. Understand the
logic behind long-term investing rather than memorizing any single
stock, fund, or strategy. What stays with you for years is not one
correct call, but an enduring wealth mindset. May this book be a
starting point on your long-term investing journey.
\end{dhquote}

\newpage

\chapter{作者的话 · A Note from the
Author}\label{ux4f5cux8005ux7684ux8bdd-a-note-from-the-author}

\begin{dhquote}
很多人第一次接触投资，都觉得它离自己很远。有人认为，要有很多钱才能开始；也有人认为，只有金融专业人士才能做好投资。于是，投资被包装成一件复杂而神秘的事情，充斥着各种预测、技巧和所谓的''财富密码''。
\end{dhquote}

但在我看来，真正的长期投资并没有那么复杂，它更像是一种生活方式，一种面对未来的态度。

\textbf{对于我们四爱的老师们来说，最大的优势，从来不是本金，而是时间。}
也许我们现在没有丰厚的积蓄，没有丰富的投资经验，甚至还会因为一次市场波动而感到焦虑。但我们拥有未来二十年、三十年，甚至更长的时间。时间不会保证每一次投资都成功，却会让复利（compound
interest）慢慢积累，让经验不断沉淀，也让今天看似微不足道的坚持，在未来产生巨大的价值。年轻，本身就是一种别人无法复制的资本。

写这本书，并不是想告诉你哪只股票会涨，哪只基金值得买，更不是教你如何快速实现财富自由。我更希望与你分享一种长期投资的思维方式：理解财富、理解风险、理解资产配置（asset
allocation），也理解为什么投资的终点，从来都不是财富本身，而是让自己拥有更多面对人生的底气与选择。

这本书写给四爱的老师们，也写给每一个愿意认真生活、认真规划未来的人。希望它能够陪伴你，从第一笔投资开始，慢慢建立属于自己的财富认知。

愿多年以后，当你回头看时，会发现真正改变人生的，并不是某一次幸运的投资，而是那些日复一日、看似普通却始终没有放弃的坚持。因为真正值得长期做多的，不只是财富，更是你的人生。

\begin{dhquote}
------ 花花｜@BerryUIKI
\end{dhquote}

\newpage

\chapter{序章 ·
为什么普通人必须理解资产}\label{ux5e8fux7ae0-ux4e3aux4ec0ux4e48ux666eux901aux4ebaux5fc5ux987bux7406ux89e3ux8d44ux4ea7}

\begin{dhquote}
\textbf{收入决定今天的生活，资产决定未来的生活。}
\end{dhquote}

\begin{center}{\color{dh-bronze}\rule{0.2\textwidth}{0.5pt}}\end{center}

我们的主人公花有财------朋友们都叫他小名「花花」------今年二十五岁，从今天开始，每个月固定拿出五百元进行投资。假设长期年化回报率为
7\%（这只是教学假设，不代表未来收益），一直坚持到六十岁退休。到那时，他总共投入了二十一万元，账户里大约会有九十万元。

他的朋友小李情况完全相同，唯一不同的是------小李等到三十五岁才开始。同样是每个月五百元、同样的收益率、同样到六十岁。小李总共投入了十五万元，账户里大约会有四十一万元。

早开始十年，只是多投入了六万元。但最终的结果，相差了大约五十万元。

这五十万元的差距，并不是来自谁更聪明、谁更努力、谁的收入更高。它来自一个很多人从学校毕业以后就不再被提醒的事实：\textbf{时间，本身就是一种可以积累的资源。}
年轻的时候，时间是最充裕的资本；年纪越大，它的价值就越昂贵。而大多数人直到真正需要用钱的时候，才开始后悔当初没有更早地开始。

这正是这本书想要讨论的问题。

我们生活在一个对''赚钱''充满焦虑的时代。社交媒体上充满了各种令人羡慕的生活方式，理财课程和投资分享随处可见。很多人一边努力工作，一边隐隐觉得，仅靠工资好像永远追不上生活成本的增长。但另一方面，真正去了解投资的人并不多。原因也很简单------投资这个词，听起来太复杂了。它和数字有关，和风险有关，和那些看起来需要读很多报告才能理解的事情有关。很多人宁愿把这件事一直往后推，等存款再多一点、等工资再高一点、等市场再好一点、等自己再多学一点。但一年又一年过去，真正开始的那一天始终没有到来。

这本书不会告诉你哪只股票会涨，也不会承诺任何短期内快速致富的方法。事实上，如果有人向你承诺这些，那大概率不是投资，而是投机，甚至可能是骗局。投资和投机的区别很简单：投资是让资产替你工作，分享企业和社会长期创造的价值；投机是猜测别人愿意以什么价格接盘。前者需要时间和耐心，后者需要运气和速度。这本书讲的，全部是前者。

它的目标也很简单：帮助你建立一套关于财富和投资的完整认知框架。从理解自己的财务状况开始，到认识不同类型的资产，再到建立属于自己的长期投资计划。你不需要成为金融专家，也不需要每天盯盘。你需要的是理解一些最基本、但最重要的原则------然后让时间替你完成剩下的事情。

当然，这本书不能替你赚钱。它能做的，是让你在未来的很多年里，面对市场的涨跌、面对别人的建议、面对自己的情绪时，能够独立做出判断，而不是被恐惧和贪婪推着走。

也许读完这本书，你依然不会记住某一只基金的名字，也不会预测下一次市场的方向。但我希望你能带走的，是几样更重要的东西：知道自己的钱到底放在哪里、为什么放在那里；知道什么时候应该坚持、什么时候应该调整；知道投资不是为了战胜别人，而是为了让未来的自己拥有更多选择。

如果这本书能够让你在今天迈出第一步------哪怕只是开始记录自己的收支、开始理解自己手中的资产、开始为未来做一次认真的规划------那么它就已经完成了它的使命。

因为真正改变人生的财富，从来不是一笔意外之财，而是那些在无数个普通日子里，一点一点积累起来的认知和习惯。

翻到下一页，我们从最基础的问题开始。

\begin{center}{\color{dh-bronze}\rule{0.2\textwidth}{0.5pt}}\end{center}

\emph{本书所有投资案例均为教学假设，不构成任何投资建议。历史收益不代表未来表现。投资有风险，入市需谨慎。}

\mainmatter

\partdivider{第一部分 · 先看清自己的财务生活}

\newpage

\chapter{财富从哪里来？}\label{ux8d22ux5bccux4eceux54eaux91ccux6765}

\begin{center}{\color{dh-bronze}\rule{0.2\textwidth}{0.5pt}}\end{center}

\begin{dhquote}
\textbf{你不需要成为投资专家才能积累财富。你只需要理解一个简单的逻辑：收入让你活下去，但只有资产能让你活出选择。}
\end{dhquote}

\begin{center}{\color{dh-bronze}\rule{0.2\textwidth}{0.5pt}}\end{center}

2014年，美国佛蒙特州一位名叫罗纳德·里德（Ronald
Read）的老人去世，享年92岁。他在当地一家加油站做了大半辈子的维修工，也当过清洁工，从未做过什么引人注目的事。然而当他去世后，人们发现这位修理工留下了将近800万美元的遗产------其中一部分捐给了当地图书馆和医院。他的邻居、同事，甚至家人都不知道他有这么多钱。

他是怎么做到的？没有中彩票，没有继承遗产，没有创业暴富。他只是用几十年的时间，从工资里省下一小部分，持续地买入他了解和相信的优质公司股票，然后什么都不做，让时间来完成剩下的工作。

罗纳德·里德的一生证明了一件事：财富和收入，是两个完全不同的概念。你可以赚得不多但积累很多，也可以赚得很多但一无所有。分界线不在于你拿到多少工资，而在于你拿工资之后做了什么。

\subsection{三十五年之后，差距会有多大？}\label{ux4e09ux5341ux4e94ux5e74ux4e4bux540eux5deeux8dddux4f1aux6709ux591aux5927}

让我们做一个有具体数字的思想实验。

假设两个人从25岁开始工作，月薪都是8000元------在一个普通二线城市，这是非常典型的起点。第一个行动者，每个月的发薪日，把500块钱放进一只跟踪沪深300的指数基金里。他选这只基金没有特别的理由------只是因为它的费率低、覆盖广、不需要花时间研究个股。500块钱大约是他月收入的6\%。剩下的7500元用来租房、吃饭、交通、偶尔和朋友聚餐。

他坚持了35年。在这35年里，这只指数基金的平均年化收益率大约是7\%------并不是一个惊人的数字，但它是一个长期合理的假设。到60岁的时候，他一共投入了21万元的本金（500元
× 12个月 × 35年）。但由于复利的作用，这笔钱的最终规模大约为90万元。

第二个同龄人，月薪同样是8000元。他觉得每个月省500块毫无意义------``就这么点钱，能干什么呢？''他每次发工资都花得干干净净。他过得并不奢侈------只是从没想过''留下一点''这件事------到60岁时，他的银行账户余额大致就是最后几笔工资加社保。

这两个人的智商、能力、运气和起点完全相同。三十五年后，第一个人拥有近90万元的投资资产，而且这笔钱在退休后还在为他产生收入；第二个人几乎没有可自由支配的金融资产。这个差距不是由''谁更能赚钱''决定的------他们的月薪完全一样------而是由''谁更早理解了一个简单的财富逻辑''决定的。

这个逻辑就是：\textbf{收入是流量，财富是存量。}
工资是每个月流进池子里的水；财富是池子里积累的水位。如果你每个月把流进来的水全部用光------无论进水口有多大------池子里的水位永远不会上升。

\subsection{三种收入，只有一种不用你亲自到岗}\label{ux4e09ux79cdux6536ux5165ux53eaux6709ux4e00ux79cdux4e0dux7528ux4f60ux4eb2ux81eaux5230ux5c97}

经济学家通常把收入分成三类。理解这三者的本质区别，是建立财富认知的第一个关键节点。

第一种是\textbf{工资收入（earned
income）}：用时间和技能交换报酬。月薪五千和月薪五万，在这个分类里没有本质区别------都是''你停止工作，收入就停止''。工资收入最致命的弱点不是数额大小，而是它与你的肉身牢牢绑定。你生病了，它停了。你想休息半年，它停了。你到了退休年龄，它也必须停。

第二种是\textbf{经营收入（business
income）}：来自拥有或运营一个实体------自由职业、小生意、创业公司。它打破了工资的天花板------一个成功的店主可以比一个高级工程师赚得多得多。但它仍然绑定你的时间。开餐馆的老板和写代码的程序员，在''必须亲自到场才有收入''这一点上，本质上属于同一个类别。

第三种，也是这本书要重点讨论的，是\textbf{资本收入（capital
income）}：来自你持有的资产所产生的回报------股票的股息和增值、债券的利息、房产的租金、基金的收益。这笔收入最大的不同是：它不需要你亲自出现。无论你今天在上班、在旅行、在睡觉、还是在医院，只要你持有能创造价值的资产，它们就在持续为你工作。

三种收入之间最核心的转换逻辑不是''一夜之间放弃工资去投资''------那既不现实也不安全。真正的路径是：在你仍然拥有工资收入的阶段，有意识地把一部分工资收入，持续地、有纪律地转化为能够产生资本收入的资产。这个\textbf{转化率}------你每个月把工资的百分之多少变成了资产------比你工资的绝对数字，更准确地预示着你二十年后的财务自由度。

一个每月赚两万但转化率为零的人，二十年后拥有的资产可能远不如一个每月赚八千但转化率为15\%的人。这就是为什么本章开头那位加油站的维修工，可以比很多高薪白领更富有。

\subsection{通货膨胀：最安静的财富杀手}\label{ux901aux8d27ux81a8ux80c0ux6700ux5b89ux9759ux7684ux8d22ux5bccux6740ux624b}

在讨论如何积累资产之前，还有一个更根本的问题：为什么钱不能老老实实地放在银行里？

答案是\textbf{通货膨胀（inflation）}：货币购买力持续下降的宏观经济现象。它不是一次性的冲击，不是偶尔发生的危机，而是以一种安静到几乎不被察觉的方式，年复一年地侵蚀你的储蓄。

以中国的消费者价格指数（CPI）来看，过去三十多年间，年均通胀率大约在2\%到5\%的区间内波动，某些年份更高。换算成更直观的画面：1990年初的100块钱，如果只是压在床垫底下，到2024年，它还能买到的东西大约只剩下原来的三分之一。三十四年前能买一双不错的皮鞋的钱，今天只够买一双拖鞋。

这并不是中国特色。大洋彼岸的美国，通胀看起来温和得多------过去三十年的年均CPI大约在2.5\%的水平。但即便是这个''温和''的数字，累积三十四年之后，也让1990年的100美元，到2024年需要大约241美元才能买到相同的一篮子商品和服务。

这个数学的残酷在于它的不对称性：2\%到3\%的年通胀率------看起来不值一提对吧？但它遵循的是指数增长曲线，而不是线性增长。每年2.5\%的通胀，累积三十年，购买力减半；累积五十年，只剩下不到三分之一。

这就是为什么对于任何一个投资期限还有二三十年的人来说，``只存不投''的选择，不是在''保本''，而是在用一种确定的、缓慢的方式接受一种确切的损失。通货膨胀最有效的对冲策略，就是把钱配置到那些长期回报率能够超越通胀的资产上------股票、债券、不动产、黄金------而不仅仅是依赖银行存款的微薄利息。存款的作用是安全垫和流动性------它当然要有------但它不应该是你所有钱的唯一去处。

\subsection{资产和消费品：一道改变命运的选择题}\label{ux8d44ux4ea7ux548cux6d88ux8d39ux54c1ux4e00ux9053ux6539ux53d8ux547dux8fd0ux7684ux9009ux62e9ux9898}

日常对话中，我们经常把''买了一个很贵的东西''说成''拥有了一项资产''。但从财务本质上讲，这两个词的含义截然相反。

\textbf{资产（asset）}，是那些能够持续为你创造价值、带来未来经济收益的东西。一只指数基金是资产，因为它代表着一篮子上市公司的部分所有权------那些公司每天都在经营、竞争、盈利，而你作为股东分享这些成果。一张国债是资产，因为国家承诺按期支付利息，到期归还本金。一套出租用房是资产，因为它每个月能带来持续的租金。资产的核心特征，是它会随着时间\textbf{增值}或\textbf{产生现金流}。

\textbf{消费品（consumption
good）}，是你购买来使用或享受的东西------新车、新手机、一顿精致的晚餐、一只名牌包。消费品让你的当下生活变得更好，这对生活品质当然重要。但它们通常不会为你产生一分钱现金流，而且大部分从你买到的那一刻起就开始贬值。一辆新车开出4S店，残值可能已经打了九折；一部手机用了两年，转手价格往往不到原价的一半。消费品的核心特征，是它会随着时间\textbf{折旧}或被\textbf{消耗}。

两种选择都没有道德上的对错------享受生活当然应该。但这里有一个认知层面的分水岭：你在做出每一笔较大支出时，有没有意识到自己在做一道有长期后果的选择题？

举两个例子。第一个，更贴近日常：你打算花七千块钱换一部最新款的手机。你也可以不换------把这七千块钱放进一只年化收益率7\%的指数基金。三十年后，这笔钱大约变成了五万三千元。而那部手机------除非它是一台会被收藏的初代iPhone------大概率已经躺在抽屉里，屏幕可能裂了，电池可能鼓包了，连开机都做不到了。

第二个，金额再大一些：两万块钱买一只名牌包和两万块钱买一只指数基金。在刷卡的那一瞬间，感觉可能一样------都是一笔''花出去''的钱。三十年后，基金变成了大约十五万元；而那只包------最好的情况是它还在衣柜的某个角落，也许已经好几年没有背过了。

两个例子讲的其实是同一道算术题：你付出去的价格，到底买到了一个会随时间增值的资产，还是一个会随时间折旧的消费品？金额大小在这道题里没有区别------它考的不是你''有多少钱''，而是你''在做选择的那一刻，脑子里有没有这笔账''。

当然，我并不想建议你把所有的钱都拿来投资而完全不消费------既不现实也不快乐。我想建议的是：把''每一笔较大的支出，都是一道资产还是消费品的选择题''这个意识，变成你财务决策的默认起点。当你开始用这种眼光审视自己的每一笔支出时，你在未来十年、二十年里积累下来的差距，会超过你的想象。

\subsection{财富的真正尺度：选择的自由}\label{ux8d22ux5bccux7684ux771fux6b63ux5c3aux5ea6ux9009ux62e9ux7684ux81eaux7531}

如果问一个真正拥有充足资产的人''什么是财富''，他大概率不会回答''财富就是有很多钱''。他会说：\textbf{财富，是你可以对不想做的事情说不的能力。}

当你拥有足够的资本收入来覆盖基本生活开支时，你获得了这个世界上大多数人一辈子都没有的东西：选择权。你可以选择接受一份薪水更低但真正让你有热情的工作------因为工资已经不是你的主要收入来源。你可以在家人最需要你的时候放下工作去陪伴------因为你的资产还在不间断地为你的家庭产生现金流。你可以在某一天决定，职业生涯到此为止了，接下来的时间完全属于自己------不是''被迫退休''然后为每月的支出精打细算，而是主动选择进入人生的下一个章节。你也可以选择什么都不改变，继续做现在的工作------但那是你的选择，不是你的被迫。

这才是财富的真正功能：不是用来炫耀的数字，不是社会地位的标签，而是\textbf{购买自由的资本}。很多人花一辈子去追逐更高的工资、更好的车、更大的房子、更贵的消费------但他们从来没有停下来问过一个更核心的问题：我是在积累''选择的自由''，还是在积累''更多需要维护的负担''？

理解收入与财富之间的这一道分界线，是这本书的起点。积累财富不需要你出身富贵，不需要你运气过人，不需要你比同事更聪明。它需要的条件比你想象中要简单得多：在你的工资还在稳定流入的阶段，把这道分界线看清楚------然后开始把你每个月收入的一部分，一砖一瓦地，变成能为你创造更多收入的资产。

从下一章开始，我们将用最简单的方法，先把你的真实财务状况算清楚------你有多少、欠多少、你的起点在哪里。先把地图打开，再谈怎么走。

\newpage

\chapter{先计算自己的净资产}\label{ux5148ux8ba1ux7b97ux81eaux5df1ux7684ux51c0ux8d44ux4ea7}

\begin{center}{\color{dh-bronze}\rule{0.2\textwidth}{0.5pt}}\end{center}

\begin{dhquote}
\textbf{知道自己拥有什么，比拥有什么更重要。}
\end{dhquote}

\begin{center}{\color{dh-bronze}\rule{0.2\textwidth}{0.5pt}}\end{center}

如果把个人理财比作体检，那么计算净资产就是第一项基础检查。你不一定需要每天都盯着这个数字，但一年至少应该认真看一次------它会告诉你，过去一年的努力，究竟有没有让财务状况变得更加健康。

\subsection{资产：你拥有的，不全是资产}\label{ux8d44ux4ea7ux4f60ux62e5ux6709ux7684ux4e0dux5168ux662fux8d44ux4ea7}

从财务的角度，资产可以分为几类，每一类在人生中扮演的角色都不相同。

\textbf{现金及现金等价物（cash and cash
equivalents）}，包括银行活期存款、定期存款和货币基金。它们最大的优势是流动性极高，随时可以取用。代价是收益通常很低，长期持有过多现金反而会面临购买力下降的问题。这类资产的职责，是覆盖日常开支和应对突发情况，而不是替你创造财富。

\textbf{金融资产（financial assets）}，包括股票、债券、基金、ETF
以及各种理财产品。这类资产是长期财富增长的主力。它们的价格会波动，有时波动还不小，但长期来看，优秀企业创造的价值会逐渐反映到金融资产的价格上。

\textbf{实物资产（physical
assets）}，包括自住房、投资性房产、汽车、贵重物品等。实物资产看得见摸得着，给人一种踏实的安全感。但它们的流动性通常较差------卖一套房子可能需要几个月，而在这期间如果急需用钱，就可能不得不接受一个不理想的价格。

在做房产估值时，需要根据所在城市的不同层级有所区分。一线城市（北京、上海、广州、深圳）的房产具有更强的流动性和更稳定的价格基础，通常比较容易找到市场价格参考。二线城市（如杭州、成都、武汉、南京等省会及强地级市）的房产市场活跃度次之，估值时应适当参考近期同小区的实际成交价而非挂牌价。三、四线及以下城市的房产，流动性相对较低，变现周期更长，在计算净资产时应当更加保守------可以将估值打一个折扣（比如参考价的
80\%---90\%），以反映在市场低迷时可能出现的折价风险。如果房产位于人口持续流出的地区，长期来看资产增值空间有限，更不宜在净资产计算中给予过高权重。

\textbf{退休账户和其他长期储蓄（retirement
accounts）}，包括个人养老金账户、企业年金、公积金等。这些资金通常有特定的提取限制，但也可能享受税收优惠。在做财务规划时，它们应当单独列出，因为退休账户里的钱在退休前通常不应被动用，它也不是你日常投资组合的一部分。

\textbf{人力资本（human
capital）}，这是一个容易被忽略、但实际上可能是年轻人最重要的资产。简单来说，人力资本就是你未来所有劳动收入的现值。一个三十岁的软件工程师，就算现在银行账户里没有太多存款，但未来几十年的职业生涯将带来数百万的工资收入。从这个角度看，年轻人的净资产往往被严重低估。理解这一点，也能帮助你更理性地评估自己能承受多少投资风险------只要人力资本还在增长，你就有能力承受短期内的市场波动，也有能力从错误中学习和调整。

\begin{center}{\color{dh-bronze}\rule{0.2\textwidth}{0.5pt}}\end{center}

\subsection{负债：欠别人的，要分清楚优先级}\label{ux8d1fux503aux6b20ux522bux4ebaux7684ux8981ux5206ux6e05ux695aux4f18ux5148ux7ea7}

负债听起来让人不舒服，但并不是所有负债都是坏事。关键在于区分负债的类型和成本。

\textbf{高息负债（high-interest
debt）}，最典型的是信用卡透支和某些消费贷款。它们的年化利率往往在 15\%
到 20\%
甚至更高。长期持有这种负债，相当于你每年都在以远超投资收益率的成本''养''着债务。在开始投资之前，处理高息负债应该排在第一位------因为没有任何投资能够稳定地提供
18\% 的年化回报，但提前还清信用卡债务，实际上就等于''赚''到了 18\%。

\textbf{住房贷款（mortgage）}，是最常见的中长期负债。房贷的利率通常远低于消费贷，而且房产本身的价值有可能随着时间增长。在大多数情况下，不必急于提前还清所有房贷------如果房贷利率低于你的长期投资预期回报，将多余的现金流用于投资可能更划算。但这也不是绝对的，每个人的风险承受能力和心理舒适度不同，有些人不喜欢背债，提前还清也是一种合理的选择。关键是根据自己的情况做决定，而不是听别人说''房贷一定要提前还''或者''房贷一定不要还''。

\textbf{学生贷款（student
loan）}，在中国主要包括国家助学贷款和各类教育贷款。助学贷款的利率通常较低，还款期限也比较灵活，是国家为了支持教育而提供的福利性金融工具。如果你的助学贷款利率很低（比如基准利率甚至免息期间），不必急于提前偿还------把有限的现金流用于建立应急资金和开始投资，在财务上往往是更优的顺序。但要注意，助学贷款虽然利率低，它仍然是一笔真实的负债，需要在毕业后按约定偿还，不应因为''利率低''就不放在心上。对于已经毕业并开始工作的读者，建议把助学贷款的月供金额纳入固定支出预算，在计算每月可投资金额时一并考虑。

\begin{center}{\color{dh-bronze}\rule{0.2\textwidth}{0.5pt}}\end{center}

\subsection{净资产：最后那个数字}\label{ux51c0ux8d44ux4ea7ux6700ux540eux90a3ux4e2aux6570ux5b57}

把你列出的所有资产加在一起，减去所有负债的总额，得到的就是你的净资产。

如果你刚工作不久，净资产是负数，这并不奇怪。助学贷款、刚买房的按揭贷款，都会在早期让净资产看起来不太好看。真正需要关注的不是今天的数字，而是它每年的变化方向。

如果你工作了很多年，收入也不错，但净资产增长缓慢，那就值得审视一下自己的消费习惯了。每个月花掉的钱，有多少是真正让自己生活变好的，又有多少是在不知不觉中消失的？

净资产的变化，是个人财务健康最诚实的指标。它不给任何人留情面------赚得多但花得更多的，净资产不会说谎；收入不高但坚持储蓄投资的，年复一年，净资产会替你证明这一切的努力都是值得的。

\begin{center}{\color{dh-bronze}\rule{0.2\textwidth}{0.5pt}}\end{center}

下面是一个简化的个人资产负债表模板。你可以拿出纸笔，或者打开手机备忘录，花十分钟填一下。不用追求精确到分，大致准确就够了。

\begin{longtable}[]{@{}llll@{}}
\toprule\noalign{}\rowcolor{dh-pink!25}
资产 & 金额（元） & 负债 & 金额（元） \\
\midrule\noalign{}
\endhead
\bottomrule\noalign{}
\endlastfoot
现金和活期存款 & \_\_\_\_\_\_\_\_\_ & 信用卡透支 & \_\_\_\_\_\_\_\_\_ \\
定期存款和货币基金 & \_\_\_\_\_\_\_\_\_ & 消费贷款 &
\_\_\_\_\_\_\_\_\_ \\
股票 / 基金 / ETF & \_\_\_\_\_\_\_\_\_ & 汽车贷款 &
\_\_\_\_\_\_\_\_\_ \\
投资性房产（市值） & \_\_\_\_\_\_\_\_\_ & 住房贷款余额 &
\_\_\_\_\_\_\_\_\_ \\
自住房产（市值） & \_\_\_\_\_\_\_\_\_ & 其他借款 & \_\_\_\_\_\_\_\_\_ \\
退休账户 / 公积金 & \_\_\_\_\_\_\_\_\_ & & \\
其他资产 & \_\_\_\_\_\_\_\_\_ & & \\
\textbf{资产总计} & \_\_\_\_\_\_\_\_\_ & \textbf{负债总计} &
\_\_\_\_\_\_\_\_\_ \\
& & \textbf{净资产（资产 - 负债）} & \_\_\_\_\_\_\_\_\_ \\
\end{longtable}

\begin{center}{\color{dh-bronze}\rule{0.2\textwidth}{0.5pt}}\end{center}

自住房要不要列入这个表格，其实取决于你的目的。如果你只是想了解自己''有多少钱''，把它放进去没问题。但如果你是在为投资做规划，建议单独列出可投资资产------也就是扣除自住房和退休账户之后、可以用来长期投资的资金。毕竟你不会因为股价跌了就把卧室卖掉，也不应该在临近退休的时候把养老金全部投进股票市场。

\begin{center}{\color{dh-bronze}\rule{0.2\textwidth}{0.5pt}}\end{center}

计算净资产不是为了和任何人比较。它只是给自己的一个参照点。你今天填下这些数字，一年以后再填一次，如果净资产变多了，而且增长主要来自资产的增值而不是新借的债务，那么你就在正确的方向上。

这个参照点，是所有长期投资规划的起点。因为只有知道自己站在哪里，才知道下一步应该往哪个方向走。

\newpage

\chapter{储蓄率比收入数字更重要}\label{ux50a8ux84c4ux7387ux6bd4ux6536ux5165ux6570ux5b57ux66f4ux91cdux8981}

\begin{center}{\color{dh-bronze}\rule{0.2\textwidth}{0.5pt}}\end{center}

\begin{dhquote}
\textbf{赚多少钱决定你的起点，留下多少决定你能走多远。}
\end{dhquote}

\begin{center}{\color{dh-bronze}\rule{0.2\textwidth}{0.5pt}}\end{center}

社交场合里，人们习惯用收入来衡量一个人的财务状况。月薪两万、年薪五十万，这些数字听起来直观又有冲击力。但如果你认真观察就会发现，收入相近的人，十年后的财富状况可能天差地别。有人收入不低却一直存不下钱，有人工资平平却慢慢积累起可观的资产。拉开差距的往往不是赚钱的能力，而是留钱的能力。

这个能力有一个简单的衡量指标，叫做\textbf{储蓄率（savings rate）}。

\subsection{储蓄率，是你对自己未来的投资比例}\label{ux50a8ux84c4ux7387ux662fux4f60ux5bf9ux81eaux5df1ux672aux6765ux7684ux6295ux8d44ux6bd4ux4f8b}

储蓄率的计算方法很简单：每个月的储蓄金额，除以税后到手收入。如果你每月到手一万，花掉七千，存下三千，储蓄率就是
30\%。如果全部花光，储蓄率就是
0\%。如果花的比赚的还多，储蓄率就是负数------这种情况在信用卡普及的时代并不少见。

有人会觉得，月薪三千的时候谈储蓄率没有意义，存下来的钱太少，不如等收入高了再说。这种思维的问题在于，它假设储蓄是一种''收入高了自然就会做的事情''。但现实是，如果月薪三千的时候没有建立储蓄习惯，月薪三万的时候大概率也不会突然变得节俭。人的消费是会随着收入一起增长的------这在行为经济学里叫做\textbf{生活方式通胀（lifestyle
inflation）}。涨薪之后，外卖从三十元升级到一百元、租房从合租升级到整租、出行从公交升级到打车------每一样看起来都合情合理，但加起来，储蓄率可能和涨薪前没有任何区别，甚至更低了。

储蓄率真正重要的地方在于，它是你未来投资的本金来源。不论你的投资能力有多强，如果没有本金，一切都无从谈起。而本金从哪里来？不是从天上掉下来的，是从你每个月主动留下的那笔钱里，一点一点积累出来的。

\subsection{搞清楚你的钱去了哪里}\label{ux641eux6e05ux695aux4f60ux7684ux94b1ux53bbux4e86ux54eaux91cc}

很多人不是不愿意存钱，而是根本不知道自己每个月花掉了多少。

固定支出，比如房租、房贷、物业费、保险、通讯费，这些是有规律可循的，通常按月或按年支付，金额也比较稳定。可变支出则不同------吃饭、购物、娱乐、交通、社交、临时的各种开销，零零散散，但加起来可能比固定支出还多。很多人到月底一看账户余额，只会觉得''好像也没买什么大件''，但钱确实不见了。

记录一个月完整的开支，是了解自己消费习惯最直接的方式。不需要什么专业软件，手机备忘录就够了。每一笔都记下来，哪怕是十几元的外卖、几元钱的会员续费------一个月以后你可能会惊讶地发现，有些支出类别比你想象中大得多。而这些''隐形支出''往往是最容易优化的地方。每天一杯咖啡，一年下来可能是一趟短途旅行的费用；每天一顿外卖比自己做贵出来的差额，一年下来可能够买一份不错的保险。

当然，这不是在倡导一种清苦的生活方式。削减开支的目的不是为了过得拮据，而是为了把省下来的钱放在更重要的事情上。如果你真的喜欢每天喝一杯咖啡，那就继续喝------但只要你知道每个月这笔开销是多少，并且确信它对你很重要，那它就是一笔合理的开支。理财从来不是不让花钱，而是让你知道每一笔钱花得值不值。

\subsection{多少储蓄率算合理}\label{ux591aux5c11ux50a8ux84c4ux7387ux7b97ux5408ux7406}

这个问题没有标准答案，但有一些可以参考的框架。

对于刚开始工作的年轻人来说，如果能够把储蓄率保持在 20\% 到
30\%，已经是一个非常不错的起点。如果刚毕业、收入还不高，能做到 10\%
也很值得鼓励------关键在于养成习惯，而不是追求数字。随着收入增长，理想情况下储蓄率应该逐步提高，而不是和几年前保持一致。这是因为基本生活支出有一个相对固定的下限，收入提升之后，新增的部分不应该全部被消费吃掉。

有一种更激进的理念叫做 \textbf{FIRE（Financial Independence, Retire
Early，财务独立、提前退休）}，倡导者会把储蓄率推到 50\%
甚至更高，目标是在十年到二十年内积累足够的资产，实现不依赖工资的生活。这种生活方式并不适合大多数人，但它至少揭示了一个基本道理：储蓄率越高，达到财务独立所需的时间越短。即便你并不打算提前退休，只是为了让自己在四五十岁的时候拥有更多选择权，提高几个百分点的储蓄率，放在二十年、三十年的时间里，积累下来的差异也会非常惊人。

\subsection{收入不稳定的人怎么办}\label{ux6536ux5165ux4e0dux7a33ux5b9aux7684ux4ebaux600eux4e48ux529e}

前面讨论的都是按月领工资的场景。但并不是所有人都有稳定的现金流。自由职业者、个体经营者、佣金制销售、季节性行业从业者------他们的收入可能某个月很高，某个月很低，甚至没有。

对于收入波动较大的人，储蓄的逻辑需要调整一下。传统的那种''每月固定存一笔''的方法可能会在某些月份根本执行不了，而在收入较高的月份又容易因为''这个月赚得多''而放松消费。

一个适合不规律收入的储蓄方法是''收入分账''。每当一笔收入到账，先按照一个固定的比例把储蓄部分划走------比如
30\%------剩下的才是当月或当时可以支配的金额。收入高的月份，储蓄的绝对金额自然会更多；收入低的月份，储蓄也会相应减少。但关键是，这个动作要在花钱之前完成，而不是等月底看剩下多少再决定存多少。因为如果等到月底，``剩下的''往往已经不剩了。

对于那些收入本身就覆盖不了基本支出的特殊时期，储蓄当然可以暂停。这并不意味着失败，而是一种务实的短期调整。重要的是，当收入恢复之后，重新回到存钱的轨道上来，而不是因为中断过一次就彻底放弃。

\begin{center}{\color{dh-bronze}\rule{0.2\textwidth}{0.5pt}}\end{center}

储蓄率不是一个用来比较的数字。它是一个你留给未来的自己的比例。今天多留下一个百分点，未来就多收到一份来自过去的礼物。而这份礼物会在投资的作用下越变越大------关于这一点，我们在后面的章节里详细展开。

现在，你已经清点过自己的净资产，也看清楚了自己每个月能留下多少钱。接下来的问题是：在正式开始投资之前，还有什么需要先处理好的事情？

\newpage

\chapter{投资之前，先处理负债和应急资金}\label{ux6295ux8d44ux4e4bux524dux5148ux5904ux7406ux8d1fux503aux548cux5e94ux6025ux8d44ux91d1}

\begin{center}{\color{dh-bronze}\rule{0.2\textwidth}{0.5pt}}\end{center}

\begin{dhquote}
\textbf{最好的投资决策，有时候是不投资------先把地基打牢。}
\end{dhquote}

\begin{center}{\color{dh-bronze}\rule{0.2\textwidth}{0.5pt}}\end{center}

在投资世界里，大多数人最关心的问题是''买什么、什么时候买、买多少''。但对于一个刚开始规划财务生活的人来说，这些问题都应该排在一些更基本的事情之后。就像盖房子，你不能跳过地基直接搭屋顶------地基不稳，上面盖得再漂亮也没有意义。

投资的''地基''有两件事：清理高成本负债，建立应急资金。

\subsection{先还掉最贵的债}\label{ux5148ux8fd8ux6389ux6700ux8d35ux7684ux503a}

负债本身并不可怕。一个刚买房的人可能背负着几百万元的房贷，但他的净资产未必是负数------房产本身也是一项资产。一个刚毕业的学生可能还有几万元的助学贷款，但未来几十年的工资收入足以覆盖这些债务。

真正需要警惕的，是\textbf{高息负债（high-interest debt）}。

在中国市场，信用卡透支或分期付款的实际年化利率常常在 15\% 到 18\%
之间。一些消费金融产品的利率可能更高。你不需要成为一个投资专家也能看出来：几乎没有任何合法投资能够稳定地提供
15\% 以上的年化回报。这就意味着，如果你一边持有年利率 18\%
的信用卡债务，一边每个月拿出几百元买基金，你实际上是在做一笔稳赔不赚的买卖------基金可能赚
7\%，但债务在确定地吃掉 18\%，净损失 11\%。

因此，投资之前的第一个动作，应该是把高息负债尽可能清零。这比任何投资都更''划算''。因为提前还清一笔年利率
18\% 的贷款，等价于你获得了一笔年化 18\%
的无风险收益------市场上根本不存在这样的投资产品。

不同类型的负债，处理优先级也不一样。

\textbf{第一优先级：信用卡透支和高息消费贷。}
这些应该尽快还清，不要在付高利息的同时还想着投资。

\textbf{第二优先级：利率适中的贷款。} 比如一些银行的消费贷款，年利率在
5\% 到 10\%
之间。是否提前偿还，可以做一个简单的比较：如果贷款利率明显高于你预计的长期投资回报率，提前还贷是理性的选择；如果两者差不多或贷款利率更低，可以考虑按正常节奏还款，把多余的现金流用于投资。但这里的''预计回报率''必须保守估计------用
5\% 到 7\% 作为长期股票投资的参考区间，而不是用某一年牛市里的 30\%
来做比较。

\textbf{第三优先级：住房贷款。} 在中国，首套房的公积金贷款利率在 3\%
左右，商业贷款利率在过去几年也经历了明显下降。如果你的房贷利率在 3\% 到
4\%
之间，从纯数学的角度，把钱用于长期投资而非提前还贷，在历史上往往是更优的选择。但财务决策不完全是数学题------有些人天生不喜欢背债，提前还清房贷带来的心理舒适感，也是一种真实的回报。两种选择都没有错，取决于你对风险的感受。

\subsection{应急资金：你人生的安全气囊}\label{ux5e94ux6025ux8d44ux91d1ux4f60ux4ebaux751fux7684ux5b89ux5168ux6c14ux56ca}

处理完高息负债之后，下一件事不是急着去买基金，而是先准备一笔应急资金（emergency
fund）。

应急资金的作用很简单：让你的生活不会因为一次意外而被迫改变长期投资计划。你突然失业了，家人突然生病了，房子突然需要大修------这些事情每一年都在无数家庭中发生，但它们大多是不可预测的。如果你没有一笔随时可以动用的现金储备，你就可能不得不在最糟糕的时机卖出投资资产。比如刚好赶上市场大跌，你本来只需要套现一万，结果因为资产价格缩水，实际损失远超一万。

应急资金应该准备多少？一般的建议是 \textbf{3 到 6
个月的基本生活支出}。所谓''基本生活支出''，指的是在没有收入的情况下，维持最低生活标准所需的开销------房租、房贷、食物、水电、交通、基本医疗。不包括旅行、娱乐、餐厅消费。这不是要求你过苦日子，而是让你在最坏情况下有足够的时间找到新的收入来源，而不必因为缺钱而贱卖资产。

三个月还是六个月，取决于你的收入稳定性。如果你在一家稳定的公司、有着明确的职业晋升路径，三个月的备用金可能就够了。如果你是自由职业者、收入波动较大、或者所在的行业变化很快，那么准备六个月甚至更久会更安心。同理，如果你有家庭需要供养，应急资金也应该比单身时留得更多。

具体来说，可以参考以下比例：

\begin{itemize}
\tightlist
\item
  \textbf{收入非常稳定}（公务员、事业单位、大型国企等）：3
  个月的基本支出作为应急资金通常足够。这类职业的裁员风险极低，即使发生调动或岗位调整，收入也不会完全中断。
\item
  \textbf{收入较为稳定}（正规企业、有固定劳动合同）：建议准备 4 到 5
  个月的应急资金。虽然有劳动合同保障，但行业波动或企业调整仍可能影响收入。
\item
  \textbf{收入不够稳定}（自由职业者、佣金制销售、初创企业员工、季节性行业）：建议至少准备
  6 个月，如果条件允许可以逐步积累到 9 到 12
  个月。收入越不稳定，应急资金就应该越厚实。这不是浪费------它本质上是一笔''反脆弱''的投资，让你在收入低谷时不至于被迫贱卖长期资产，也不必因为一时的现金流紧张而接受不合理的借贷条件。
\end{itemize}

一个实用的评估方法是：问问自己，如果明天突然失去所有收入来源，你需要多长时间才能找到一份维持基本生活的新收入？把那个时间乘以两倍，就是你应该准备的应急资金月数。因为现实中找工作往往比预期的时间更长，而且你不太可能在心态最焦虑的时候做出最好的职业决策。

应急资金应该放在哪里？放在收益率最高的地方吗？不是。应急资金的核心要求是\textbf{安全和流动性}------你需要它的时候，它必须在，而且不能缩水。银行活期存款和货币基金是两种最常见的存放方式。它们的收益不高，但这是它们应该担任的角色------就像你不会让消防员去帮你送外卖，每一样工具都有它本来的分工。

\subsection{哪些钱不能进入股票市场}\label{ux54eaux4e9bux94b1ux4e0dux80fdux8fdbux5165ux80a1ux7968ux5e02ux573a}

有了应急资金之后，还有一笔钱同样不能用来投资：那些在未来较短时间里有明确用途的资金。

如果你计划明年买房，首付款不应该放在股票或股票型基金里。如果你下半年要交一笔大额学费，这笔钱同样不应该进入市场。如果你在攒一笔一年内必须使用的医疗费用，也不应该想着靠投资来''增值''一点。

原因很简单：短期市场的走向是完全不可预测的。某一年的市场可能涨
30\%，也可能跌
30\%。如果你刚好计划用钱的那一年碰到了下跌，你就没有足够的时间等待市场恢复------你只能被迫卖出，把浮亏变成真实的亏损。

一个简单的原则是：\textbf{五年内确定需要用到的钱，不适合投入到高波动的权益类资产中。}
如果期限在三到五年之间，可以考虑债券或偏债类的稳健产品。如果期限在三年以内，最安全的方式就是银行存款或货币基金。这不是保守，这是让每笔钱的期限和目标匹配。

\subsection{保险不是投资}\label{ux4fddux9669ux4e0dux662fux6295ux8d44}

还有一个经常被混为一谈的话题：保险和投资的关系。

保险（insurance）的本质是风险转移------你用一笔确定的保费，把未来某种小概率但大损失的风险转移给保险公司。医疗险、重疾险、意外险、定期寿险，都属于这一类。它们的功能是防止你在人生的重要阶段因为一场意外而陷入财务困境。

投资型保险产品（比如万能险、分红险、投资连结险）则试图把保障和投资功能捆绑在一起。听起来很方便，但实际上这类产品通常既不如纯保障型保险来得实惠，也不如你自己直接投资来得透明和灵活。保障和投资是两件完全不同的事，把它们混在一起，往往两头都做不好。

对于大多数年轻投资者来说，一个更清晰的策略是：用费用较低的纯保障型保险覆盖重大风险（医疗、重疾、意外、定期寿险），用自己构建的投资组合来实现财富增值。前者是防守，后者是进攻------两者各司其职，而不是让一个产品试图同时做两件它都不擅长的事。

\begin{center}{\color{dh-bronze}\rule{0.2\textwidth}{0.5pt}}\end{center}

现在，你的高息负债已经清零（或者正在清零），你的应急资金已经准备到位，你也知道了哪些钱不适合进入市场。你就像盖好了一座房子的地基------虽然还看不到上面的建筑，但你知道，这座房子不会轻易倒塌。

接下来，我们开始正式讨论投资中最核心的概念之一：复利。

它比你想象中更强大，也比你想象中更需要耐心。

\partdivider{第二部分 · 理解收益、复利与风险}

\newpage

\chapter{复利究竟如何工作}\label{ux590dux5229ux7a76ux7adfux5982ux4f55ux5de5ux4f5c}

\begin{center}{\color{dh-bronze}\rule{0.2\textwidth}{0.5pt}}\end{center}

\begin{dhquote}
\textbf{复利不是魔法，它只是一道被时间放大了的数学题。}
\end{dhquote}

\begin{center}{\color{dh-bronze}\rule{0.2\textwidth}{0.5pt}}\end{center}

如果说投资世界里有一个概念，每个人都知道它很重要，但很少有人真正感受过它的力量------那一定是复利（compounding）。大多数人对复利的理解停留在''利滚利''三个字上，但很少有人动手算过一笔账：三十年、四十年之后，复利到底能带来多大的差别。

\subsection{复利的数学，其实只有一行}\label{ux590dux5229ux7684ux6570ux5b66ux5176ux5b9eux53eaux6709ux4e00ux884c}

复利的公式非常简单：

\textbf{FV = PV × (1 + r)\^{}n}

其中，FV 是未来的终值（future value），PV 是现在的本金（present
value），r 是每一期的收益率，n
是期数。这个公式的意义在于：你的收益不仅仅来自最初的本金，每一次赚到的钱也会在下一期继续产生新的收益。这就是''利滚利''的本质。

来看一个简单的例子。假设你有一笔 10,000 元的本金，年化收益率为
7\%（教学假设，不代表未来收益）。如果只靠单利（simple
interest）------也就是每年只对本金计算收益、赚到的收益不再投入------三十年后你会得到
10,000 + 10,000 × 7\% × 30 = 31,000 元。其中收益部分是 21,000 元。

但如果是复利------每年的收益也加入本金继续滚动------三十年后你的 10,000
元会变成 10,000 × (1 + 7\%)\^{}30 ≈ 76,123 元。收益部分约为 66,123 元。

同样是 10,000 元本金，单利和复利的最终差距是 45,000
多元------超过本金的四倍。而这还只是一次性投入的情况。如果每个月持续追加投入，差距会更大。

\begin{center}{\color{dh-bronze}\rule{0.2\textwidth}{0.5pt}}\end{center}

\subsection{定投下的复利：每个月存五百元，三十年之后}\label{ux5b9aux6295ux4e0bux7684ux590dux5229ux6bcfux4e2aux6708ux5b58ux4e94ux767eux5143ux4e09ux5341ux5e74ux4e4bux540e}

现实生活中，大多数人的投资不是一次性投入一笔大钱，而是每个月从工资中拿出一部分持续投入------也就是定投。在定投场景下，复利的公式稍微复杂一些，但原理完全一样：每一笔投入都在为你工作，赚到的钱也在继续为你工作。

假设你每个月投资 500 元，年化收益率仍为
7\%（教学假设），连续投入三十年。三十年间你总共投入了 500 × 12 × 30 =
180,000 元。在复利的作用下，三十年后你的账户里大约会有 610,000
元。其中约 430,000 元来自收益，超过总投入的两倍还多。

如果你能够每个月投入 1,000 元，同样的条件下，三十年后大约是 122
万元。投入了 36 万元，收益约 86 万元。

如果你每个月只能投入 200 元，三十年后的结果大约是 24.4 万元。投入了 7.2
万元，收益约 17.2 万元。

这些数字当然依赖那个 7\% 的假设。没有人能保证未来三十年一定能获得 7\%
的年化回报------可能更高，也可能更低。但复利给我们的启示并不依赖于某一个具体的数字：\textbf{在投资中，时间是最可靠的放大器。投入金额越大、时间越长、收益率越稳定，复利的效果就越显著。}

\begin{center}{\color{dh-bronze}\rule{0.2\textwidth}{0.5pt}}\end{center}

\subsection{早十年开始，差距有多大}\label{ux65e9ux5341ux5e74ux5f00ux59cbux5deeux8dddux6709ux591aux5927}

这是复利最令人震撼的地方。

假设有两个人------花有财（小名花花）和小李。花花从 25 岁开始，每个月投资
500 元，年化收益率 7\%，坚持到 60 岁。小李从 35
岁开始，其他条件完全一样------同样的金额、同样的收益率、同样到 60 岁。

花花总共投入了 210,000 元（25 岁到 60 岁共 35 年），最终账户约 90 万元。

小李总共投入了 150,000 元（35 岁到 60 岁共 25 年），最终账户约 41 万元。

花花只比小李多投入了 60,000 元，但最终多出了约 50 万元。这 50
万元的差距，不是因为花花更努力、更聪明、或者收入更高------仅仅是因为他早了十年开始。那十年里，他投入的每一分钱都有更多的时间去工作和生长。

有一个经验法则可以帮助你快速估算复利的效果，叫做 \textbf{72 法则（Rule
of 72）}。用 72
除以年化收益率，得到的结果就是本金翻倍大约需要的年数。比如年化收益率
7\%，72 ÷ 7 ≈ 10.3，意味着大约十年你的本金可以翻一倍。如果收益率是
6\%，则需要大约十二年翻倍；如果收益率是 10\%，只需要大约七年。72
法则是一个非常粗略的估算工具，但它足够直观地揭示了收益率和时间的互动关系。

顺带一提：这个法则常被误传为爱因斯坦的发现，但它的来源更早------最早可以追溯到意大利数学家卢卡·帕乔利（Luca
Pacioli）在 1494 年出版的《算术、几何、比与比例概要》（Summa de
arithmetica）中的记载。流传最广的版本是「72
法则」这个名称，不过无论它出自谁手，背后的数学道理才是真正有用的部分。

\begin{center}{\color{dh-bronze}\rule{0.2\textwidth}{0.5pt}}\end{center}

\subsection{复利为什么前期慢、后期快}\label{ux590dux5229ux4e3aux4ec0ux4e48ux524dux671fux6162ux540eux671fux5feb}

很多人在开始投资的头几年会感到失望。明明已经在坚持定投了，为什么账户的增长看起来一点也不令人兴奋？

这是因为复利的增长并不是一条直线，而是一条不断加速上弯的曲线。在最初的十年，你的账户里大部分是自己的投入，收益只占一小部分，看起来确实没什么特别。但到了第二个十年，前十年积累的收益开始产生自己的收益，账户的增长速度明显加快。到了第三个、第四个十年，收益的收益、收益的收益的收益------一层一层叠加起来，增长速度会远远超过你当初的想象。

如果你把复利过程拆开来看：第 1 年你投入的 500 元，到第 30 年会变成约 500
× (1+7\%)\^{}29 ≈ 3,800 元；而第 29 年投入的 500 元，到第 30
年只会变成约 535 元。同样是 500 元，早投入 29
年的那一笔，最终价值是晚投入那笔的七倍。这就是时间的价值------它免费，也公平，但它不会等人。

\subsection{什么会破坏复利}\label{ux4ec0ux4e48ux4f1aux7834ux574fux590dux5229}

复利最大的敌人有两个：频繁交易和中断投资。

如果你在市场中不停地买进卖出，每一次交易都意味着你打断了复利的积累过程。而且每次交易都可能产生佣金和税费，这些成本看起来微不足道，但几十年下来，累积的摩擦成本同样惊人。

更大的破坏来自中断投资。如果在市场大跌的时候，你因为恐惧而卖出了所有持仓，那么你不仅将浮亏变成了真实亏损，还放弃了在市场恢复时重新参与复利的机会。复利不会对中途下车的人网开一面------它只奖励那些全程在场的人。

\begin{center}{\color{dh-bronze}\rule{0.2\textwidth}{0.5pt}}\end{center}

复利不是一种能让你暴富的技巧，也不是一种神秘的力量。它只是一道简单的数学题，被时间拉长了而已。当复利遇到几十年，那几十年的每一天看上去都很普通。但当所有普通的日子叠加在一起，最终的结果往往令人难以置信。

很多人高估了一年能做的事情，却低估了三十年能带来的变化。这就是复利给我们的终极启示。

我自己刚开始理财的那几年，也犯过同样的错误------总想在最短的时间里看到账户翻倍，结果追涨杀跌，收益反而不如安安静静地放着。后来我才慢慢想明白：复利这件事，难的从来不是那个百分之几的数学，而是你在头十年里，看着账户不紧不慢地增长，还愿不愿意继续往里面放钱。时间会给耐心的人发工资，只是它发工资的方式，从来不声张。

\newpage

\chapter{通货膨胀如何侵蚀购买力}\label{ux901aux8d27ux81a8ux80c0ux5982ux4f55ux4fb5ux8680ux8d2dux4e70ux529b}

\begin{center}{\color{dh-bronze}\rule{0.2\textwidth}{0.5pt}}\end{center}

\begin{dhquote}
\textbf{现金看似安全，但在时间面前，安全本身就是风险。}
\end{dhquote}

\begin{center}{\color{dh-bronze}\rule{0.2\textwidth}{0.5pt}}\end{center}

很多人觉得，把钱放在银行账户里是最保守、最安全的做法。从某个角度说，这没错------你不会在某一天打开账户发现数字变少了。但如果把时间拉长到十年、二十年，你就会发现，虽然账户上的数字没有变少，但这些数字能买到的东西，却在慢慢变少。

这就是\textbf{通货膨胀（inflation）}------物价普遍持续上涨，导致货币的购买力下降。

\subsection{一百元的三十年}\label{ux4e00ux767eux5143ux7684ux4e09ux5341ux5e74}

以一组具体数据来感受这种侵蚀。根据国家统计局数据，1995
年全国城镇居民人均可支配收入为 4,283 元/年（约 357
元/月），人均消费支出为 3,537 元/年（约 295 元/月）。到了 2025
年，城镇居民人均可支配收入达到 56,502 元/年（约 4,709
元/月），人均消费支出达到 35,869 元/年（约 2,989
元/月）------名义金额分别增长了约 12 倍和 9 倍。但如果用 CPI
剔除通胀因素，实际购买力的增长远没有数字看起来那么显著。日常消费品中，1995
年一碗普通面条大约 2---3 元，2025 年同样的面条大约 15---25 元；1995
年公交票价 0.5---1 元，2025 年普遍 2---4
元。三十年间的价格变迁，正是通货膨胀在日常生活中留下的印记。

（数据来源：国家统计局年度统计公报，1995 年与 2025
年城镇居民人均可支配收入与消费支出数据；食品与交通价格参考公开物价记录。）

更直观的方式是看通货膨胀如何影响购买力。假设每年的通货膨胀率是
3\%。你有一笔 10
万元的存款，每年什么都不做，只是放在那里。第一年后，这笔钱的''名义金额''仍然是
10 万元，但它的''实际购买力''只相当于大约 97,000
元。十年后，实际购买力下降到约 74,000 元。二十年后，只剩约 54,000
元。三十年后，不到 40,000 元。

也就是说，你什么都没做，三十年后你的钱缩水了六成。这不是因为银行偷了你的钱，而是因为物价一直在悄悄上涨，而你手中的现金没有能力跟上这种上涨。

\begin{center}{\color{dh-bronze}\rule{0.2\textwidth}{0.5pt}}\end{center}

\subsection{CPI
不等于你感受到的通胀}\label{cpi-ux4e0dux7b49ux4e8eux4f60ux611fux53d7ux5230ux7684ux901aux80c0}

国家统计局每个月会发布 \textbf{CPI（Consumer Price
Index，居民消费价格指数）}，这个数字代表的是''一篮子商品和服务''价格的平均变化。但每个人的消费结构不同，所以每个人感受到的通胀也不同。

如果你不抽烟、不喝酒，那么烟酒价格上涨对你没有影响。如果你需要租房而不是已经拥有自住房，那么房租上涨对你的冲击远大于对有房者。如果你有孩子在上学，教育支出的增长可能才是你最大的压力来源。如果你或家人有医疗需求，医疗费用的增幅往往远超
CPI 平均数。

因此，在评估通胀对个人财务的影响时，与其盯着新闻里的 CPI
数据，不如回顾一下自己家庭的真实支出变化。三年前你每个月的开销是多少？今天是多少？这个个人化的''家庭通胀率''，才是对你真正有意义的数字。

\begin{center}{\color{dh-bronze}\rule{0.2\textwidth}{0.5pt}}\end{center}

\subsection{现金的安全与购买力风险}\label{ux73b0ux91d1ux7684ux5b89ux5168ux4e0eux8d2dux4e70ux529bux98ceux9669}

每一种资产都有风险，只不过风险的形式不同。股票的风险是价格波动------今天可能涨，明天可能跌，账面价值忽上忽下。债券的风险是利率变化和信用违约。而现金的风险，就是通货膨胀导致的购买力下降。

很多人只看到了波动的风险，却忽视了通胀的风险------因为前者看得见，后者看不见。股票跌了
20\%，你打开账户就能看到。但现金的购买力每年缩水
3\%，你看不到任何红色的数字，只有等你去超市、去交房租、去付孩子的学费时，才会隐隐约约感觉到：``钱好像越来越不经花了。''

这就是为什么，把所有积蓄都放在现金或活期存款里，看似安全，实则是在承担一种''安静的亏损''。长期来看，现金无法为你保值，更无法为你增值。

\subsection{如何计算实际收益率}\label{ux5982ux4f55ux8ba1ux7b97ux5b9eux9645ux6536ux76caux7387}

当你看到一项投资的名义收益率时，不能直接用它来和通货膨胀率比较。你需要计算\textbf{实际收益率（real
rate of return）}。

有一个简化的算法，叫做\textbf{费雪方程（Fisher Equation）}：实际收益率 ≈
名义收益率 - 通货膨胀率。

如果一只基金去年的名义收益率是 5\%，而同期的通货膨胀率是
3\%，那么你的实际收益率大约是 2\%。这 2\%
才是你购买力真正的增长。如果名义收益率是 3\%，通胀率也是
3\%，那么你的实际收益率为零------你的购买力和一年前完全一样。如果名义收益率低于通胀率，那么你的购买力实际上在下降，尽管账户里的数字看起来在增加。

因此，在评估任何一项投资的时候，不能只看它''赚了多少''，还要看它''跑赢了通胀多少''。一个年化
3\% 的银行理财产品，在通胀率为 2\% 的时候看起来还不错；但如果通胀上升到
4\%，它实际上在让你悄悄变穷。

\begin{center}{\color{dh-bronze}\rule{0.2\textwidth}{0.5pt}}\end{center}

通货膨胀不会在一天之内改变你的生活，但它会在几十年里重新定义你手中每一分钱的价值。那些今天看起来''足够''的积蓄，如果没有任何增值，十年后可能只能覆盖一半的生活。这就是为什么我们需要投资------不是为了追求刺激，而是为了让我们的财富至少跟上物价的步伐，最好还能跑赢它，让未来拥有更多实实在在的购买力。

\begin{figure}
\centering
\pandocbounded{\includegraphics[keepaspectratio]{ch06_inflation.png}}
\caption{1995-2025年中国CPI通胀走势}
\end{figure}

接下来的问题是：既然我们愿意承受一定的风险去追求更高的回报，那么风险到底是什么？它比大多数人想象的更复杂。

\newpage

\chapter{风险不只是价格下跌}\label{ux98ceux9669ux4e0dux53eaux662fux4ef7ux683cux4e0bux8dcc}

\begin{center}{\color{dh-bronze}\rule{0.2\textwidth}{0.5pt}}\end{center}

\begin{dhquote}
\textbf{真正的风险，不是资产的价格在变化，而是在你不了解的时候做出了决定。}
\end{dhquote}

\begin{center}{\color{dh-bronze}\rule{0.2\textwidth}{0.5pt}}\end{center}

``这只股票风险太大了。''``这个理财应该没有风险。''``我不想承担任何风险。''------在日常对话中，人们习惯把风险简单地等同于''可能会亏钱''。但在投资的世界里，风险是一个远比''亏钱''复杂得多的概念。如果不能对风险建立一个更完整的认识，就很难做出理性的投资决策。

\subsection{波动不是风险，永久损失才是}\label{ux6ce2ux52a8ux4e0dux662fux98ceux9669ux6c38ux4e45ux635fux5931ux624dux662f}

假���你今天买入了一只基金。第二天，市场整体下跌，你的账户亏了
3\%。第三天又跌了 2\%。你是否面临了''风险''？

大部分人会回答''是''。但让我们看看另一种情况：你没有卖出，一年之后这只基金涨了回来，还创了新高。那之前那两天
5\% 的下跌，到底是不是''真正的损失''？

\textbf{波动（volatility）}
指的是资产价格上下起伏的幅度。它是市场正常运转的一部分------就像海浪有起有落，股票和基金也会随着每天的交易而涨跌。波动让人感到不舒服，但它不等于永久性的损失。

\textbf{永久损失（permanent loss of capital）}
则是另外一回事。如果你的投资标的本身发生了根本性的恶化------企业失去了核心竞争力，破产清算；债券发行人违约，无法偿还本息------那么即使你继续持有，资产的价值也可能永远无法恢复。这才是投资中真正应该警惕的风险。

投资大师霍华德·马克斯（Howard
Marks）在他的著作《投资最重要的事》中将风险定义为''永久损失资本的可能性''，同时明确指出，风险并不等于波动。这一区分是理解投资风险的关键起点。

理解这一点之后，面对市场的日常波动，你可以多问自己一个问题：我持有的资产本身变差了吗？还是只是市场情绪变了？如果答案是后者，那今天的下跌就只是一个海浪，而不是海啸。

\textbf{主要市场历史最大回撤与恢复时间：}

\begin{longtable}[]{@{}
  >{\raggedright\arraybackslash}p{(\linewidth - 8\tabcolsep) * \real{0.2000}}
  >{\raggedright\arraybackslash}p{(\linewidth - 8\tabcolsep) * \real{0.2000}}
  >{\raggedright\arraybackslash}p{(\linewidth - 8\tabcolsep) * \real{0.2000}}
  >{\raggedright\arraybackslash}p{(\linewidth - 8\tabcolsep) * \real{0.2000}}
  >{\raggedright\arraybackslash}p{(\linewidth - 8\tabcolsep) * \real{0.2000}}@{}}
\toprule\noalign{}\rowcolor{dh-pink!25}
\begin{minipage}[b]{\linewidth}\raggedright
市场指数
\end{minipage} & \begin{minipage}[b]{\linewidth}\raggedright
最大回撤幅度
\end{minipage} & \begin{minipage}[b]{\linewidth}\raggedright
发生时期
\end{minipage} & \begin{minipage}[b]{\linewidth}\raggedright
见底天数
\end{minipage} & \begin{minipage}[b]{\linewidth}\raggedright
恢复至前高所需时间
\end{minipage} \\
\midrule\noalign{}
\endhead
\bottomrule\noalign{}
\endlastfoot
标普500（S\&P 500） & −56.8\% & 2007.10---2009.03 & 517 天 & 约 4.0
年（至 2013.03） \\
标普500（S\&P 500） & −49.1\% & 2000.03---2002.10 & 929 天 & 约 7.2
年（至 2007.05） \\
标普500（S\&P 500） & −33.9\% & 2020.02---2020.03 & 33 天 & 约 5
个月（至 2020.08） \\
沪深300（CSI 300） & −72.3\% & 2007.10---2008.11 & 约 13 个月 & 约
7---10 年（含分红再投资至 2015 年才实质超越前高） \\
沪深300（CSI 300） & −46.1\% & 2015.06---2016.02 & 约 8 个月 & 约 4.5
年（至 2020.12） \\
\end{longtable}

从上表可以看出两个规律：第一，回撤幅度越大，恢复时间往往越长------``跌得越快、反弹越快''（如
2020 年疫情冲击）只是一个例外而非规律；第二，A
股的回撤幅度通常大于美股，但恢复窗口也更不可预测。因此，最大回撤不仅是选择投资标的时的参考数据，更应该用来校准自己投资组合的权益敞口------如果你的组合中有
60\% 的股票，那就意味着在历史上最糟糕的年份，你账面上可能承受
30\%---40\%
以上的浮亏。如果你不能接受这个数字出现在自己的账户里，那么你应该在买入之前就降低股票的比例，而不是在下跌之后被迫做出情绪化的决定。

（数据来源：标普500 回撤与恢复 --- History of Market / S\&P Dow Jones
Indices；沪深300 回撤 --- Bloomberg / Wind
终端，按价格指数计算回撤幅度，恢复时间含分红再投资。）

\begin{center}{\color{dh-bronze}\rule{0.2\textwidth}{0.5pt}}\end{center}

\subsection{风险的九种面貌}\label{ux98ceux9669ux7684ux4e5dux79cdux9762ux8c8c}

除了市场价格波动，投资者实际上还面临着许多不同层面的风险。它们并不总是同时出现，但每一种都可能对你的投资组合产生影响。

\textbf{信用风险（credit
risk）}，指的是债券或其他债务工具的发行人无法按时支付利息或偿还本金的风险。国债的信用风险通常被认为接近于零，而企业债的信用风险取决于发行企业的财务状况。评级机构（如标普、穆迪、惠誉）会为债券提供信用评级，但评级并非保证------历史上也出现过评级下调甚至违约的事件。

\textbf{利率风险（interest rate
risk）}，主要影响债券和债券基金。当市场利率上升时，已发行的债券价格会下跌。原因很简单：新发行的债券利率更高，旧债券的相对吸引力下降，价格自然回落。我们将在第十章详细讨论这个机制。

\textbf{流动性风险（liquidity
risk）}，指的是无法在需要的时候以合理价格快速卖出资产。比如一套房产，可能挂了半年还卖不掉；一只冷门的小盘股，可能因为成交清淡而出现大幅折扣才能脱手。流动性风险在极端市场环境下尤其明显------当所有人都在抛售时，你会发现买家突然消失了。

\textbf{汇率风险（currency
risk）}，是在投资海外资产时必须考虑的因素。如果你用人民币买入了一只美元计价的资产，即使资产本身的价格没有变化，人民币对美元汇率的波动也会影响你的实际收益。汇率升值时，你的海外资产折算成人民币会缩水；汇率贬值时，则会放大你的收益。关于这一点，我们在全球投资的相关章节中还会展开。

\textbf{集中度风险（concentration
risk）}，来自过度集中于某一种资产、某一个行业或某一个国家。持有十只科技股，并不等于''分散投资''------它们面对的风险因素高度相似。真正的分散，是在不同风险来源之间分配资金。

\textbf{杠杆风险（leverage
risk）}，来自使用借来的钱进行投资。杠杆可以放大收益------如果投资上涨
10\%，两倍杠杆可以让你的收益变成 20\%。但如果投资下跌
10\%，两倍杠杆意味着损失
20\%。杠杆是双刃剑，对于新手投资者来说，最安全的建议是：在完全理解杠杆的成本与风险之前，不要使用它。

\textbf{期限错配风险（maturity mismatch
risk）}，发生在投资期限与资金需求不匹配的时候。如果你用一笔明年就要使用的资金投资了二十年期限的长期债券，当市场利率上升、债券价格下跌时，你无法等到债券到期收回本金------不得不在亏损时卖出。

\textbf{托管与平台风险（custody and platform
risk）}，指的是投资资金存放的平台或托管机构出现问题的情况。选择合规、受监管的金融机构和交易平台，是降低这类风险最基本也是最重要的一步。不在本书的具体推荐范围内，但读者应当自行确认所使用的平台是否持有合法牌照。

\begin{center}{\color{dh-bronze}\rule{0.2\textwidth}{0.5pt}}\end{center}

\subsection{亏损之后，需要涨更多才能回本}\label{ux4e8fux635fux4e4bux540eux9700ux8981ux6da8ux66f4ux591aux624dux80fdux56deux672c}

很多人低估了亏损的恢复难度。下跌和上涨在数学上并不是对称的。

\begin{longtable}[]{@{}ll@{}}
\toprule\noalign{}\rowcolor{dh-pink!25}
亏损幅度 & 需要上涨多少才能回本 \\
\midrule\noalign{}
\endhead
\bottomrule\noalign{}
\endlastfoot
-10\% & +11.1\% \\
-20\% & +25.0\% \\
-30\% & +42.9\% \\
-40\% & +66.7\% \\
-50\% & +100.0\% \\
-60\% & +150.0\% \\
-70\% & +233.3\% \\
\end{longtable}

这张表最值得关注的是：亏损超过 50\%
以后，回本所需的涨幅变得极其陡峭。亏损
70\%，需要涨两倍多才能回本------这在短期内几乎是不可能的。这就是为什么控制大幅回撤比追求短期高收益更为重要。一次致命的亏损，可能耗尽你很多年积累的所有收益。

\begin{center}{\color{dh-bronze}\rule{0.2\textwidth}{0.5pt}}\end{center}

风险管理不是要消灭风险，而是要理解你承担的是什么风险、为什么承担它，以及这个风险是否在你的可承受范围之内。用投资术语来说，这叫''为风险定价''------你不是在回避风险，而是在确保你承担的每一份风险，都有合理的机会带来相应的回报。

\begin{figure}
\centering
\pandocbounded{\includegraphics[keepaspectratio]{ch07_drawdowns.png}}
\caption{主要市场历史最大回撤对比}
\end{figure}

接下来的问题是：你怎么知道，自己到底能承受多大的风险？

\newpage

\chapter{你真正能承受多少风险}\label{ux4f60ux771fux6b63ux80fdux627fux53d7ux591aux5c11ux98ceux9669}

\begin{center}{\color{dh-bronze}\rule{0.2\textwidth}{0.5pt}}\end{center}

\begin{dhquote}
\textbf{别人能承受 30\%
的回撤，不意味着你也可以。风险承受力，是一道关于你自己的算术题。}
\end{dhquote}

\begin{center}{\color{dh-bronze}\rule{0.2\textwidth}{0.5pt}}\end{center}

在投资中，有一个很容易被忽略但极其重要的区别：你能承受多少风险，和你愿意承受多少风险，以及你需要承受多少风险------这是三件不同的事。把它们混在一起，是很多人投资失败的根本原因。

\subsection{你能承受多少风险}\label{ux4f60ux80fdux627fux53d7ux591aux5c11ux98ceux9669}

\textbf{风险承受能力（risk capacity）}
是一个客观的财务指标，与你的情绪和意愿无关。它取决于几个硬性条件。

第一个是时间。你的投资期限越长，你承受短期波动的能力就越强。一个二十五岁的年轻人，距离退休还有四十年，有足够的时间等待市场从任何一次下跌中恢复。一个五十五岁、距离退休只有五年的人，即使心态比年轻人更沉稳，从财务角度看也无法承受同样的波动------因为他没有足够的时间来等待恢复。

第二个是人力资本。我们在第二章提到过这个概念。一个人如果还处于职业生涯的上升期，未来还有几十年的工资收入，即使投资账户暂时缩水
30\%，也不至于影响生活。但一个已经退休、不再有劳动收入的人，投资账户就是他的全部收入来源，同样的
30\% 下跌对他来说是完全不同的量级。

第三个是家庭现金流。无论你的投资账户里有多少钱，如果你的每个月都在亏损------工资不够还房贷和日常支出------那么你就没有真正的风险承受能力。只有正的现金流，才能为你提供投资中''熬过低谷''的底气。

\subsection{你愿意承受多少风险}\label{ux4f60ux613fux610fux627fux53d7ux591aux5c11ux98ceux9669}

\textbf{风险承受意愿（risk tolerance）}
则是一个主观的心理指标。有些人看到账户跌了 5\% 就睡不着觉，有些人跌了
20\%
依然能吃能睡。这没有对错------每个人的心理特质不同，成长经历不同，对风险的感受也不同。

问题是，很多人对自己的风险承受意愿存在严重的误判。在牛市中，市场天天涨，每个人都觉得自己可以承受很大的波动------``跌
20\%
我也能扛住''。但当熊市真的到来、媒体上全是悲观消息、身边的人都在割肉离场的时候，当初的信心往往会迅速瓦解。一个人在牛市中对自己风险偏好的判断，往往和熊市中的实际表现大相径庭。

因此，评估风险承受意愿不能只是在风和日丽的时候想当然。你应该回顾一下自己过去的经历：上一次面对真正的金钱损失是什么感觉？你有没有因为账户下跌而改变过自己的行为？你有没有在市场最恐慌的时候卖过资产？这些真实发生的事情，比任何问卷调查都更诚实。

\subsection{你需要承受多少风险}\label{ux4f60ux9700ux8981ux627fux53d7ux591aux5c11ux98ceux9669}

\textbf{风险承受需求（risk required）}
来自你的投资目标。如果你希望在三十年后积累足够的养老金，而当前的储蓄率只能支撑一个温和的退休生活，那么你可能需要承担更高的风险来追求更高的长期回报------否则你的目标可能无法实现。

但这是一个需要格外谨慎的领域。为了追求目标而承担超出自己能力的风险，是目前投资中最常见的错误之一，也是最危险的行为偏差之一。如果你的投资目标需要一个
15\% 的年化回报率才能实现，而市场在长期内只能提供 6\% 到
8\%，那么问题不在于你的风险偏好------问题在于你的目标设定得不够现实，或者你的储蓄率需要提高。

投资的正确顺序是：先根据你的风险承受能力确定你能安全承担多大的波动；再根据你的风险承受意愿微调（如果你真的无法忍受波动，可以适当降低高风险资产的比例，但要接受长期回报可能降低的结果）；最后，如果发现风险承受能力无法支撑你的投资目标，那么调整的应该是目标（推迟退休、降低预期支出）或储蓄率（增加每月投入），而不是冒险加大风险敞口。

\begin{center}{\color{dh-bronze}\rule{0.2\textwidth}{0.5pt}}\end{center}

\subsection{一份简单的风险自测}\label{ux4e00ux4efdux7b80ux5355ux7684ux98ceux9669ux81eaux6d4b}

以下是一个简化的风险自测框架。它不是精确的科学测评，但可以帮助你对自己有一个基本的判断。

\textbf{1. 你的投资期限大约是多少年？}

\begin{itemize}
\tightlist
\item
  超过 20 年：你可以承受较大的短期波动，适合较高比例的权益类资产。
\item
  10---20 年：你的投资期限仍然较长，但需要考虑中期目标（如购房、教育）。
\item
  5---10 年：你的投资期限中等，波动较大的资产比例应当适当降低。
\item
  不到 5 年：你的投资期限较短，不应承担过多的权益风险。
\end{itemize}

\textbf{2. 你的收入稳定性如何？}

\begin{itemize}
\tightlist
\item
  非常稳定且处于增长期：你可以承担较高的风险敞口。
\item
  比较稳定，但存在一定的不确定性：应当保留适当的缓冲空间。
\item
  波动较大或行业风险较高：应当更保守地配置风险资产。
\end{itemize}

\textbf{3. 你的家庭现金流状况如何？}

\begin{itemize}
\tightlist
\item
  每月有稳定的正向现金流：可以按照投资计划正常执行。
\item
  收支基本平衡：应当优先建立应急资金，投资以保守为主。
\item
  每月入不敷出：需要先解决现金流问题，再考虑投资。
\end{itemize}

\textbf{4. 如果市场下跌 20\%，你的第一反应是？}

\begin{itemize}
\tightlist
\item
  继续按计划定投，甚至希望抓住机会：你可能有较高的风险承受意愿。
\item
  有些不安，但不会做出大的改变：你的风险承受意愿处于中等水平。
\item
  想要卖出止损：你可能倾向于低风险偏好，需要为投资组合增加稳定资产。
\end{itemize}

\textbf{5. 你是否曾经因为市场波动而卖出过投资资产？}

\begin{itemize}
\tightlist
\item
  从来没有：你对波动有较强的心理弹性，或投资经历较短。
\item
  偶尔有过，但后来恢复了投资：你在学习和调整中。
\item
  经常因为害怕亏损而卖出：你可能不适合高比例的权益类资产，这不代表失败------了解自己，恰恰是投资中最重要的一步。
\end{itemize}

\subsection{将自测结果映射到资产配置}\label{ux5c06ux81eaux6d4bux7ed3ux679cux6620ux5c04ux5230ux8d44ux4ea7ux914dux7f6e}

如果你对上面五个问题做了一个大概的统计，可以用以下框架给自己一个初始参考方向。

\begin{longtable}[]{@{}
  >{\raggedright\arraybackslash}p{(\linewidth - 6\tabcolsep) * \real{0.2500}}
  >{\raggedright\arraybackslash}p{(\linewidth - 6\tabcolsep) * \real{0.2500}}
  >{\raggedright\arraybackslash}p{(\linewidth - 6\tabcolsep) * \real{0.2500}}
  >{\raggedright\arraybackslash}p{(\linewidth - 6\tabcolsep) * \real{0.2500}}@{}}
\toprule\noalign{}\rowcolor{dh-pink!25}
\begin{minipage}[b]{\linewidth}\raggedright
如果你的答案大多是
\end{minipage} & \begin{minipage}[b]{\linewidth}\raggedright
风险类型
\end{minipage} & \begin{minipage}[b]{\linewidth}\raggedright
建议权益类资产比例
\end{minipage} & \begin{minipage}[b]{\linewidth}\raggedright
说明
\end{minipage} \\
\midrule\noalign{}
\endhead
\bottomrule\noalign{}
\endlastfoot
第一个选项（A） & 进取型 & 60\%--70\% &
你的风险承受能力和意愿都较高，可以考虑以权益类资产为主的配置。 \\
第二个选项（B） & 稳健型 & 50\%--60\% &
你可以在追求增长的同时保留一定比例的稳健资产作为缓冲。 \\
第三个选项（C） & 保守型 & 30\%--40\% &
你可能需要更多债券和现金类资产来稳定组合，接受长期回报相对温和的结果。 \\
第四个选项（D） & 优先修复���务基础 & 10\% 或更低 &
你可能需要先关注储蓄率、应急资金和债务管理，暂不适合高比例权益投资。 \\
\end{longtable}

这个映射只是一个起点，不是公式。你在不同问题上的答案可能不是一致的------比如投资期限很长（A）但恐惧下跌（C）。这种矛盾恰恰说明你需要找到一个让自己既舒适又能坚持的平衡点，而不是逼迫自己接受一个''理论上最优''的方案。

\begin{center}{\color{dh-bronze}\rule{0.2\textwidth}{0.5pt}}\end{center}

这些问题的答案没有正确或错误之分。它们的目标不是给你打分，而是帮助你诚实地面对自己。投资中最大的风险，常常不是市场的下跌，而是在不了解自己的情况下，做出了超出自己能力范围的决定。

在下一章，我们将进入一个更具体的话题：当你清楚了风险的全貌之后，市场上到底有哪些工具可以用来配置你的资金。

\partdivider{第三部分 · 钱到底可以放在哪里}

\newpage

\chapter{现金、存款与货币类工具}\label{ux73b0ux91d1ux5b58ux6b3eux4e0eux8d27ux5e01ux7c7bux5de5ux5177}

\begin{center}{\color{dh-bronze}\rule{0.2\textwidth}{0.5pt}}\end{center}

\begin{dhquote}
\textbf{现金不是无用的，它只是在等待它的任务。}
\end{dhquote}

\begin{center}{\color{dh-bronze}\rule{0.2\textwidth}{0.5pt}}\end{center}

在追求投资收益的路上，现金常常被当作反面教材------收益低、跑不赢通胀、放在账户里就是在贬值。这些批评并非全无道理：长期持有大量现金确实不是明智之举。但如果我们因此得出''现金无用''的结论，就走到了另一个极端。

现金在投资组合中承担着不可替代的角色：提供流动性、保障应急资金、平衡整体风险。就像一支球队不能只有前锋没有守门员，一个健康的财务体系也不能只有高波动资产而没有现金缓冲。

\subsection{现金的四种形态}\label{ux73b0ux91d1ux7684ux56dbux79cdux5f62ux6001}

\textbf{活期存款（demand deposit）}
是最基础的形式。随时可以存取，没有任何限制，利率也最低。它的用途是覆盖日常开支------这个月的房租、水电、吃饭、交通。你不需要从活期存款中获得收益，你只需要知道自己随时可以用。

\textbf{定期存款（time deposit / certificate of deposit）}
锁定期限，换取比活期更高的利率。三年期定期存款的利率通常高于一年期，一年期又高于活期。代价是如果提前取出，利息会按活期利率计算甚至损失全部利息。因此定期存款只适合存放确定在近期内不会动用的资金------比如三个月后要交的保费、半年后需要的装修款，而不是你的长期投资。

\textbf{货币基金（money market fund，简称货基）}
是近十年在中国普及速度最快的现金管理工具。余额宝、零钱通等产品的底层资产，本质上都是货币基金。货币基金主要投资于短期国债、银行定期存单、高等级短期债券等低风险工具，因此收益率通常高于活期存款而略低于较长期限的定期存款，但流动性远好于定期存款------大多数货基支持当天赎回、快速到账。

近年来，中国货币基金的 7
日年化收益率经历了从高位回落到低位企稳的过程。以市场上规模最大的货币基金之一天弘余额宝（000198）为例：2013---2014
年间，7 日年化收益率一度高达 4.5\%---4.8\%；2015---2018 年间回落至
2.5\%---3.9\% 区间；2019---2020 年进一步下行至
1.8\%---2.3\%；2021---2026 年间基本维持在 1.4\%---2.0\%
的低位区间。这一下降趋势背后是全球低利率环境的延续，以及国内货币政策保持宽松的结果。对投资者而言，这意味着货币基金已从''能跑赢通胀的现金管理工具''转变为''仅能略微减缓购买力侵蚀的短期停泊工具''。

（数据来源：天弘余额宝货币基金（000198）历史 7 日年化收益率，天天基金网
https://fund.eastmoney.com；中国人民银行活期存款基准利率
http://www.pbc.gov.cn。）

\textbf{短期国债类工具（short-term government bond
instruments）}，比如国债逆回购或短期国债ETF，兼具安全性和一定收益，但它们通常需要证券账户才能操作。对于初学者来说，货基和银行存款已经足够覆盖现金管理的需求，不必一开始就为了多几个点的收益而把简单的事情变复杂。

\subsection{现金在组合中的职责}\label{ux73b0ux91d1ux5728ux7ec4ux5408ux4e2dux7684ux804cux8d23}

很多人会犯一个错误：把资产配置等同于''如何分配投资资金''，而忽略了现金本身就是一种配置。

现金的职责不是赚钱，而是给你三个东西：安全感、选择权、操作空间。当你确定自己有三个月到六个月的生活费用随时可以动用时，你在面对市场下跌时的心态会完全不同------因为你不需要被迫卖出资产。当市场出现难得的低估机会时，拥有现金的人可以选择买入，而全部资金被锁在资产里的人只能旁观。

但现金也不宜过多。超过应急资金和中短期目标所需的现金，长期闲置在低收益账户里，实际上是在承担通胀风险。一个简单的原则是：应急资金和中短期目标所需资金放在现金或类现金工具中；其余的资金，根据自己的投资期限和风险承受能力，配置到能够产生长期回报的资产中。

\begin{center}{\color{dh-bronze}\rule{0.2\textwidth}{0.5pt}}\end{center}

现金很无趣，它不会给你讲令人兴奋的增长故事。但它是你所有投资中最忠实的那一部分------当市场上其他东西都在跌的时候，它依然在那里，一分不少。

\begin{figure}
\centering
\pandocbounded{\includegraphics[keepaspectratio]{ch09_money_market.png}}
\caption{货币基金7日年化收益率历史走势}
\end{figure}

在下一章，我们将讨论现金之外的第一个''可投资资产''：债券。

\newpage

\chapter{债券到底是什么}\label{ux503aux5238ux5230ux5e95ux662fux4ec0ux4e48}

\begin{center}{\color{dh-bronze}\rule{0.2\textwidth}{0.5pt}}\end{center}

\begin{dhquote}
\textbf{债券不是''更安全的股票''，它是另一类完全不同的资产。}
\end{dhquote}

\begin{center}{\color{dh-bronze}\rule{0.2\textwidth}{0.5pt}}\end{center}

如果把投资世界比作交通工具，股票是有时快有时慢的汽车------可能遇上堵车，也可能开得飞快。现金是站在地面上不动。而债券，则更像一列按照时刻表行驶的火车：速度不快，但路线基本固定，到达时间也有大致的预期。

债券在投资组合中的角色，不是追求高收益，而是提供稳定性和可预期性。理解债券，也是对利率、信用和期限结构的一次基础训练------这些概念会反复出现在后面的资产配置讨论中。

\subsection{债券的本质：借钱给别人，约定还钱的时间和方式}\label{ux503aux5238ux7684ux672cux8d28ux501fux94b1ux7ed9ux522bux4ebaux7ea6ux5b9aux8fd8ux94b1ux7684ux65f6ux95f4ux548cux65b9ux5f0f}

购买一份债券，本质上就是你把钱借给了一个实体------可能是政府（国债），可能是地方政府（地方债），也可能是企业（企业债）------对方承诺在未来某个时间归还本金，并在持有期间按照约定的利率支付利息。

债券有三个关键要素。\textbf{面值（face value / par value）}
是债券到期时你会收回的本金金额，也是发行时计算利息的基础。\textbf{票面利率（coupon
rate）} 是发行人承诺每年支付利息的比率。\textbf{到期日（maturity date）}
则是本金归还的最后期限。

举个例子：你购买了一张面值 100 元、票面利率
4\%、期限为五年的国债。按照约定，你每年会收到 4
元的利息。五年到期时，发行方归还你 100
元本金。整个过程中你的现金流入是完全确定的------只要发行方不违约。这个确定性，是股票无法提供的。

\subsection{收益率：你实际赚了多少}\label{ux6536ux76caux7387ux4f60ux5b9eux9645ux8d5aux4e86ux591aux5c11}

债券在发行之后，可以在二级市场（secondary
market）上买卖。一旦进入交易，它的市场价格就不再等于面值了，而你实际的投资收益率也不再等于票面利率。

\textbf{到期收益率（Yield to Maturity，简称 YTM）}
是衡量债券投资回报的核心指标。它综合考虑了你买入债券时支付的价格、持有期间收到的利息、以及到期时收回的本金，计算出的年化回报率。

如果你以面值（100 元）买入一只票面利率为 4\%
的债券，持有到期，你的到期收益率就是 4\%。但如果你以 95
元买入同一只债券，到期时你不仅每年收到 4 元利息，到期还会收回 100
元本金------你赚了 5 元价差。这时候你的实际年化回报就高于
4\%。反过来，如果你以 105 元买入，实际回报就低于 4\%。

这个机制揭示了一个重要的规律：\textbf{债券的收益率和价格，永远反向运动。}
价格涨了，收益率就降了；价格跌了，收益率就升了。

\subsection{利率上升，为什么债券会亏钱}\label{ux5229ux7387ux4e0aux5347ux4e3aux4ec0ux4e48ux503aux5238ux4f1aux4e8fux94b1}

这是债券投资中最容易被误解的地方。很多人以为''买了债券就不会亏钱''，但实际上，如果你在市场利率上升期间卖出手中的债券，可能亏得不少。

假设你持有一张票面利率为 3\%
的十年期国债。第二年，市场上新发行的十年期国债票面利率涨到了
5\%。现在如果有人想买你手里的债券，他为什么要接受 3\%
的利率，而不是去买新发行的 5\%
的债券？答案是他不会------除非你的债券价格降到足够低，使得实际收益率等同于新债券的水平。于是你的债券价格就会下跌。

\textbf{利率风险（interest rate risk）}
正是通过这个机制影响债券投资者。一般来说，债券的期限越长，对利率变化就越敏感。衡量这种敏感度的指标叫\textbf{久期（duration）}。粗略地说，久期越大，利率每变化
1\%，债券价格变化的幅度就越大。例如一只久期为 7
年的债券基金，当市场利率上升 1\% 时，基金净值大约会下跌 7\%。

\textbf{中国国债关键期限收益率（2026 年中，中国债券信息网）：}

\begin{longtable}[]{@{}
  >{\raggedright\arraybackslash}p{(\linewidth - 16\tabcolsep) * \real{0.1111}}
  >{\raggedright\arraybackslash}p{(\linewidth - 16\tabcolsep) * \real{0.1111}}
  >{\raggedright\arraybackslash}p{(\linewidth - 16\tabcolsep) * \real{0.1111}}
  >{\raggedright\arraybackslash}p{(\linewidth - 16\tabcolsep) * \real{0.1111}}
  >{\raggedright\arraybackslash}p{(\linewidth - 16\tabcolsep) * \real{0.1111}}
  >{\raggedright\arraybackslash}p{(\linewidth - 16\tabcolsep) * \real{0.1111}}
  >{\raggedright\arraybackslash}p{(\linewidth - 16\tabcolsep) * \real{0.1111}}
  >{\raggedright\arraybackslash}p{(\linewidth - 16\tabcolsep) * \real{0.1111}}
  >{\raggedright\arraybackslash}p{(\linewidth - 16\tabcolsep) * \real{0.1111}}@{}}
\toprule\noalign{}\rowcolor{dh-pink!25}
\begin{minipage}[b]{\linewidth}\raggedright
期限
\end{minipage} & \begin{minipage}[b]{\linewidth}\raggedright
1 年
\end{minipage} & \begin{minipage}[b]{\linewidth}\raggedright
2 年
\end{minipage} & \begin{minipage}[b]{\linewidth}\raggedright
3 年
\end{minipage} & \begin{minipage}[b]{\linewidth}\raggedright
5 年
\end{minipage} & \begin{minipage}[b]{\linewidth}\raggedright
7 年
\end{minipage} & \begin{minipage}[b]{\linewidth}\raggedright
10 年
\end{minipage} & \begin{minipage}[b]{\linewidth}\raggedright
20 年
\end{minipage} & \begin{minipage}[b]{\linewidth}\raggedright
30 年
\end{minipage} \\
\midrule\noalign{}
\endhead
\bottomrule\noalign{}
\endlastfoot
收益率 & 1.18\% & 1.22\% & 1.28\% & 1.42\% & 1.55\% & 1.62\% & 1.80\% &
1.88\% \\
\end{longtable}

作为对比，2024 年中的 10 年期国债收益率约为 2.28\%，2025 年中约为
1.85\%------利率在过去两年持续下行。这意味着新发行的国债利率越来越低，已持有的老债券（票面利率更高的）在二级市场上的价格反而上涨；但对于新进入债市的投资者来说，当前买入并持有到期的收益率已处于历史低位。

（数据来源：中国债券信息网（www.chinabond.com.cn）中债国债到期收益率曲线，2024/2025/2026
年 6 月 30 日前后最近交易日数据。）

\subsection{国债与企业债}\label{ux56fdux503aux4e0eux4f01ux4e1aux503a}

\textbf{国债（government bond / sovereign bond）}
由国家信用背书，通常被认为是最安全的债券类型。中国的国债，美国的国债（U.S.
Treasury
bonds），都享有接近''无风险''的信用评价------因为一个国家只要还能征税和印钞，就几乎不可能对本国货币计价的债务违约。但需要指出，``无风险''指的是违约风险，而不是价格风险------国债仍然面临利率风险。

\textbf{企业债（corporate bond）}
由企业发行，信用质量取决于发行企业的财务状况。为了补偿更高的信用风险，企业债通常提供比同期限国债更高的利率------这个差额叫做\textbf{信用利差（credit
spread）}。评级机构如标普（S\&P）、穆迪（Moody's）和惠誉（Fitch）会对企业债进行信用评级，AAA
是最高等级，BBB 以下通常被称为''高收益债''（high-yield
bond）或''垃圾债''（junk
bond）------顾名思义，收益更高，但违约风险也更大。

\subsection{债券基金为什么也可能亏损}\label{ux503aux5238ux57faux91d1ux4e3aux4ec0ux4e48ux4e5fux53efux80fdux4e8fux635f}

很多投资者通过债券基金（bond
fund）而不是直接购买单只债券来参与债市。债券基金持有一篮子债券，提供分散化的优势。但债券基金没有''到期日''------基金经理会不断买卖债券来维持基金的久期目标。这就意味着，债券基金的净值会随着市场利率的变化而波动，不像你直接持有单只债券到到期日那样能够锁定最终的本金回收。

在利率上升周期，债券基金出现亏损是完全正常的。这并不说明基金管理得不好------它只是在反映整个债券市场的基本规律。

\begin{center}{\color{dh-bronze}\rule{0.2\textwidth}{0.5pt}}\end{center}

债券在投资组合中的角色，用一个词来概括就是''稳定器''。它不会让你的财富快速增长，但它能在股票市场剧烈震荡的时候，帮你撑住整个组合不至于崩盘。关于它和其他资产之间的搭配，我们将在资产配置的部分详细展开。

\begin{figure}
\centering
\pandocbounded{\includegraphics[keepaspectratio]{ch10_bond_yields.png}}
\caption{中国国债关键期限收益率曲线}
\end{figure}

接下来，看看债券的对立面------那个承担了组合中几乎全部''进攻''任务的角色：股票。

\newpage

\chapter{股票收益从哪里来}\label{ux80a1ux7968ux6536ux76caux4eceux54eaux91ccux6765}

\begin{center}{\color{dh-bronze}\rule{0.2\textwidth}{0.5pt}}\end{center}

\begin{dhquote}
\textbf{买股票不是买一串代码，而是买一家企业未来创造价值的能力。}
\end{dhquote}

\begin{center}{\color{dh-bronze}\rule{0.2\textwidth}{0.5pt}}\end{center}

对于很多第一次接触股票的人来说，股票市场给人的印象就是红红绿绿的数字、上下跳动的曲线，以及社交媒体上那些令人羡慕或令人恐惧的收益截图。这种印象不能说完全错误------股价确实每天都在波动。但它掩盖了一个更基本的事实：股票代表的是一家企业的部分所有权。长期来看，股票的价格最终会反映这家企业所创造的价值。短期来看，股价更多反映的是市场参与者的情绪。

理解股票收益的来源，是区分投资和投机的第一道分水岭。

\subsection{股票收益，来自三个地方}\label{ux80a1ux7968ux6536ux76caux6765ux81eaux4e09ux4e2aux5730ux65b9}

如果你持有一只股票，你的总回报（total return）来自三个部分。

第一部分是\textbf{股息（dividend）}。企业赚了钱之后，可以选择把一部分利润以现金的形式分配给股东。股息直接以现金进入你的账户，是最实在的回报形式。成熟型的大企业------比如银行、公用事业公司、消费品公司------往往有相对稳定的分红政策。成长型企业则通常将利润再投资于业务扩张，因此分红较少甚至完全不分红。

第二部分是\textbf{盈利增长（earnings
growth）}。企业的利润在增长，它作为一家公司的内在价值也在增长。如果市场给予的估值（市盈率等指标下详）保持不变，那么股价的涨幅就应该大致等于盈利的涨幅。这是股票长期收益的最根本来源------企业创造的真实价值。

第三部分是\textbf{估值变化（valuation
change）}。市场给同一家企业''出价''的高低，会受到利率、情绪、政策、宏观经济等多种因素的影响。有时一家企业盈利能力并未发生显著变化，但由于市场整体变得乐观或悲观，它的股价仍然可能大幅波动。估值变化是短期股价波动的主要驱动因素，但它在长期中的影响通常小于盈利增长。

以下数据是标普 500 和沪深 300 在过去约 20
年间（2005---2025）的年化总回报拆解：

\begin{longtable}[]{@{}lll@{}}
\toprule\noalign{}\rowcolor{dh-pink!25}
收益来源 & 标普500（S\&P 500） & 沪深300（CSI 300） \\
\midrule\noalign{}
\endhead
\bottomrule\noalign{}
\endlastfoot
盈利增长（earnings growth） & +6.8\%/年 & +8.5\%/年 \\
股息收益（dividend yield） & +2.0\%/年 & +2.2\%/年 \\
估值变化（valuation change） & +1.0\%/年 & −2.8\%/年 \\
\textbf{年化总回报（total return）} & \textbf{+9.8\%/年} &
\textbf{+7.9\%/年} \\
\end{longtable}

从拆解中可以读出两个关键信息。第一，在长期中，股票回报的最主要驱动力是盈利增长------标普
500 的 9.8\% 年化回报中，盈利增长贡献了约 70\%，沪深 300
的盈利增长贡献甚至超过
100\%（因为估值变化为负拖累了总回报）。第二，估值变化在长期中的贡献相对有限，而在中短期中却可能成为主导因素------尤其是在
A 股市场，2007---2008 年极端高估值到极端低估值的切换，直接吞噬了沪深 300
约 2.8\%
的年化回报。这再次印证了''短期看情绪、中期看估值、长期看盈利''的结论。

（数据来源：标普500 --- S\&P Dow Jones Indices 年度市场报告及 Bloomberg
终端共识数据；沪深300 ---
中证指数有限公司（www.csindex.com.cn）指数月报及 Wind
终端，盈利增长基于成分股汇总 EPS
增长率、股息收益基于指数股息率、估值变化基于 PE
变动。回报均为含分红再投资的总回报。）

一个简单的结论是：\textbf{短期看情绪，中期看估值，长期看盈利。}
如果你只持有一支股票几个星期或几个月，那么你赚或亏主要取决于别人愿意以什么价格接盘------这更接近于投机。如果你持有一支股票十年以上，它的大部分回报将来自企业在这十年间赚了多少钱、成长了多少。这就是为什么长期投资更接近分享企业创造的价值，而短期交易更接近于猜测市场参与者的行为。

\subsection{好公司不一定是好投资}\label{ux597dux516cux53f8ux4e0dux4e00ux5b9aux662fux597dux6295ux8d44}

这是一个初学者经常感到困惑的地方。一家企业产品质量优秀、品牌知名度高、管理层能力强------它是一家''好公司''。但你买入它的股票，是否就一定是一笔''好投资''？

不一定。因为股票的价格可能已经包含了市场对这家企业所有优点的预期。如果一家优秀企业的股价已经被推到了一个非常高的水平------反映的是市场假定它未来五年都能保持
20\% 的年增长率------那么即使它确实做到了 20\%
的年增长，股价也可能不涨反跌。原因是''20\%
的增长''已经被定价了，而市场原本期待的是 25\%。

反过来，一家看起来平平无奇的企业，如果经营稳定、现金流充裕、估值被市场长期低估，也可能在多年后为耐心的投资者带来可观的回报。

因此，投资股票不只是研究企业本身，还要研究你付出的价格。价格是你付出的，价值是你得到的。两者之间差距越大，你的安全边际（margin
of safety）就越宽。

\subsection{市值与股价：不要被表面数字迷惑}\label{ux5e02ux503cux4e0eux80a1ux4ef7ux4e0dux8981ux88abux8868ux9762ux6570ux5b57ux8ff7ux60d1}

很多人习惯用''一股多少钱''来判断股票贵不贵。茅台一股一两千元，某家小公司一股才几块钱------是不是小公司更便宜？这个类比是错误的。

一只股票贵不贵，不看它的每股价格，而看它的总市值相对于它创造利润的能力。一家总市值
10 亿元、每年赚 2 亿元的企业，其股价可能比一家总市值 100 亿元、每年只赚
1 亿元的企业更''便宜''。你需要看的是估值指标------比如市盈率（P/E
ratio），而不是股价本身。

\subsection{长期收益不等于短期价格的放大}\label{ux957fux671fux6536ux76caux4e0dux7b49ux4e8eux77edux671fux4ef7ux683cux7684ux653eux5927}

把时间轴拉长，股票确实是历史上表现最好的大众化投资工具之一。但这里的''长期''，指的是二十年、三十年甚至更长。如果你拿过去一百年的美国股票市场数据来看，年化回报大约在
6\% 到 7\%
左右（扣除通胀后）。这是经历了两次世界大战、多次经济衰退、数次股灾之后的结果。

问题是，在这''长期''的每一年之中，市场的表现可能完全不同。某一年涨
30\%，某一年跌
40\%，某一年波澜不惊。如果你计划用五年内需要的钱去买股票，你就可能在市场刚好低迷的那一年被迫以一个糟糕的价格卖出------这就是为什么我们在第四章强调了''投资期限''的重要性。

\begin{center}{\color{dh-bronze}\rule{0.2\textwidth}{0.5pt}}\end{center}

\begin{figure}
\centering
\pandocbounded{\includegraphics[keepaspectratio]{ch11_stock_returns.png}}
\caption{标普500与沪深300收益来源拆解}
\end{figure}

股票是投资组合中承担''增长引擎''角色的资产。没有股票，你的长期财富几乎不可能跑赢通胀。但持有股票需要两样东西------对企业价值的基本理解，以及对波动保持理性的能力。幸运的是，不是所有人都需要自己挑选个股。在本书的后半部分，我们将用大量篇幅讨论一种更简单但也同样有效的方式：通过指数基金和
ETF 来拥有整个市场，而不是赌某一家企业的命运。

在此之前，我们先把目光转向另外两类资产------黄金和
REITs------它们在不同的人生组合中扮演着各自独特的角色。

\newpage

\chapter{黄金、REITs
与其他资产}\label{ux9ec4ux91d1reits-ux4e0eux5176ux4ed6ux8d44ux4ea7}

\begin{center}{\color{dh-bronze}\rule{0.2\textwidth}{0.5pt}}\end{center}

\begin{dhquote}
\textbf{每一种资产都有它存在的理由，也有它不该出现的地方。}
\end{dhquote}

\begin{center}{\color{dh-bronze}\rule{0.2\textwidth}{0.5pt}}\end{center}

除了现金、债券和股票这三类最核心的资产之外，市场上还有一些重要的工具经常出现在家庭投资组合中。它们各有各的逻辑，但如果理解不清，也容易被过度使用或者被彻底忽视。这一章简要介绍黄金、REITs
和其他另类资产------不是为了让你立刻买入，而是让你在未来的资产配置中，能够清楚地判断：这个工具在我的组合里，到底应该承担什么角色。

\subsection{黄金：没有现金流，但有历史}\label{ux9ec4ux91d1ux6ca1ux6709ux73b0ux91d1ux6d41ux4f46ux6709ux5386ux53f2}

黄金是投资世界中最独特的一类资产。它不像股票一样有企业利润和股息，也不像债券一样有固定利息和到期还本。黄金本身不产生任何现金流。你买入一克黄金，放十年，它仍然是一克黄金------不会变成更多黄金。

那么黄金的价值到底来自哪里？

黄金的价值主要来自三点。第一，它的稀缺性------全球已探明的黄金储量有限，每年的开采量也相对稳定，不受某个国家货币政策的直接影响。第二，它的避险属性------在战争、金融危机、严重通胀或货币体系动荡时期，人们倾向于持有黄金，因为它不依赖于任何一个国家政府的信用。第三，它的历史惯性------几千年来，黄金在几乎所有文明中都被视为财富的象征和储存手段，这种文化惯性本身就是一种''共识价值''。

\textbf{黄金在重大危机期间的价格表现：}

\begin{longtable}[]{@{}
  >{\raggedright\arraybackslash}p{(\linewidth - 4\tabcolsep) * \real{0.3333}}
  >{\raggedright\arraybackslash}p{(\linewidth - 4\tabcolsep) * \real{0.3333}}
  >{\raggedright\arraybackslash}p{(\linewidth - 4\tabcolsep) * \real{0.3333}}@{}}
\toprule\noalign{}\rowcolor{dh-pink!25}
\begin{minipage}[b]{\linewidth}\raggedright
危机事件
\end{minipage} & \begin{minipage}[b]{\linewidth}\raggedright
期间
\end{minipage} & \begin{minipage}[b]{\linewidth}\raggedright
黄金表现（伦敦金，USD/oz）
\end{minipage} \\
\midrule\noalign{}
\endhead
\bottomrule\noalign{}
\endlastfoot
2008 年全球金融危机 & 2008 全年 & \textbf{+4.3\%}（同期标普500 跌
−38.5\%） \\
2011 年欧债危机 & 2011 全年 & \textbf{+9.5\%} \\
2020 年新冠疫情冲击 & 2020 全年 & \textbf{+24.6\%}（同期沪深300 全年仅
+27.2\%，但在 3 月单月一度暴跌） \\
2013 年（黄金大幅回撤年份） & 2013 全年 & \textbf{−28.3\%}（美联储缩减
QE 预期引发的''黄金闪崩''） \\
\end{longtable}

从长期来看，2005---2025 年的 21 年间，黄金的年化回报率约为
\textbf{+10.1\%}（按伦敦金定盘价 USD/oz 计算），略高于同期标普 500 的约
+9.8\%。但这一数字的''含金量''需要审慎看待：黄金约三分之一的累计回报集中在
2005---2011 年的大宗商品超级周期中。2012---2025 年间黄金的年化回报降至约
+5.9\%，且包含了 2013 年 −28.3\% 和 2021---2022
年连续下跌的年份。黄金具有显著的''阶段性爆发、长期温和''特征，不应简单用过去二十年的高
CAGR 线性外推。

（数据来源：LBMA（London Bullion Market
Association）伦敦金定盘价年度涨跌幅，https://www.lbma.org.uk；World Gold
Council 补充数据。）

黄金在投资组合中的角色，可以用两个词来概括：\textbf{防御和保险。}
它不是让你变富的工具，而是在其他资产表现不佳的时候提供缓冲。大多数长期投资者只把黄金作为总资产的一小部分（通常不超过
5\% 到 10\%），而且把它看作长期持有、不频繁交易的资产。

黄金也有局限性。因为不产生现金流，你无法用传统的估值方法来评估它''值多少钱''------它的价格主要由市场情绪和供需关系驱动，波动本身并不小。因此，把大部分资金投入黄金，不仅放弃了其他资产的长期增值潜力，也没有真正降低整体的波动性。

\subsection{REITs：投资房地产，但不用买房}\label{reitsux6295ux8d44ux623fux5730ux4ea7ux4f46ux4e0dux7528ux4e70ux623f}

在中国，住房资产在家庭财富中的占比非常高。很多人把大部分净资产都压在了自住房上，然后再想''有没有办法多配置一些房地产相关的资产''。

\textbf{REITs（Real Estate Investment Trusts，房地产投资信托基金）}
提供了一种不需要直接购房就能参与不动产投资的途径。它把大量商业不动产------写字楼、购物中心、物流仓库、数据中心、公寓楼等------打包成一个资金池，在证券市场上发行份额。

REITs 最吸引人的地方在于它的分红机制。各国法规通常要求 REITs
将绝大部分租金收入（在美国通常是应税收入的 90\%
以上）以分红的形式分配给投资者。因此，REITs
的投资者获得的回报主要由两部分构成：定期收到的分红，以及 REIT
本身价格的变动。

直接拥有一套出租房和持有一只 REITs 之间有几个重要区别。REITs
的门槛极低------几百元就能买一份，不需要凑首付；流动性远好于实物房产------在交易时间随时可以买卖；而且
REITs 投资的是多种物业的组合，不依赖某一套房、某一个租客。但 REITs
的价格也会波动------受利率、经济周期、商业地产供需等多种因素影响，波动幅度可能在一段时间内比宽基股票指数还大。因此，把它当作''稳定的收租工具''可能会低估它的风险和波动。

REITs
在投资组合中的职责，是在传统股票和债券之外增加一类与它们相关性较低的资产------让整个组合更多元。但它不适合作为新手唯一或最主要的投资标的，更不应该因为''觉得房价一直涨''而把大量资金集中到
REITs 中。

\subsection{商品、衍生品与另类资产：为什么不建议新手碰}\label{ux5546ux54c1ux884dux751fux54c1ux4e0eux53e6ux7c7bux8d44ux4ea7ux4e3aux4ec0ux4e48ux4e0dux5efaux8baeux65b0ux624bux78b0}

还有一些资产经常出现在投资讨论中，但不太适合刚开始建立长期组合的新手：大宗商品期货（原油、铜、大豆等）、加密货币、杠杆或反向
ETF、结构化理财产品、私募股权产品，等等。

这些工具之所以不适合新手，原因各不相同：有的涉及复杂的杠杆和保证金机制（亏损可能超过本金），有的底层资产缺乏现金流支撑（价格完全由市场情绪驱动），有的流动性较差且信息不透明（难以在需要时以合理价格卖出），有的结构复杂到销售方自己都未必完全理解。

一个诚实的建议是：在你能够用简单工具（股票指数基金、债券、现金）建立一个稳定运行的长期投资组合，并且真正经历了至少一轮完整的牛熊周期之前，不需要急于接触这些复杂的工具。投资不是收集名片------你不需要把什么都买一遍。真正有效的投资，往往是用少数几个你充分理解的工具，长期坚持，而不是用一大堆你不了解的工具频繁操作。

\begin{center}{\color{dh-bronze}\rule{0.2\textwidth}{0.5pt}}\end{center}

到这里，你已经把市场上最主要的资产类别过了一遍。从现金到债券，从股票到黄金和
REITs，每一种工具都有它存在的理由，也没有一种工具是完美的。

\begin{figure}
\centering
\pandocbounded{\includegraphics[keepaspectratio]{ch12_gold.png}}
\caption{2005-2025年黄金价格走势与危机表现}
\end{figure}

在下一部分，我们将进入全书最重要的主题之一：如何把这些工具组合起来，变成一套为你的目标服务的长期投资方案。

\partdivider{第四部分 · ETF 与基金的完整操作逻辑}

\newpage

\chapter{指数是如何编制出来的}\label{ux6307ux6570ux662fux5982ux4f55ux7f16ux5236ux51faux6765ux7684}

\begin{center}{\color{dh-bronze}\rule{0.2\textwidth}{0.5pt}}\end{center}

\begin{dhquote}
\textbf{指数不是一面''客观的镜子''，它是一套有偏好的规则。}
\end{dhquote}

\begin{center}{\color{dh-bronze}\rule{0.2\textwidth}{0.5pt}}\end{center}

在讨论 ETF 之前，我们必须先理解指数------因为大部分 ETF
都是在追踪某个指数。如果你不知道指数是怎么编出来的、里面装着什么、多久换一次成分，你就无法真正理解你买入的
ETF 到底持有了什么。

\subsection{指数是一套选股和加权规则}\label{ux6307ux6570ux662fux4e00ux5957ux9009ux80a1ux548cux52a0ux6743ux89c4ux5219}

想象你要跟踪中国最大的一些上市公司。你可以制定这样一套规则：在上海证券交易所和深圳证券交易所上市的所有公司中，选出总市值最大的
300
家，把它们按照市值大小分配权重------市值越大的公司，在指数中占的比例越高。每个月或每季度检查一次，如果某家公司市值跌出了前
300 名，就把它换掉，让新进入的公司补位。

这就是\textbf{沪深 300 指数（CSI 300 Index）}
的基本编制逻辑。它不是什么神秘的东西，就是一套公开的、机械执行的选股和加权规则。标普
500（S\&P 500）、纳斯达克 100（Nasdaq-100）、中证
500，都是基于类似逻辑的不同版本。

\subsection{三种主要的加权方式}\label{ux4e09ux79cdux4e3bux8981ux7684ux52a0ux6743ux65b9ux5f0f}

不同的指数使用不同的加权（weighting）方法，这会直接影响指数中哪些公司占据更大的比重。

\textbf{市值加权（market-cap weighting）}
是最常见的方法。公司的市值越大，在指数中的权重就越高。这是沪深 300、标普
500
等大多数主流指数采用的方式。它的优点是不需要频繁调仓------股价上涨自然导致权重上升，股价下跌自然导致权重下降。缺点是当某些大市值公司处于严重高估状态时，它们的权重也相应最高，投资者实际上在被动地''追涨''。

\textbf{等权重（equal weighting）}
给每一家成分公司分配相同的权重。如果指数有 100 家公司，每家公司占
1\%。这种方式避免了''大公司主导一切''的问题，但需要定期再平衡------因为股价波动会导致权重偏离，必须定期把涨多的卖掉、跌多的买入来恢复相等权重。等权指数的调仓成本通常高于市值加权指数。

\textbf{价格加权（price-weighted）}
按每一股的股价来决定权重------股价越高的公司，权重越大。这是道琼斯工业平均指数（Dow
Jones Industrial
Average）采用的方法。股价加权在今天被视为一种过时的方法，因为股价本身并不反映公司的真实规模------一家股价
1000 元但总市值只有 10
亿元的小公司，在价格加权指数中的权重可能超过一家股价 50 元但总市值 5000
亿元的大公司。

\subsection{指数不是天然客观的}\label{ux6307ux6570ux4e0dux662fux5929ux7136ux5ba2ux89c2ux7684}

很多人把指数想象成一架绝对中立的摄像机，忠实地记录市场发生了什么。事实上，每一只指数都是按照人为制定的规则运作的。不同指数提供商（如中证指数公司、MSCI、FTSE
Russell、S\&P Dow
Jones）对同一类市场可能编制出不同的指数，因为它们的规则------选股标准、权重上限、行业分类方式、调整频率------各不相同。

甚至同一只指数的成分股变化，也不完全反映市场的''自然优胜劣汰''。当有些公司由于临时性原因（比如被调查、被制裁、财务造假）被剔除出指数时，指数基金必须在限定的时间内按照规则被动卖出------无论价格是否合理。这种机械性操作有时也会带来额外的成本或风险。

理解这一点，不是为了让你抛弃指数投资，而是为了让你在投资 ETF
时能多问一个问题：\textbf{我买入的这只
ETF，背后追踪的是什么指数？这个指数是怎么选股、怎么加权的？其中有没有我需要注意的集中度或行业偏向？}

\subsection{两个具体案例：沪深 300 与标普 500
的成分调整}\label{ux4e24ux4e2aux5177ux4f53ux6848ux4f8bux6caaux6df1-300-ux4e0eux6807ux666e-500-ux7684ux6210ux5206ux8c03ux6574}

让我们通过两个具体的指数编制案例，更直观地理解''规则决定持仓''这件事。

\textbf{沪深 300：基于排名的机械化调整。} 沪深 300
由中证指数公司编制，每年 6 月和 12 月各进行一次定期审核（semi-annual
review），根据最近一年的日均总市值排名来筛选成分股。它的规则几乎完全是量化的：在沪深两市所有上市时间超过一个季度、非
ST 或 *ST、经营状况无重大异常的股票中，按日均总市值从高到低排序，取前
300
名。同时，为了保证指数的稳定性，排名设有一个''缓冲区''------如果一只股票已经是指数成分股，只有当它的市值排名跌出前
360 名时才会被剔除；而新进入指数需要排名进入前 240
名。这种缓冲区机制避免了股票在排名附近反复进出导致的频繁调仓。

由于沪深 300
的调整几乎完全依赖量化排名，它的调仓具有很强的可预测性。在每次审核之前，市场分析师通常可以提前几个月推算出哪些股票可能被纳入、哪些可能被踢出。这种透明度对
ETF
管理者来说是有利的------他们可以提前规划交易，减少临时调仓带来的市场冲击。

\textbf{标普 500：委员会说了算。} 与沪深 300 的机械排名不同，标普 500
的成分股出入由标普道琼斯指数公司的一个专门的指数委员会（Index
Committee）决定。委员会不公开具体的决策会议记录，但有一套公开的准入标准：公司必须是在美国注册、在纽约证券交易所或纳斯达克上市，最近四个季度的
GAAP 净利润之和为正且最近一个季度为正，流通市值至少达到约 146
亿美元（门槛会定期调整），且公众持股比例不低于 50\%。

即便一家公司满足了上述所有量化标准，它仍然可能被委员会拒绝纳入------或者反过来，一家公司不完全满足某些标准，但委员会认为它''代表了某个重要行业''，也可能被破格纳入。最著名的例子是特斯拉（Tesla）：尽管它在
2020 年初期已经满足了标普 500 的市值和盈利标准，委员会仍然等待了数月才在
2020 年 12
月正式将其纳入。委员会对盈利的可持续性有自身判断，这种''人的裁量权''既是指数质量的保障，也带来了不确定性。

\textbf{两个案例给 ETF 投资者的启示。} 如果你持有一只跟踪沪深 300 的
ETF，你可以清楚地知道每半年会有哪些股票进出------规则是透明的、可计算的。但如果你持有一只跟踪标普
500 的 ETF，成分变化的时间点和人选在委员会公布之前是无法确定的。这对 ETF
的运作有直接影响：沪深 300 ETF
的管理人可以更精准地规划调仓时间和交易策略；标普 500 ETF
的管理人则需要在委员会公告后迅速执行交易------往往恰好在其他跟踪同一指数的基金也同时执行交易的时候，可能带来短暂的价格波动。

\begin{center}{\color{dh-bronze}\rule{0.2\textwidth}{0.5pt}}\end{center}

指数是 ETF 的''原材料''。下一章我们将看到，ETF
这个金融工具虽然依附于指数运行，但它有自己的交易规则、自己的定价机制、自己的隐性成本。理解了指数的编制逻辑，你才能判断一只
ETF 是否忠实地反映了你想投资的标的------而不是被名称和营销宣传误导。

\newpage

\chapter{ETF
到底是怎么运作的}\label{etf-ux5230ux5e95ux662fux600eux4e48ux8fd0ux4f5cux7684}

\begin{center}{\color{dh-bronze}\rule{0.2\textwidth}{0.5pt}}\end{center}

\begin{dhquote}
\textbf{ETF
的价格不是由它追踪的指数决定的，而是由市场上买家和卖家的博弈决定的------但有一群专业玩家在确保它不会跑偏太远。}
\end{dhquote}

\begin{center}{\color{dh-bronze}\rule{0.2\textwidth}{0.5pt}}\end{center}

在前一章中我们讨论了指数是如何编制出来的。ETF（Exchange-Traded
Fund，交易型开放式指数基金）是一种在证券交易所上市交易的基金，它通常追踪某个指数的表现。但''追踪指数''这件事是怎么发生的？为什么有些
ETF 的价格和指数之间会有微小的偏离？为什么 ETF
的流动性比它的成交量更高？这是本章要回答的问题。

\subsection{ETF
与普通指数基金的区别}\label{etf-ux4e0eux666eux901aux6307ux6570ux57faux91d1ux7684ux533aux522b}

一只传统的指数基金（index mutual
fund），你只能在每个交易日收盘后，按照当天的基金净值（NAV，Net Asset
Value）进行申购或赎回。你的交易不会立刻完成------通常需要一到两个工作日。你买入的价格，是当天收盘后计算出来的净值，而不是你下单那一刻的价格。

一只 ETF
不同。它像股票一样，在交易时间内持续交易，价格由买卖双方的出价和要价决定，每秒钟都在变化。你可以在上午十点买，也可以在下午两点半卖。这就是''交易所交易''（exchange-traded）的含义。

交易方式的不同，带来了一系列运作机制上的差异。

\subsection{一级市场与二级市场：两套并行的系统}\label{ux4e00ux7ea7ux5e02ux573aux4e0eux4e8cux7ea7ux5e02ux573aux4e24ux5957ux5e76ux884cux7684ux7cfbux7edf}

理解 ETF，最关键的是区分它的两个市场。

\textbf{二级市场（secondary market）}，就是你和我这样的普通投资者买卖
ETF 的地方。��在券商 APP
里下单，输入数量、价格，点击买入------这就是在二级市场交易。你的交易对手是另一个普通投资者（或者做市商），而不是基金公司本身。二级市场上的
ETF
价格，完全由买卖双方的力量对比决定------买的人多，价格就高于净值；卖的人多，价格就低于净值。

\textbf{一级市场（primary market）}，是 ETF
份额被''创造''和''销毁''的地方。一级市场的参与者不是普通投资者，而是被称为\textbf{授权参与商（Authorized
Participant，简称 AP）} 的大型金融机构。他们与 ETF
发行方之间有特殊的协议。

当 AP 看到 ETF 的二级市场价格明显高于净值（比如溢价
1\%）时，他们可以做一件事：在一级市场按照净值从基金公司那里''申购''新的
ETF 份额------用一篮子成分股票交换 ETF
份额------然后在二级市场上以更高的价格卖出。这个过程增加了市场供给，把价格压回净值附近，同时
AP 赚到了差价。

反过来，当 ETF 价格明显低于净值（折价）时，AP 在二级市场上低价买入 ETF
份额，在一级市场上将其''赎回''为一篮子股票，再把股票卖出获利。这个过程减少了市场供给，把价格拉回净值附近。

这就是 ETF 定价最核心的机制：\textbf{AP 的套利行为，让 ETF
的二级市场价格始终在净值附近运行。}
不需要任何人为干预，纯粹的利润驱动就能完成这个''纠偏''过程。

\subsection{流动性从哪里来}\label{ux6d41ux52a8ux6027ux4eceux54eaux91ccux6765}

很多新手在选择 ETF 时会犯一个错误：看它的日均成交量。成交量小的
ETF，会不会''卖不出去''？

答案是：不一定。因为 ETF
的流动性，并不完全来自二级市场上的买卖成交量。它的流动性最终来自\textbf{底层资产的流动性}。

即使一只 ETF
在二级市场上一天只有几十万元成交，只要它持有的底层资产------比如沪深 300
的成分股------流动性非常好，AP 就可以随时在一级市场创造或赎回 ETF
份额来满足任何规模的买卖需求。因为 AP 不需要自己持有 ETF
库存------他们只需要在接到你的订单后，去市场上买入一篮子成分股，然后换成
ETF 份额交给你。

因此，评估 ETF 流动性的正确方式，不是看 ETF
本身的成交量，而是看\textbf{它持有的底层资产的流动性}。一只追踪沪深 300
的
ETF，即使自身成交量不大，流动性也几乎不是问题。而一只追踪小众行业或低流动性市场的
ETF，底层股票本身很难买卖，那么这只 ETF
的流动性就可能真的有问题------买卖价差较大，大额交易也可能推高或压低价格。

\subsection{IOPV：盘中参考净值}\label{iopvux76d8ux4e2dux53c2ux8003ux51c0ux503c}

一只传统基金的净值（NAV）每天只公布一次------在收盘后，根据当天持仓的收盘价计算。但在交易时间，ETF
的价格在不断变化，你怎么知道它是贵了还是便宜了？

交易所通常会提供\textbf{IOPV（Indicative Optimized Portfolio
Value，参考净值）}，也叫做盘中净值，每 15 秒更新一次。它是根据 ETF
持仓的实时价格估算出来的参考值。你把 ETF 的当前市价和这个 IOPV
对比，就能大致判断它是在溢价交易还是在折价交易。

IOPV
是一个参考值，不是精确值------它可能因为底层资产的价格延迟、外汇折算差异等各种原因而存在误差。但它已经足够让普通投资者在交易时多一个判断依据。

\begin{center}{\color{dh-bronze}\rule{0.2\textwidth}{0.5pt}}\end{center}

理解 ETF
的运作机制，并不会让你多赚多少钱。但它能帮你避免一些不必要的损失------比如在没有搞清楚溢价率的情况下买入了高溢价的
QDII ETF。在下一章，我们专门讨论这个问题：ETF
为什么会溢价或折价？高溢价买入到底有多危险？

\newpage

\chapter{ETF
的溢价率与折价率}\label{etf-ux7684ux6ea2ux4ef7ux7387ux4e0eux6298ux4ef7ux7387}

\begin{center}{\color{dh-bronze}\rule{0.2\textwidth}{0.5pt}}\end{center}

\begin{dhquote}
\textbf{高溢价买入 ETF，可能是在花 120 元买一个只值 100
元的篮子------即使篮子里的东西没变，你已经在入口处亏了钱。}
\end{dhquote}

\begin{center}{\color{dh-bronze}\rule{0.2\textwidth}{0.5pt}}\end{center}

在第十四章中我们提到，ETF
的二级市场价格会围绕净值波动，而授权参与商的套利机制会防止偏离过大。但这个机制并不是完美无缺的------在某些情况下，ETF
的价格可能会显著高于或低于净值。对于投资者来说，理解这种现象并学会计算溢价率，是避免不必要亏损的基本功。

\subsection{溢价和折价的计算}\label{ux6ea2ux4ef7ux548cux6298ux4ef7ux7684ux8ba1ux7b97}

溢价率和折价率的计算公式非常简单：

\textbf{溢价率 = (市价 - 净值) ÷ 净值 × 100\%}

如果一只 ETF 的净值是 1.000 元，而市场上的成交价是 1.200 元，溢价率就是
(1.200 - 1.000) ÷ 1.000 × 100\% = 20\%。这意味着你正在花 1.20
元买一个只值 1.00 元的资产篮子。

反过来，如果市价是 0.900 元，净值是 1.000 元，那就是折价
10\%。你在以九折的价格买入资产篮子------听起来像好事，但需要判断折价会不会继续扩大，以及你是否能够承受进一步的下跌。

\subsection{溢价买入为什么危险}\label{ux6ea2ux4ef7ux4e70ux5165ux4e3aux4ec0ux4e48ux5371ux9669}

高溢价买入的危险并不复杂。你多付了钱，而那个多付的部分并不是因为资产本身变值钱了。

假设你以 1.200 元买入了一只溢价率为 20\% 的
ETF。过了一段时间，溢价率回归到零------市价回到了净值 1.000
元。你亏损了多少？

很多人的直觉回答是 20\%。但让我们来算一下：(1.200 - 1.000) ÷ 1.200 =
0.200 ÷ 1.200 ≈ 16.7\%。你实际的亏损比例是 16.7\%，而不是
20\%。原因是溢价率（以净值为分母）和你的实际亏损比例（以你的买入成本为分母）计算基数不同。

这个细微的差别并不改变一个核心事实：\textbf{你在一笔投资还没开始的时候，就已经因为价格偏差而处于亏损状态。}
即便 ETF
持仓的底层资产一分不涨、一分不跌，只要溢价率收窄或消失，高溢价买入的投资者就会亏损。如果底层资产同时下跌，亏损就是双重的------资产跌了，溢价也没了。

\subsection{为什么会产生溢价和折价}\label{ux4e3aux4ec0ux4e48ux4f1aux4ea7ux751fux6ea2ux4ef7ux548cux6298ux4ef7}

AP
的套利机制本来应该把溢价率压制在一个很小的范围内。那么，什么情况下它会失效？

\textbf{跨市场时差} 是最常见的原因之一。一只追踪美国股市的 QDII ETF
在中国交易所上市交易。中国股市开盘的时候，美国市场已经收盘了。ETF
今天的净值是多少？不知道------因为底层资产的美股昨天收盘价是旧的，今天还没开盘。这时候，投资者只能根据对今晚美股走势的预期来出价。如果大家普遍预期美股今晚会大涨，ETF
的市价就可能被推到高于''旧净值''的位置，形成溢价。如果预期今晚大跌，则可能出现折价。

\textbf{停牌资产} 也会导致定价失准。如果 ETF
持仓中的某只成分股因为重大事项停牌了，那么 ETF
的净值计算就不得不继续使用停牌前的价格------这个价格可能已经严重过时。投资者对这只成分股真实价值的猜测，会反映在
ETF 的市价中，从而产生与''官方净值''之间的偏离。

\textbf{外汇额度限制} 是 QDII 类 ETF 特有的问题。QDII
基金有外汇额度上限。当大量资金涌入一只热门的 QDII
ETF，而基金管理人已经用完了外汇额度时，就无法在一级市场创造新的 ETF
份额。供给被锁住了，需求继续涌入二级市场，溢价就会越推越高。在一些极端案例中，某些
QDII ETF 的溢价率曾达到 20\% 甚至更高。

\textbf{流动性不足} 也会导致溢价或折价。对于成交量极其清淡的
ETF，做市商给出的买卖报价可能本身就包含了较大的差价------买价偏低、卖价偏高，以补偿自己承担的风险。这也会体现为某种程度上的''结构性溢价''。

\textbf{市场情绪}
是最不可控的因素。牛市中，追涨情绪高涨，投资者愿意支付溢价买入热门
ETF；熊市中，恐慌性抛售可能导致 ETF
跌破净值。情绪驱动的溢价和折价通常来得快、去得也快。

\subsection{高溢价买入后的几种可能走向}\label{ux9ad8ux6ea2ux4ef7ux4e70ux5165ux540eux7684ux51e0ux79cdux53efux80fdux8d70ux5411}

当你买入了一只高溢价的 ETF 之后，未来可能发生以下几种情况：

第一，溢价继续扩大。你买的时候溢价 15\%，后来涨到
25\%。如果你选择卖出，你在溢价层面是赚钱的。但这本质上是在赌博------赌别人愿意以更高的溢价接盘。

第二，溢价逐渐收窄。这是最常见的情况。随着套利机制发挥作用和市场情绪冷却，溢价慢慢回到正常水平（通常在
1\% 以内）。你账面上出现了亏损------即使底层资产本身没有变化。

第三，基金公司增加供给。如果基金管理人获得了新的外汇额度或者找到了其他方式增加
ETF
份额，一级市场供给恢复，溢价可能迅速消失。那些在高溢价时买入的投资者，会发现自己持有的资产在一天之内大幅''缩水''。

第四，底层资产大幅上涨。如果溢价消失带来的损失被底层资产本身的上涨所覆盖------比如你买入时溢价
10\%，但资产本身在后续涨了
20\%------那么你的总账仍然是赚钱的。但这并不意味着高溢价买入是正确的决策。只是运气眷顾了你。

\subsection{如何避免高溢价陷阱}\label{ux5982ux4f55ux907fux514dux9ad8ux6ea2ux4ef7ux9677ux9631}

在买入任何 ETF 之前，花一分钟查看它的溢价率。大多数券商 APP
和基金平台都会实时显示 ETF 的 IOPV（参考净值）和溢价率。如果溢价率超过
3\%，尤其是 QDII 类
ETF，就需要认真考虑------你是否愿意为这笔交易额外支付至少 3\%
的''入场费''。如果溢价率超过
5\%，对于长期投资者来说，通常更明智的选择是等待溢价回落再买入，或者寻找追踪同类指数的替代产品（如另一只规模更大、供给更充足的
ETF）。

当然，如果溢价率是负数------也就是折价------而且折价的幅度大于正常水平，有时反而是买入的机会。但前��是你要搞清楚为什么会折价------是因为流动性问题、底层资产风险、还是纯粹的市场情绪。

\begin{center}{\color{dh-bronze}\rule{0.2\textwidth}{0.5pt}}\end{center}

溢价率是 ETF
操作中最容易被新手忽略、也最容易造成''不知道为什么就亏了钱''的问题。它并不复杂，只是很少被重视。

\begin{figure}
\centering
\pandocbounded{\includegraphics[keepaspectratio]{ch15_qdii_premium.png}}
\caption{QDII ETF溢价率变化示意}
\end{figure}

在下一章，我们将讨论另一个同样容易被忽视的话题：那些看似微小的费率，如何在长期中吞噬你的收益。

\newpage

\chapter{ETF
的费率、跟踪误差与隐性成本}\label{etf-ux7684ux8d39ux7387ux8ddfux8e2aux8befux5deeux4e0eux9690ux6027ux6210ux672c}

\begin{center}{\color{dh-bronze}\rule{0.2\textwidth}{0.5pt}}\end{center}

\begin{dhquote}
\textbf{你看到的费率，只是一个百分比；你看不到的那个百分比，乘以三十年之后，可能是几十万元。}
\end{dhquote}

\begin{center}{\color{dh-bronze}\rule{0.2\textwidth}{0.5pt}}\end{center}

在挑选 ETF
的时候，很多人首先关注的是它的历史收益------过去一年涨了多少、过去三年涨了多少。但收益是过去的事，无法代表未来。而有一个东西是确定会发生的，而且会年复一年地从你的账户中抽取：费用。

\subsection{看得见的费用：管理费与托管费}\label{ux770bux5f97ux89c1ux7684ux8d39ux7528ux7ba1ux7406ux8d39ux4e0eux6258ux7ba1ux8d39}

每一只 ETF 都有两项最基本的费用。

\textbf{管理费（management fee）}
是基金管理公司收取的费用，用于覆盖基金的日常运作------基金经理的薪酬、研究支持、合规成本、信息披露等。管理费通常按年化比率计算，每天从基金净资产中扣除一小部分。比如一只管理费为
0.15\% 的 ETF，每年会从你的持有份额中扣除相当于资产规模 0.15\% 的费用。

\textbf{托管费（custodian fee）}
是基金资产托管银行收取的费用，用于保管基金持有的底层证券------股票、债券、现金等。托管费同样按年化比率收取，通常比管理费更低。

你不需要手动支付这些费用------它们会在基金的净值中自动扣除，你看到的所有净值数据已经是费用扣除后的结果。这种''无声扣除''的好处是不用操心操作，坏处是很多人根本不知道自己付了多少钱。

\subsection{费率的长期复利效应}\label{ux8d39ux7387ux7684ux957fux671fux590dux5229ux6548ux5e94}

假设你投资 10 万元，长期年化收益率是
7\%（教学假设），持有三十年。如果选择的 ETF 综合费率为
0.20\%，三十年后你实际得到的年化收益率约为 6.8\%，最终金额约为 72 万元。

如果选择另一只跟踪同一指数、但综合费率为 0.60\% 的
ETF，你的实际年化收益率降到约 6.4\%，最终金额约为 64 万元。同样是 10
万元本金，0.40\% 的费率差异，三十年后相差了约 8 万元。

下表以 10 万元本金、7\% 年化毛收益（教学假设）为前提，展示四种费率水平在
10 年、20 年、30 年后的终值差异：

\begin{longtable}[]{@{}
  >{\raggedright\arraybackslash}p{(\linewidth - 8\tabcolsep) * \real{0.2000}}
  >{\raggedright\arraybackslash}p{(\linewidth - 8\tabcolsep) * \real{0.2000}}
  >{\raggedright\arraybackslash}p{(\linewidth - 8\tabcolsep) * \real{0.2000}}
  >{\raggedright\arraybackslash}p{(\linewidth - 8\tabcolsep) * \real{0.2000}}
  >{\raggedright\arraybackslash}p{(\linewidth - 8\tabcolsep) * \real{0.2000}}@{}}
\toprule\noalign{}\rowcolor{dh-pink!25}
\begin{minipage}[b]{\linewidth}\raggedright
综合年费率
\end{minipage} & \begin{minipage}[b]{\linewidth}\raggedright
10 年终值
\end{minipage} & \begin{minipage}[b]{\linewidth}\raggedright
20 年终值
\end{minipage} & \begin{minipage}[b]{\linewidth}\raggedright
30 年终值
\end{minipage} & \begin{minipage}[b]{\linewidth}\raggedright
与零费率相比 30 年被费率吞噬的金额
\end{minipage} \\
\midrule\noalign{}
\endhead
\bottomrule\noalign{}
\endlastfoot
0.15\%（低成本 ETF） & ¥193,975 & ¥376,262 & ¥729,854 & ¥31,372 \\
0.30\%（中低费率 ETF） & ¥191,269 & ¥365,838 & ¥699,733 & ¥61,492 \\
0.50\%（中等费率 ETF） & ¥187,714 & ¥352,365 & ¥661,437 & ¥99,789 \\
1.00\%（主动基金常见费率） & ¥179,085 & ¥320,714 & ¥574,349 &
¥186,876 \\
0.00\%（零费率理论基准） & ¥196,715 & ¥386,968 & ¥761,226 & --- \\
\end{longtable}

从上表可以看到，0.15\% 和 1.00\% 两个费率水平之间，30 年终值相差约
¥155,000------相当于本金的 1.55
倍被费率吞噬。这不是基金经理''赚不到钱''，而是数学：哪怕年化毛收益完全一致，费率差异在复利作用下的终值差距会随着时间急剧放大。因此，在长期投资中，费率是少数几个你''确定能控制''且''确定会影响收益''的变量。

费率差异对长期收益的影响，在初始阶段几乎看不见。第一年，0.40\%
的费率差异只差了几十块钱。但三十年之后，这个差异被复利放大了几十倍。这就是为什么在长期投资中，费率是少数几个''确定会对收益产生负面影响''的因素之一，也是为什么我们反复强调''低成本''原则。

\subsection{看不见的成本：买卖价差与交易佣金}\label{ux770bux4e0dux89c1ux7684ux6210ux672cux4e70ux5356ux4ef7ux5deeux4e0eux4ea4ux6613ux4f63ux91d1}

在费率之外，ETF 的交易还有一些隐形成本。

\textbf{买卖价差（bid-ask spread）}
是你买入和卖出之间存在的价格差异。做市商给出的买价（你卖出时的成交价）总是略低于卖价（你买入时的成交价），这个差额就是做市商的收入来源之一。对于流动性充足的
ETF，买卖价差可能只有万分之一甚至更低；对于交易清淡的
ETF，买卖价差可能达到千分之几甚至更高。

每次买卖，你都在''不知不觉中''支付这个差价。如果你每年只操作一两次，这个成本可以忽略。如果你频繁交易，买卖价差的累积效应可能超过管理费本身。

\textbf{交易佣金（commission）}
在大部分主流券商平台上已经降到了很低的水平，有些甚至对 ETF
免佣金。但并非所有平台都如此，每家券商的费率结构不同，需要自行确认。

\subsection{跟踪误差与跟踪差异}\label{ux8ddfux8e2aux8befux5deeux4e0eux8ddfux8e2aux5deeux5f02}

选择了''跟踪沪深 300 指数''的 ETF，并不意味着这只 ETF 的表现会和沪深 300
指数一模一样。两者之间的偏差，叫做跟踪误差或跟踪差异。

\textbf{跟踪误差（tracking error）} 通常用标准差来衡量------它告诉你 ETF
收益与指数收益之间偏离的''波动程度''有多大。\textbf{跟踪差异（tracking
difference）} 则直接给出 ETF
收益与指数收益之间的实际差额------比如指数涨了 10\%，ETF 涨了
9.7\%，跟踪差异就是 -0.3\%。

跟踪差异的主要来源包括：管理费和托管费（这是最大的来源）、现金拖累（ETF
持有一小部分现金以备赎回）、指数调仓时 ETF
调仓的时间差、证券借贷收入（可能弥补部分费用）、以及复制方式的限制。

对普通投资者来说，最重要的指标不是跟踪误差，而是跟踪差异------它直接告诉你，你相对于''完美复制指数''少赚了多少钱。这个数字加上费率，就是一只
ETF 的真实运行成本。

\subsection{其他隐性成本}\label{ux5176ux4ed6ux9690ux6027ux6210ux672c}

\textbf{税费（tax impact）} 在不同地区差异巨大。在中国市场，ETF
交易目前暂不征收印花税，但分红收入可能需要缴纳个人所得税，具体取决于 ETF
的类型和投资者的税务身份。跨境 ETF
还可能涉及外国预扣税------例如投资美股的
ETF，其持有的美国股票分红通常会被美国国税局扣缴一部分税款。

\textbf{换汇成本（currency conversion cost）} 在 QDII ETF
中尤为显著。基金将人民币兑换为外币（如美元、港元）购买底层资产，再将外币收益换回人民币。每次换汇都涉及汇率差价和手续费，这些成本最终会反映在基金的净值中。

\textbf{证券借贷收入（securities lending income）}
可以部分抵消费用。基金可以将持有的股票暂时借给做空者，收取借贷费用。一些大型
ETF
通过证券借贷获得的收入可以覆盖相当一部分管理费------但这取决于市场环境和基金的具体做法，并非稳定预期。

\begin{center}{\color{dh-bronze}\rule{0.2\textwidth}{0.5pt}}\end{center}

费率和成本不是选 ETF
的唯一标准------一只费率极低但流动性极差、或者追踪指数不合你目标的
ETF，并不适合你。但在同等条件或近似条件下，费率越低，长期收益越接近于你所投资市场的真实回报。

\begin{figure}
\centering
\pandocbounded{\includegraphics[keepaspectratio]{ch16_etf_fees.png}}
\caption{不同费率水平对30年复合收益的影响}
\end{figure}

在下一章，我们将整理一份完整的 ETF
筛选清单------帮助你从投资目标、指数、费率、规模、流动性到税务，逐项审视一只
ETF 是否值得长期持有。

\newpage

\chapter{如何真正筛选一只
ETF}\label{ux5982ux4f55ux771fux6b63ux7b5bux9009ux4e00ux53ea-etf}

\begin{center}{\color{dh-bronze}\rule{0.2\textwidth}{0.5pt}}\end{center}

\begin{dhquote}
\textbf{你不需要把市场上所有 ETF 都研究一遍。你只需要对你想买的每一只
ETF，回答这二十个问题。}
\end{dhquote}

\begin{center}{\color{dh-bronze}\rule{0.2\textwidth}{0.5pt}}\end{center}

在了解了指数编制、ETF
运作机制、溢价折价和费率之后，现在我们可以把所有这些知识整合成一份实用的筛选工具。以下二十个检查项目，不是要你每一次买入都逐条填写------而是在你第一次研究一只
ETF 的时候，帮助你避免遗漏重要信息。

\begin{center}{\color{dh-bronze}\rule{0.2\textwidth}{0.5pt}}\end{center}

\subsection{ETF
筛选二十项检查清单}\label{etf-ux7b5bux9009ux4e8cux5341ux9879ux68c0ux67e5ux6e05ux5355}

\textbf{关于投资目标与指数}

\textbf{1. 投资目标匹配度：} 这只 ETF
在我的投资组合中应当承担什么角色？是作为核心配置（宽基指数），还是卫星配置（行业/主题），还是防御配置（债券/黄金）？如果回答不清楚，就不要买。

\textbf{2. 跟踪指数：} 这只 ETF
跟踪的是哪一个具体指数？指数的全称、代码、编制机构是什么？不要只看 ETF
的名字------同一类名称可能对应不同的指数。

\textbf{3. 指数编制规则：} 指数是如何选股的（市值排名？流动性筛选？ESG
标准？）、如何加权的（市值加权？等权重？因子加权？）、多久调整一次成分股？这些规则直接影响你的持仓结构。

\textbf{4. 成分股数量：} 指数包含多少家公司？成分股太少（例如少于 50
家），集中度可能过高；成分股数量很大（上千家），分散度较好但尾部公司影响微乎其微。

\textbf{5. 前十大持仓集中度：} 前十大持仓合计占比多少？如果超过
50\%，说明这只 ETF
的权重高度集中在少数几家公司------尽管它名义上持有很多股票。

\textbf{6. 行业集中度：}
第一大行业占比多少？前三大行业合计占比多少？如果某只''宽基指数''的第一大行业占比超过
30\%，它实际上带有显著的行业倾斜。

\textbf{7. 地区集中度：} 如果是一只覆盖多个国家或地区的
ETF，主要市场占比多少？美国占比多少、中国占比多少？单一市场权重过高的''全球''ETF，实际上承担了较大的地域集中风险。

\textbf{关于成本与效率}

\textbf{8. 综合费率（总费率，Total Expense Ratio）：} 管理费 + 托管费 +
其他费用合计多少？这个数字应该能从基金公告或招募说明书中找到。

\textbf{9. 跟踪差异（tracking difference）：} 过去三年或五年，ETF
的实际年化回报与指数年化回报之间差距是多少？这个数字已经包含了费率、交易成本和证券借贷收入的全部影响------是''真实成本''的最佳衡量指标。

\textbf{10. 买卖价差（bid-ask spread）：} 这只 ETF
在日常交易中的买价和卖价之间差了多少？流动性好的 ETF
通常只有极小的价差。

\textbf{11. 溢价折价历史：} 这只 ETF
在过去是否存在持续性或大幅度的溢价或折价？如果是 QDII
ETF，历史上最高溢价率是多少？

\textbf{12. 复制方式：} 它是通过实际持有全部成分股（完全复制，full
replication），还是通过持有部分成分股加衍生品（抽样复制，sampling /
synthetic replication）来追踪指数？对于新手投资者来说，完全复制的 ETF
通常更容易理解其持仓和风险。

\textbf{13. 分红政策：} 这只 ETF
是累积型（accumulating，分红自动再投资）还是分配型（distributing，分红以现金形式发放）？你的投资目标更适合哪一种？

\textbf{关于规模与流动性}

\textbf{14. 基金规模：} 这只 ETF 的总资产规模（AUM）多大？规模过小的 ETF
存在被清盘的风险------如果基金公司认为运营成本高于收入，可能决定终止该产品。

\textbf{15. 日均成交量：} 虽然 ETF
的流动性不完全取决于二级市场成交量（见第十四章），但成交量仍然反映了交易活跃度。极度清淡的
ETF 可能导致买卖价差较大、大额交易时产生滑点。

\textbf{16. 底层资产流动性：} ETF
持有的股票或债券本身流动性好吗？如果投资的是一些小型、新兴或受限市场的资产，底层流动性可能成为真正的瓶颈。

\textbf{关于风险与合规}

\textbf{17. 税率影响：} 这只 ETF 的分红收益需要缴纳哪些税？如果 ETF
注册在海外，是否存在预扣税（withholding
tax）？不同地区的税务处理差异需要通过自己的尽职调查（due
diligence）来确认。

\textbf{18. 汇率风险：} 如果 ETF
以非人民币计价或投资于海外资产，汇率波动会对实际收益产生多大影响？你是否愿意承担这个额外的波动来源？

\textbf{19. 杠杆或反向结构：} 这只 ETF
是否含有杠杆（2x、3x）或反向（-1x）结构？如果是，你的投资目标是否真正需要这种工具？杠杆
ETF
由于每日再平衡机制的影响，长期持有的收益可能与指数乘以杠杆倍数完全不同，不适合作为长期投资工具。

\textbf{20. 产品终止风险：} 这只 ETF 的管理人是否有过清盘或合并 ETF
的历史？规模较小、成立不久、或者发行方规模不大的
ETF，面临更大的终止风险。产品终止时你通常可以收回净值对应的资金，但可能被迫在不合适的时间点退出投资。

\begin{center}{\color{dh-bronze}\rule{0.2\textwidth}{0.5pt}}\end{center}

\subsection{示例：用清单看一只
ETF}\label{ux793aux4f8bux7528ux6e05ux5355ux770bux4e00ux53ea-etf}

以下是将上述清单应用于一只实际 ETF 的简化示例------\textbf{华泰柏瑞沪深
300 ETF（代码
510300）}。这只是一个教学演示，帮助理解如何使用这份清单；不构成买入建议。

\begin{longtable}[]{@{}
  >{\raggedright\arraybackslash}p{(\linewidth - 4\tabcolsep) * \real{0.3333}}
  >{\raggedright\arraybackslash}p{(\linewidth - 4\tabcolsep) * \real{0.3333}}
  >{\raggedright\arraybackslash}p{(\linewidth - 4\tabcolsep) * \real{0.3333}}@{}}
\toprule\noalign{}\rowcolor{dh-pink!25}
\begin{minipage}[b]{\linewidth}\raggedright
检查项
\end{minipage} & \begin{minipage}[b]{\linewidth}\raggedright
510300 的情况
\end{minipage} & \begin{minipage}[b]{\linewidth}\raggedright
备注
\end{minipage} \\
\midrule\noalign{}
\endhead
\bottomrule\noalign{}
\endlastfoot
1. 投资目标匹配 & 核心配置（宽基指数） & 沪深 300 覆盖 A
股大盘，适合作权益类核心持仓 \\
2. 跟踪指数 & 沪深 300 指数（CSI 300） & 编制机构：中证指数公司 \\
3. 编制规则 & 市值排名前 300，市值加权，每半年调整 &
缓冲区规则：成分股跌出前 360 才剔除 \\
4. 成分股数量 & 300 只 & 覆盖沪深两市大市值公司 \\
5. 前十大集中度 & 约 20--25\%（随行情波动） &
集中度适中，不过度依赖少数公司 \\
6. 行业集中度 & 金融 + 消费合计约 40\% & 沪深 300
历史上金融权重偏高，近年有所分散 \\
8. 综合费率 & 管理费 0.15\% + 托管费 0.05\% = \textbf{0.20\%} &
在国内宽基 ETF 中属于低费率档位 \\
14. 基金规模 & \textasciitilde1000 亿元（市场最大之一） &
流动性充足，无清盘风险 \\
12. 复制方式 & 完全复制（full replication） & 实际持有全部 300
只成分股 \\
19. 杠杆/反向 & 无 & 标准 ETF，不含杠杆结构 \\
\end{longtable}

上面只填写了最核心的十个检查项（整个清单有二十项）。你不需要每次都填满------对新
ETF
第一次研究时覆盖全部项目，后续跟踪检查时只关注几个关键变化（费率、规模、集中度、跟踪差异）就足够了。

\begin{center}{\color{dh-bronze}\rule{0.2\textwidth}{0.5pt}}\end{center}

这份清单不是要求每一只 ETF
都完美通过所有项目。它的目的是确保你在买入之前，已经对自己买的是什么有了清晰的认识。投资中最贵的不是费率，而是因为不了解而犯下的错误。

在下一章，我们将离开被动投资的领域，看看主动管理基金是否值得配置------以及如何判断一只基金的业绩是来自能力还是运气。

\newpage

\chapter{主动基金应该怎么看}\label{ux4e3bux52a8ux57faux91d1ux5e94ux8be5ux600eux4e48ux770b}

\begin{center}{\color{dh-bronze}\rule{0.2\textwidth}{0.5pt}}\end{center}

\begin{dhquote}
\textbf{买主动基金，不是买一只基金------是选择一个人、一个团队、一种投资哲学。}
\end{dhquote}

\begin{center}{\color{dh-bronze}\rule{0.2\textwidth}{0.5pt}}\end{center}

在前面几个章节中，我们花了大量篇幅讨论指数
ETF。但在很多人的投资组合中，主动管理基金仍然占有重要的一席之地。主动基金与被动基金之间并不是''一个绝对好、一个绝对坏''。它们只是在用不同的方式解决同一个问题：如何让投资者的资金增值。

\subsection{主动与被动：根本的分歧}\label{ux4e3bux52a8ux4e0eux88abux52a8ux6839ux672cux7684ux5206ux6b67}

主动基金的基金经理试图通过自己的研究和判断，选出能够战胜市场的股票或债券------在好的价格买入，在更好的价格卖出，实现超越市场平均水平的回报。被动基金（如指数
ETF）的目标则简单得多：尽可能准确地复制某个市场指数的表现，不主动做任何判断，只求不输给市场。

主动基金要跑赢被动基金，必须克服两个障碍：一是费用------主动基金的管理费通常远高于被动基金；二是市场效率------在一个由成千上万专业投资者参与的市场中，持续击败对手是一件极其困难的事情。

SPIVA（S\&P Indices Versus
Active）是全球最权威的主动基金与被动指数基金对比研究。根据其最新年度报告（截至
2024
年年终数据），主动管理型基金在多个市场和不同时间维度下跑输基准的比例令人警醒：

\begin{longtable}[]{@{}lll@{}}
\toprule\noalign{}\rowcolor{dh-pink!25}
市场 / 类别 & 统计周期 & 跑输基准的主动基金占比 \\
\midrule\noalign{}
\endhead
\bottomrule\noalign{}
\endlastfoot
美股大盘基金（US Large-Cap） & 1 年 & 51\% \\
美股大盘基金（US Large-Cap） & 5 年 & 75\% \\
美股大盘基金（US Large-Cap） & \textbf{10 年} & \textbf{87\%} \\
美股中盘基金（US Mid-Cap） & 5 年 & 58\% \\
美股小盘基金（US Small-Cap） & 5 年 & 62\% \\
新兴市场基金（Emerging Markets） & 5 年 & 84\% \\
国际股票基金（International） & 5 年 & 65\% \\
中国 A 股主动基金 & 1 年 & 42\% \\
中国 A 股主动基金 & 3 年 & 55\% \\
中国 A 股主动基金 & \textbf{5 年} & \textbf{68\%} \\
\end{longtable}

数据揭示了一个清晰且一致的规律：时间越长，主动基金战胜指数的概率越低。在美股市场，10
年期视角下近九成的主动基金跑输基准；即使在 A
股（常被认为''散户多、主动基金容易产生超额收益''），五年期也有近七成的主动基金未能战胜指数。这不是对基金经理个人能力的否定，而是对一个结构性事实的印证：在扣除管理费、交易成本和税收拖累后，持续战胜市场是一个概率极低的游戏。

（数据来源：S\&P SPIVA
Scorecard（最新年度报告），https://www.spglobal.com/spdji/en/research-insights/spiva/；中国
A 股数据对应报告中的 China Equity 分类。）

\subsection{评估一只主动基金的五个维度}\label{ux8bc4ux4f30ux4e00ux53eaux4e3bux52a8ux57faux91d1ux7684ux4e94ux4e2aux7ef4ux5ea6}

如果你决定配置主动基金，以下几个评估维度比''近一年收益排名''重要得多。

\textbf{基金经理。}
一只主动基金的灵魂是管理它的那个人（或团队）。现任基金经理管理这只基金多久了？他之前的投资记录如何？他的投资理念是否清晰、前后一致？如果一只基金频繁更换基金经理，它的历史业绩就不具备任何参考价值------因为过去的成绩是前任留下的，和现在毫无关系。

\textbf{超额收益（Alpha）。} Alpha
衡量的是基金相对于其基准指数所创造（或损失）的额外回报。一只年化回报
10\% 的基金听起来不错，但如果它对应的基准指数同期年化回报是 12\%，它的
Alpha 就是 -2\%------实际上跑输了市场。反过来，一只年化回报 8\%
的基金，如果基准是 5\%，即使绝对回报不高，也说明基金经理创造了正 Alpha。

\textbf{风格漂移（style drift）。}
一只基金宣称自己是''大盘价值型基金''，但某段时间突然重仓了成长股或小盘股------这就是风格漂移。风格漂移意味着你买的和你以为自己在买的东西不一样。对于投资者来说，风格漂移会让资产配置变得不可控，因为你文件夹里本来打算作为''稳健配置''的基金，可能在不知不觉中变成了''激进配置''。

\textbf{最大回撤（maximum drawdown）。}
这只基金在历史上最大的''从高峰到低谷''的跌幅是多少？回撤幅度是否在你的风险承受范围内？一个年化回报很高的基金，如果历史上最大回撤达到了
50\%，它可能不适合大多数非专业投资者------因为在那个回撤的过程中，绝大多数人会在最低点附近恐慌卖出。

\textbf{换手率（turnover ratio）。}
基金在一定时期内买卖股票的频率。换手率过高意味着频繁交易，产生更高的交易成本（包括佣金、印花税和买卖价差），这些成本最终由基金持有人承担。高换手率也暗示基金经理缺乏坚定的长期判断，而是在不断追逐短期机会。

\subsection{幸存者偏差：你不能只看''活着''的基金}\label{ux5e78ux5b58ux8005ux504fux5deeux4f60ux4e0dux80fdux53eaux770bux6d3bux7740ux7684ux57faux91d1}

翻阅基金平台的排行榜，你会看到很多过去三五年表现优秀的基金。但那些过去表现糟糕、甚至已经被清盘或合并的基金，通常不会出现在排行榜上。这就是\textbf{幸存者偏差（survivorship
bias）}------你只看到了''活下来的赢家''，而看不到已经消失的输家。

因此，排行榜本身无法告诉你选到一只优秀基金的概率有多高。如果你在五年前随机挑选了十只主动基金，到今天可能只剩七八只还在运营，而这七八只中真正跑赢基准的可能只有两三只。这个概率，才是你面对的真实选择环境。

\subsection{如何阅读一份基金报告}\label{ux5982ux4f55ux9605ux8bfbux4e00ux4efdux57faux91d1ux62a5ux544a}

每只公募基金都会定期发布季报和年报。你不必像专业人士一样逐字逐句地研读，但有几个地方值得花上十分钟。

首先是''基金经理报告''或''管理人报告''部分。基金经理通常会在这里阐述他对过去一个季度市场环境的判断、自己的操作逻辑以及未来展望。读完后你可以问自己：他的判断是否有逻辑？他的操作是否与他的理念一致？他是否诚实地承认了错误，还是把一切问题都归咎于''市场环境不好''？

其次是''资产组合''部分。看看十大重仓股票或债券是什么、各自占比多少、与上一期相比有什么变化。这能帮你判断基金经理是否出现了风格漂移。

最后是''费用''部分。确认你实际支付的管理费、托管费以及总的费率水平。如果一只主动基金的年化费率超过
1.5\%，它需要产生远超市场水平的 Alpha
才能让费率物有所值------这在长期内对大多数基金来说都是一个艰难的目标。

\begin{center}{\color{dh-bronze}\rule{0.2\textwidth}{0.5pt}}\end{center}

主动基金可以为投资组合提供一些指数 ETF
无法替代的东西：在特定市场环境下的保护、对特定机遇的捕捉、以及某种特定的投资理念。但主动基金也是一件需要持续维护的工具------基金经理可能会离职、风格可能会漂移、超额收益可能会消失。和
ETF
的''买入并持有''不同，持有主动基金意味着你需要定期检查它是否仍然是当初你选择的那只基金。

\begin{figure}
\centering
\pandocbounded{\includegraphics[keepaspectratio]{ch18_spiva.png}}
\caption{SPIVA主动基金跑输基准比例}
\end{figure}

在下一章，我们将讨论另外几种常见的基金类型------QDII、FOF
和养老目标基金------帮助你把视野从单一市场扩展到更广阔的领域。

\newpage

\chapter{场外基金、QDII、FOF
与养老目标基金}\label{ux573aux5916ux57faux91d1qdiifof-ux4e0eux517bux8001ux76eeux6807ux57faux91d1}

\begin{center}{\color{dh-bronze}\rule{0.2\textwidth}{0.5pt}}\end{center}

\begin{dhquote}
\textbf{不是所有的基金都叫
ETF，不是所有的基金都能随时买卖。弄清它们之间的区别，才能把正确的工具放在正确的任务上。}
\end{dhquote}

\begin{center}{\color{dh-bronze}\rule{0.2\textwidth}{0.5pt}}\end{center}

在前面的几章中，我们以 ETF
为主线讨论了指数投资和被动管理。但在很多人的实际投资体验中，他们最先接触到的并不是
ETF，而是场外基金------通过银行
APP、基金平台或直销渠道申购的那些产品。这一章把几种常见的非 ETF
基金类型梳理清楚，帮助你在不同场景下做出更合适的选择。

\subsection{场内与场外：交易方式决定了你的操作自由度}\label{ux573aux5185ux4e0eux573aux5916ux4ea4ux6613ux65b9ux5f0fux51b3ux5b9aux4e86ux4f60ux7684ux64cdux4f5cux81eaux7531ux5ea6}

\textbf{场内交易} 是指通过证券账户在交易所买卖，就像买卖股票一样。ETF
属于场内交易品种------你在交易时间下单，按实时价格成交。\textbf{场外交易}
是指通过基金公司或代销平台（银行、第三方基金销售平台）进行申购和赎回。

两者的核心差异体现在以下几个方面。

\begin{longtable}[]{@{}
  >{\raggedright\arraybackslash}p{(\linewidth - 6\tabcolsep) * \real{0.2500}}
  >{\raggedright\arraybackslash}p{(\linewidth - 6\tabcolsep) * \real{0.2500}}
  >{\raggedright\arraybackslash}p{(\linewidth - 6\tabcolsep) * \real{0.2500}}
  >{\raggedright\arraybackslash}p{(\linewidth - 6\tabcolsep) * \real{0.2500}}@{}}
\toprule\noalign{}\rowcolor{dh-pink!25}
\begin{minipage}[b]{\linewidth}\raggedright
对比维度
\end{minipage} & \begin{minipage}[b]{\linewidth}\raggedright
场内（ETF）
\end{minipage} & \begin{minipage}[b]{\linewidth}\raggedright
场外（场外基金）
\end{minipage} & \begin{minipage}[b]{\linewidth}\raggedright
说明
\end{minipage} \\
\midrule\noalign{}
\endhead
\bottomrule\noalign{}
\endlastfoot
成交价格 & 实时市价，下单即成交 & 未知价原则，按当日收盘净值确认 & ETF
即时确定成本；场外基金当日不知最终成交价 \\
确认时间 & 实时到账 & T+1 至 T+2 个工作日 &
��点前提交按当日净值，三点后算下一个交易日 \\
交易成本 & 佣金 + 买卖价差（A 股 ETF 暂免印花税） & 申购费 +
赎回费（随金额和持有时间递减） & ETF
单次成本低但需注意价差；场外基金持有越久赎回费越低 \\
最低购买金额 & 一手 100 份，单价高时需数百至上千元 & 通常 10 元甚至 1
元起购 & 场外基金更适合小额定投 \\
定投支持 & 需手动操作或通过券商条件单 & 平台自动扣款，一键设置 &
场外基金定投自动化程度更高 \\
\end{longtable}

哪个更好？没有绝对的答案。如果你希望随时按实时价格交易、不希望等待确认时间，选择
ETF。如果你希望每个月自动扣款定投、不在乎确认时间、也希望能买到 10
元这样的小额，场外指数基金更适合你。

\subsection{QDII
基金：用人民币投全球}\label{qdii-ux57faux91d1ux7528ux4ebaux6c11ux5e01ux6295ux5168ux7403}

\textbf{QDII（Qualified Domestic Institutional
Investor，合格境内机构投资者）}
是一种允许境内基金投资境外市场的制度安排。通过 QDII
基金，你可以用人民币购买到追踪美股、港股、欧洲市场、日本市场甚至全球市场的基金------不需要自己开立海外证券账户，也不需自己换汇。

QDII
基金的优势显而易见：方便、门槛低、用人民币就能实现全球配置。但它也有几个需要特别注意的问题。

外汇额度限制是最重要的约束。每只 QDII
基金都有外汇投资额度上限，由外汇管理局审批。当投资者大量涌入某只热门
QDII
基金时，额度可能用完，基金公司不得不暂停申购或限制大额申购。这就是为什么有些
QDII ETF 会出现高溢价------供给被锁住，需求只能涌向二级市场。

时效性是另一个实际问题。QDII
基金的交易确认和资金到账周期比国内基金更长------赎回后通常需要五个到七个工作日才能到账。如果你有急用资金的需求，这个时差可能带来不便。

汇率波动也必须纳入考虑。QDII
基金的底层资产以外国货币计价。如果你投资了一只追踪美国股票的 QDII
基金，人民币对美元升值，你的实际回报就会被汇率损失所抵消------即使美股本身涨了，折算成人民币之后可能不涨反跌。

\subsection{FOF：让基金帮你选基金}\label{fofux8ba9ux57faux91d1ux5e2eux4f60ux9009ux57faux91d1}

\textbf{FOF（Fund of Funds，基金中的基金）}
不是直接投资股票或债券，而是投资于其他基金。可以理解为，一位专业投资者帮你挑选并组合了多只基金，构建成一揽子投资方案。

FOF
的主要优势是降低选择难度。对于完全没有时间研究基金的新手来说，买入一只
FOF 相当于同时配置了多只基金。养老目标基金中的许多产品，本质上就是 FOF。

但 FOF 也有一个需要注意的成本问题：\textbf{双重费率（double layer of
fees）}。FOF
本身会收取一层管理费，而它持有的底层基金也各自收取管理费。两层费用叠加之后，总成本通常高于直接购买指数
ETF。如果 FOF 的配置结果和一只简单的''股债组合
ETF''差不太多，那么双重费率就可能不值得付出。

\subsection{养老目标基金：用基金的方式为退休做准备}\label{ux517bux8001ux76eeux6807ux57faux91d1ux7528ux57faux91d1ux7684ux65b9ux5f0fux4e3aux9000ux4f11ux505aux51c6ux5907}

\textbf{养老目标基金（target-date fund 或 target-risk fund）}
是专为长期养老金积累而设计的一类公募基金。它最大的特点是''自动调整''------不是让你几十年如一日地持有同一个配置，而是随着目标日期临近，逐步降低高风险资产的比例，增加稳定资产的比例。

以目标日期基金为例（比如''某某 2050 养老目标基金''）。这只基金是为计划在
2050 年前后退休的投资者设计的。在 2026
年，距离退休还很远，基金的权益类资产比例可能高达 70\% 或
80\%。随着年份接近
2050，基金经理会自动逐年下调权益比例、上调债券和现金比例。到 2050
年退休时，配置已经变成一个非常保守的低波动组合。这个从高权益比例逐渐下滑到低权益比例的过程，就叫做\textbf{下滑曲线（glide
path）}。

养老目标基金适合不想自己定期调整资产配置、愿意接受''一条龙服务''的长期投资者。但它并不适合所有退休规划场景------如果距离退休还有很长的路但收入不够稳定，或者退休目标日期和基金的设定不匹配，又或者你有自己的资产配置方案并且有信心自己维护，养老目标基金就不是必需品。

\begin{center}{\color{dh-bronze}\rule{0.2\textwidth}{0.5pt}}\end{center}

到这里，我们已经覆盖了从现金到
ETF、从主动基金到养老目标基金的全部主要工具。在下一部分，我们将把这些工具组合起来------不是零散地买一堆，而是按照你的目标和人生阶段，建立一套属于自己的投资方案。

\partdivider{第五部分 · 家庭资产配置}

\newpage

\chapter{资产配置为什么比选股更重要}\label{ux8d44ux4ea7ux914dux7f6eux4e3aux4ec0ux4e48ux6bd4ux9009ux80a1ux66f4ux91cdux8981}

\begin{center}{\color{dh-bronze}\rule{0.2\textwidth}{0.5pt}}\end{center}

\begin{dhquote}
\textbf{决定你长期收益的，不是哪只股票涨了，而是你把多少钱放在了哪些篮子里。}
\end{dhquote}

\begin{center}{\color{dh-bronze}\rule{0.2\textwidth}{0.5pt}}\end{center}

在投资世界里，每天都有无数的讨论关于''哪只股票值得买入''\,``哪个行业是下一个风口''。但学术界和投资行业的大量研究表明，对长期投资收益影响最大的因素，不是选股，也不是择时，而是\textbf{资产配置（asset
allocation）}------也就是你如何在不同类型的资产之间分配你的资金。

这个结论并不是说选股和时机完全不重要，而是说，在控制了其他因素之后，资产配置解释了投资组合长期收益波动的大部分来源。决定你三十年之后是宽裕还是拮据的，不是你在某一年买到了哪只翻倍股票，而是你那三十年里持续地把钱放在了哪些资产类别上。

\subsection{相关性：资产配置的核心概念}\label{ux76f8ux5173ux6027ux8d44ux4ea7ux914dux7f6eux7684ux6838ux5fc3ux6982ux5ff5}

资产配置之所以有效，不是因为某类资产''最赚钱''，而是因为不同资产之间并不总是同步波动。有些资产在股市大跌的时候反而上涨，有些资产在通胀加剧的时候提供保护，有些资产无论市场怎么变，每个月都安静地派发利息。

资产之间的联动关系，可以用\textbf{相关性（correlation）}
来衡量。相关性系数的范围从 +1（完全同步，你涨我也涨）到
-1（完全反向，你涨我跌），中间值表示两者之间没有明显的联动。

以下是五大类资产在约 20 年间（2005---2025，年化数据）的历史相关性矩阵：

\begin{longtable}[]{@{}llllll@{}}
\toprule\noalign{}\rowcolor{dh-pink!25}
& 沪深300 & 标普500 & 中国债券 & 黄金 & 货币基金 \\
\midrule\noalign{}
\endhead
\bottomrule\noalign{}
\endlastfoot
\textbf{沪深300} & 1.00 & 0.45 & −0.12 & 0.08 & 0.02 \\
\textbf{标普500} & 0.45 & 1.00 & −0.08 & 0.15 & 0.01 \\
\textbf{中国债券} & −0.12 & −0.08 & 1.00 & 0.22 & 0.35 \\
\textbf{黄金} & 0.08 & 0.15 & 0.22 & 1.00 & 0.05 \\
\textbf{货币基金} & 0.02 & 0.01 & 0.35 & 0.05 & 1.00 \\
\end{longtable}

从矩阵中可以读出几个对资产配置有直接指导意义的结论：

\begin{itemize}
\tightlist
\item
  \textbf{中国股票与债券呈微弱的负相关（−0.12）}------这是资产配置中经典的''股债跷跷板''，意味着当股票下跌时，债券往往持平或微涨，为组合提供缓冲。
\item
  \textbf{中国股票与美国股票呈中等正相关（0.45）}------并非完全同步，但在全球性危机中两者趋于同步下跌。仅持有中美股票并不能实现充分分散。
\item
  \textbf{黄金与中国股票/债券的相关性均接近零}（0.08 /
  0.22）------说明黄金在组合中起到了独立的分散作用，而不是简单的''跟涨不跟跌''或''跟跌不跟涨''。
\item
  \textbf{中国债券与货币基金呈中等正相关（0.35）}------两者都属于固定收益/类现金资产，在利率下行环境中方向一致，但波动幅度不同。
\end{itemize}

（数据来源：Bloomberg / Wind 终端年度总回报数据，使用 Pearson
相关系数计算；替代来源：Portfolio
Visualizer（https://www.portfoliovisualizer.com）资产相关性工具。沪深300、标普500、中债综合财富指数、伦敦金、货币基金指数分别代表五类资产。）

相关性不是固定不变的------它在不同市场环境下会发生变化。例如在极端市场危机期间，几乎所有资产都可能同时下跌（相关性突然上升），这就是所谓的''风险资产联动''。理解这一点，可以帮助你避免在正常市场环境下过度自信，以为自己已经足够分散。真正的分散，应该在大多数情况下有效，并且你也要知道它在什么情况下可能失效。

\subsection{分散投资的真正含义}\label{ux5206ux6563ux6295ux8d44ux7684ux771fux6b63ux542bux4e49}

持有十只科技股，算不算分散？从数量上看，你确实分散了------任何一家公司出现问题，你的损失有限。但从风险来源来看，这十家公司都面临相似的风险：科技行业周期、利率对成长股的估值压力、监管政策变化、技术路线颠覆。当科技板块整体受到冲击时，这十只股票很可能同时下跌。这是\textbf{虚假的分散}。

真正的分散，是在不同的风险来源之间分配资金。资产之间------股票 vs 债券
vs 黄金 vs 现金；地区之间------中国 vs 美国 vs 欧洲 vs 日本 vs
新兴市场；行业之间------科技 vs 消费 vs 金融 vs 能源 vs
医疗；币种之间------人民币 vs 美元 vs
其他主要货币。每增加一个不同来源的风险敞口，你就在降低整个组合对单一风险因素的依赖。

但分散也有上限。当你持有的资产数量超过一定限度之后，继续增加资产种类带来的边际分散效果非常有限------这就是\textbf{过度分散（diworsification）}。持有五十只高度相似的基金，不会比持有两到三只覆盖面广的宽基
ETF
更安全，反而会让你的投资组合变得难以管理和理解。分散的目标不是让你成为''资产收藏家''，而是用最少的工具覆盖最广的风险面。

\subsection{再平衡：让配置长期有效}\label{ux518dux5e73ux8861ux8ba9ux914dux7f6eux957fux671fux6709ux6548}

资产配置不是''配一次管一辈子''。随着市场的涨跌，你最初配置好的比例会自然偏离。假设你最初设定的配置是
60\% 股票和 40\% 债券。一年之后，股票大涨了 30\%，债券只涨了
5\%。现在股票在组合中实际占到约 65\%，债券
35\%------你承担的风险已经高于最初的设计。

\textbf{再平衡（rebalancing）}
就是把比例拉回目标值------卖出部分股票、买入债券，或者用新增资金补齐不足的部分。再平衡的价值在于：它强迫你在资产较热的时候卖出一部分（高抛），在资产较冷的时候买入一些（低吸）。这不是预测市场------你只是在恢复最初设定的风险敞口。关于再平衡的具体操作和时机，我们将在第二十五章详细讨论。

\begin{center}{\color{dh-bronze}\rule{0.2\textwidth}{0.5pt}}\end{center}

资产配置是整个投资体系中最核心、同时也最容易被忽视的环节。很多人花大量时间研究''买什么''，却从来没停下来想一想''我为什么要把这个比例放在这类资产上''。纠正这个顺序------先确定大类配置，再在每类配置中选择具体的工具------是让投资从''碰运气''走向''有体系''的关键转变。

\begin{figure}
\centering
\pandocbounded{\includegraphics[keepaspectratio]{ch20_correlation.png}}
\caption{五大类资产历史相关性矩阵}
\end{figure}

在下一章，我们将讨论如何把钱分成几个账户------每个账户各司其职，让复杂的人生财务变得清晰有序。

\newpage

\chapter{家庭资金应该分成几个账户}\label{ux5bb6ux5eadux8d44ux91d1ux5e94ux8be5ux5206ux6210ux51e0ux4e2aux8d26ux6237}

\begin{center}{\color{dh-bronze}\rule{0.2\textwidth}{0.5pt}}\end{center}

\begin{dhquote}
\textbf{把所有的钱混在一个池子里，就像把所有的文件堆在同一个桌面上------看起来很忙，实际上很乱。}
\end{dhquote}

\begin{center}{\color{dh-bronze}\rule{0.2\textwidth}{0.5pt}}\end{center}

很多人在理财时遇到的最大困惑不是''买什么''，而是''我到底有多少钱可以用来投资''。这个问题之所以难以回答，往往不是因为数学不好，而是因为钱的功能没有被分层。所有资金混在一起，既要做日常生活支出，又要应急备用，还要追求长期增值------不同目标的资金被放在了同一个思维框架里，当然会觉得无所适从。

一个更清晰的做法，是把家庭资金按照使用目标和时间期限分成几个独立的''账户''------不一定是物理上分开的银行账户，但至少在思维和规划上是明确区隔的。

\subsection{第一层：日常开支账户}\label{ux7b2cux4e00ux5c42ux65e5ux5e38ux5f00ux652fux8d26ux6237}

\textbf{职责：}
覆盖每个月的基本生活支出------房租或房贷、水电燃气、通讯、交通、饮食、日常消费。

\textbf{期限：} 当月或当月 +
次月（即预留一到两个月的生活费在这个账户中）。

\textbf{适合的存放方式：}
银行活期存款、货币基金（如余额宝、零钱通等），确保随时可取、随时可用。

\textbf{注意事项：}
这个账户不是用来赚利息的------它的职责是保证你能正常生活。不要因为货币基金这几天年化收益高了一点点就频繁挪动，也不要因为觉得''现金浪费''而把生活费投进股票或基金。日常生活需要的资金不应该承担任何本金损失的风险。

\subsection{第二层：应急资金账户}\label{ux7b2cux4e8cux5c42ux5e94ux6025ux8d44ux91d1ux8d26ux6237}

\textbf{职责：}
应对突发事件------失业、疾病、家庭紧急支出、房屋或车辆大修。我们在第四章详细讨论过应急资金的准备标准：三到六个月的''基本生活支出''，收入越不稳定，准备越要充分。

\textbf{期限：} 随时可能动用，但正常情况下不会动用。

\textbf{适合的存放方式：}
货币基金、短期（三个月以内）银行定期存款或通知存款。要求是本金不缩水、赎回速度快（最好
T+0 或 T+1 到账）。

\textbf{注意事项：}
应急资金不是投资资金------它的收益不是你需要关心的指标。你需要关心的，是这笔钱在紧急关头能不能拿出来。不要因为''放着可惜''就把它挪去买理财产品或者投入基金------如果真出了急事，那时候的损失可能远超过多赚的几百元利息。

\subsection{第三层：中短期目标账户}\label{ux7b2cux4e09ux5c42ux4e2dux77edux671fux76eeux6807ux8d26ux6237}

\textbf{职责：}
为未来一到五年内确定或大概率会发生的支出做准备------买房首付、装修、婚礼、生育、买车、学费、计划中的大额旅行。

\textbf{期限：} 一到五年，有明确的动用计划。

\textbf{适合的存放方式：}
银行存款（定期或大额存单）、国债、中短期债券基金、偏债混合基金。任何高波动资产都不适合放入这一层------如果你计划明年买房的首付款因为股市下跌缩水了
20\%，你的购房计划就可能被迫推迟。

\textbf{注意事项：}
这一层的目标越明确、时间越确定，越应该保守。五年后才可能需要但不确定是否需要、也没有明确金额的资金，可以考虑稍高一些的波动容忍度。但底线是，不要为了''多赚一点''而用确定需求的钱去搏不确定的收益。

\subsection{第四层：长期投资账户}\label{ux7b2cux56dbux5c42ux957fux671fux6295ux8d44ux8d26ux6237}

\textbf{职责：}
为十年以上的长期目标积累财富------退休生活、子女的大学教育、财务独立。这是你真正践行''长期投资''的主战场，也是你资产增长的核心引擎。

\textbf{期限：}
十年以上，越长越好。这个账户里的资金在短期（甚至中期）内都不应该被被动用。

\textbf{适合的存放方式：} 以宽基股票指数 ETF
为���体的多元化投资组合，可以配置债券、黄金和全球资产作为分散。这部分正是本书自始至终在讨论的核心内容------从第五章的复利，到第十一章的股票收益，再到
ETF 的完整操作。

\textbf{注意事项：}
长期投资账户需要承受波动------历史上它是这样，未来也会是这样。关键在于：你要确定这个账户里的钱，你真的十年不会动。如果你一边告诉自己''这是长期投资''，一边在市场下跌
10\%
的时候就开始紧张，那你实质上是在用中期心态管理长期账户。反过来，如果你能做到无论涨跌都坚持定投、保持配置、只在人生阶段发生变化时调整比例，那么这个账户最终会成为你所有财富中增长最快的那一个。

\subsection{第五层：养老账户}\label{ux7b2cux4e94ux5c42ux517bux8001ux8d26ux6237}

\textbf{职责：}
专款专用，为退休后的生活提供持续的现金流。这一层与''长期投资账户''有一定重叠------退休准备金本质上就是最长期的长期投资。但在规划中单独列出的原因是：养老资金有专属的税收优惠账户（个人养老金账户）和提取限制，不宜与其他用途的资金混同管理。

\textbf{期限：}
距离退休有二十年、三十年甚至更长。退休之后还需要这笔钱持续产生收益。

\textbf{适合的存放方式：}
个人养老金账户中的养老目标基金、养老储蓄产品，加上自行管理的指数 ETF
组合。对于临近退休的人来说，高等级债券、年金和现金的比例需要逐步提高。

\subsection{第六层：��风险探索账户（可选）}\label{ux7b2cux516dux5c42ux98ceux9669ux63a2ux7d22ux8d26ux6237ux53efux9009}

\textbf{职责：}
如果你已经满足了前五层所有的基本需求、长期投资账户也在稳定运行，你可以考虑设立一个非常小的资金池，用于探索那些你感兴趣但风险较高的机会------某只你研究了很久的个股、某个你相信有长期潜力的新兴行业
ETF，甚至某些你做好了''可能全部亏损''心理准备的另类资产。

\textbf{期限：} 不需要在特定时间动用。甚至可以接受全部损失。

\textbf{比例：} 不超过你总可投资资产的
5\%。这个比例限定的作用是：即使这层资金全部亏光，对你的整体财务也不会产生实质性影响，但你仍然有机会去实践和学习。

\textbf{注意事项：}
探索账户不是赌博账户。它的目标仍然是学习和实践，而不是追求短期暴利。如果你发现这个账户里的钱总是很快亏掉，那不是市场的问题------可能是你还没有准备好进行独立研究。

\begin{center}{\color{dh-bronze}\rule{0.2\textwidth}{0.5pt}}\end{center}

这六层资金结构的设计，本质上是把''期限匹配''这个投资基本原则应用到家庭财务管理中。每一分钱都对应着一个目标、一个期限、一类适合的工具。当你不清楚一笔钱应该投资什么的时候，先回答它不是''哪个收益更高''，而是''这笔钱什么时候要用''。答案往往就在那个问题里。

接下来，我们将讨论不同年龄阶段应该如何调整各层资金的配置比例------因为二十五岁和五十五岁，面对的是完全不同的人生资产负债表。

\newpage

\chapter{不同年龄阶段的资产配置}\label{ux4e0dux540cux5e74ux9f84ux9636ux6bb5ux7684ux8d44ux4ea7ux914dux7f6e}

\begin{center}{\color{dh-bronze}\rule{0.2\textwidth}{0.5pt}}\end{center}

\begin{dhquote}
\textbf{资产配置没有标准答案，但每个阶段都有属于自己的解题思路。}
\end{dhquote}

\begin{center}{\color{dh-bronze}\rule{0.2\textwidth}{0.5pt}}\end{center}

在上一章中，我们建立了家庭资金的六层结构。现在，我们需要回答一个更具体的问题：在不同的人生阶段，这六层资金的比例应该如何调整？

以下分析基于大多数人在各年龄阶段的共性特征，但每个个体的实际情况------收入稳定性、家庭责任、负债水平、风险偏好、职业特点------都会影响最终的配置方案。以下比例均为参考区间，并非个性化建议。

\begin{center}{\color{dh-bronze}\rule{0.2\textwidth}{0.5pt}}\end{center}

\subsection{20---30
岁：人力资本是最大的资产}\label{ux5c81ux4ebaux529bux8d44ux672cux662fux6700ux5927ux7684ux8d44ux4ea7}

\textbf{财务特征：}
刚从学校毕业或工作不久，收入可能不高，但增长空间大。家庭责任较轻，通常没有房贷和子女负担。净资产可能为负数（助学贷款），但人力资本------未来几十年劳动收入的现值------极为丰厚。

\textbf{主要风险：}
收入不稳定，容易因为换工作、行业波动或技能不匹配而短期中断；消费诱惑多，容易陷入''赚多少花多少''的生活方式；缺乏投资经验，可能在市场高点追涨入场。

\textbf{参考资产配置区间：} 权益类资产（主要是宽基指数 ETF 和少量行业
ETF）可以占到长期投资部分的 60\% 到 80\%。债券类资产不必太多------20\%
到 30\% 足够提供基本缓冲。黄金和另类资产保持较低比例（不超过 10\%）。

\textbf{为什么这样配：}
投资期限长达三四十年，有足够的时间穿越牛熊周期。年轻时期的回撤在未来的工资收入和投资复利面前只是一个小插曲。高权益比例利用的是时间这个最核心的优势------让复利有足够长的时间发挥作用。

\textbf{不适用情形：}
如果你有一大笔必须在未来两年内动用的钱（比如计划买房的积蓄），这笔钱不应该出现在
80\%
的权益比例里。如果你的收入极不稳定且没有建立应急资金（第四章），那么在完成这两个步骤之前，不要急于配置高权益比例。

\textbf{关键点：}
这个阶段最重要的投资不是股票，而是你的职业能力------持续学习、提升技能、增加工资收入。每个月能从工资中''截留''多少用于投资，比投资收益率本身更重要。

\begin{center}{\color{dh-bronze}\rule{0.2\textwidth}{0.5pt}}\end{center}

\subsection{30---40
岁：在增长与稳定之间找平衡}\label{ux5c81ux5728ux589eux957fux4e0eux7a33ux5b9aux4e4bux95f4ux627eux5e73ux8861}

\textbf{财务特征：}
收入通常进入快速上升期，但支出也在同步攀升------房贷、婚育、子女教育开始占据现金流的主要部分。家庭资产负债表变得更加复杂：房产、房贷、教育储蓄、养老准备同时存在。保险需求明显上升------定期寿险、重疾险、医疗险应优先配置。

\textbf{主要风险：}
房贷和生活支出压力大，过度压缩投资；因为家庭责任增加而过度保守，错失了本应抓住的中期成长机会；将过多资产集中于自住房上，真实的投资资产占比过低。

\textbf{参考资产配置区间：} 权益类资产占长期投资部分的 50\% 到
70\%。债券类资产提升到 25\% 到 40\%。黄金和另类资产保持 5\% 到
10\%。应急资金应当提高到 5 到 6 个月甚至更高------因为家庭刚性支出更多。

\textbf{为什么这样配：}
投资期限仍然有二十年以上，长期增长依然是重要目标。但家庭已经经不起''因一次急事被迫割肉''的财务冲击，适度的债券和现金缓冲确保了长期投资计划不会被意外打断。

\textbf{不适用情形：}
如果你计划在未来三到五年内换一套更大的房子，首付金需要在买房前从权益账户中转移出来。如果你的工作属于高风险高回报类型（如创业、销售提成制），权益比例应当向保守端倾斜。

\textbf{关键点：}
这个阶段最容易被忽视的风险，不是市场波动，而是\textbf{人身风险}------家庭的主要收入来源者如果发生意外，整个家庭的现金流会瞬间断裂。保险配置是这个阶段不可绕过的话题。纯保障型保险------定期寿险、重疾险、医疗险、意外险------应当优先于投资型保险产品（详见第四章）。

\begin{center}{\color{dh-bronze}\rule{0.2\textwidth}{0.5pt}}\end{center}

\subsection{40---50
岁：把''积累''转向''保护''}\label{ux5c81ux628aux79efux7d2fux8f6cux5411ux4fddux62a4}

\textbf{财务特征：}
收入可能达到职业生涯的顶峰，也是积累财富的黄金十年。但与此同时，养老准备开始从''未来的事''变成''当下的事''------距离退休只有十几二十年，时间不再那么充足。子女教育支出达到高峰，父母赡养需求开始出现。

\textbf{主要风险：}
过度集中于单一行业------如果家庭的大部分收入都来自同一个行业，该行业发生结构性下滑时影响是双重甚至多重的（不仅工资减少，持股也可能缩水）；养老准备不足，高估了未来的投资收益；为追求高收益而承担超出自己能力范围的风险。

\textbf{参考资产配置区间：} 权益类资产进一步降低到 40\% 到
60\%。债券类资产提高至 30\% 到
50\%。现金和类现金资产的比例也��适当增加（5\% 到
10\%），为即将到来的退休生活预留更大的流动性安全垫。黄金保持 5\% 到 10\%
作为通胀对冲工具。

\textbf{为什么这样配：}
投资期限虽然仍然不短，但恢复时间已经不像二十岁时那样充裕。如果在这个阶段遭遇一次重大回撤（比如
2008
年级别的下跌），你可能没有二十年的时间来等待恢复。这不是说你必须''逃离股票''------股票仍然是长期增长的重要组成部分------而是你要开始用更多的债券和现金来平衡它的波动。

\textbf{不适用情形：}
如果你的退休金有强大的外部保障（如公务员养老金制度），可以适当提高权益比例。如果你已经积累了足够退休使用的资产，甚至可以降低权益比例，让组合更偏向保值和现金流。

\textbf{关键点：}
降低''单一风险集中度''是这个阶段最需要关注的结构性调整。如果你的大部分净资产都集中在你所在公司的股票（员工持股、期权）、你所处行业的资产、或者你所在城市的房产之中，那么一项行业或区域的负面变化会让你的财富遭受不必要的大幅冲击。

\begin{center}{\color{dh-bronze}\rule{0.2\textwidth}{0.5pt}}\end{center}

\subsection{50---60
岁：为退休做最后的布局}\label{ux5c81ux4e3aux9000ux4f11ux505aux6700ux540eux7684ux5e03ux5c40}

\textbf{财务特征：}
工资收入即将转为退休收入，这是资产配置从''进攻''转向''防守''的关键转折期。子女逐渐独立，教育支出减轻。开始更认真地规划退休后的收入来源------退休金、年金、投资组合的现金分配。

\textbf{主要风险：} \textbf{收益顺序风险（sequence of returns
risk）}------这是退休前后时期特有而最具破坏性的风险。简单来说，如果在临近退休的几年市场出现大跌，即使之后恢复，你的本金已经被大幅消耗，并且你很快就要开始从账户中提取资金，没有时间等待市场恢复。同样的跌幅，发生在一个人三十岁时和五十五岁时，对退休生活的影响是完全不同的量级。

\textbf{参考资产配置区间：} 权益类资产降至 30\% 到
50\%。债券类资产应成为组合的主体之一------40\% 到
55\%，以高质量的中短久期债券为主。现金和货基比例提升到 5\% 到 10\%。

\textbf{为什么这样配：}
降低权益比例的目的不是为了退出市场------你退休后可能还有二三十年的生活，通胀仍然是敌人。但降低波动、保护本金的重要性，在这个阶段已经超过了追求最高收益。你需要的是退休前最后几年本金不遭受重大损失，以及退休后能稳定地提款。

\textbf{不适用情形：}
如果你的退休金非常充裕，超过了此后生活所需的数倍，那么多出来的资金可以继续保持较高的权益比例------那些钱的期限实际上是你的继承人的投资期限，而不是你自己的。如果孩子尚未经济独立或你有持续的大额支出需求，配置需要更加保守。

\textbf{关键点：}
在退休前三年到五年内，将未来退休后三到五年的基本生活支出提前从权益账户中分离出来，放入现金或短期债券中，让它们不受市场波动的影响。这样即使你退休后的第一年市场就大跌，你也不用被迫在低点卖出股票。

\begin{center}{\color{dh-bronze}\rule{0.2\textwidth}{0.5pt}}\end{center}

\subsection{60
岁以上：从投资转向现金流管理}\label{ux5c81ux4ee5ux4e0aux4eceux6295ux8d44ux8f6cux5411ux73b0ux91d1ux6d41ux7ba1ux7406}

\textbf{财务特征：}
劳动收入停止或大幅缩减，投资组合成为主要的非退休金收入来源。医疗支出可能上升。财富传承（遗产安排）从''以后再说''变成需要认真对待的议题。

\textbf{主要风险：}
过度保守导致购买力下降------把所有钱都放在银行存款和货币基金里，三十年的通胀会把购买力消耗殆尽；过度激进导致重大亏损------被高收益产品吸引，将养老金投入高波动甚至诈骗产品；忽视流动性------在需要支付大额医疗费用时发现大部分资金被锁在非流动资产中。

\textbf{参考资产配置区间：} 权益类资产保持在 20\% 到
40\%------即使是八十岁的退休者，仍然需要一部分股票来对抗可能长达二三十年的通胀。债券和现金类资产合计
55\% 到 75\%。

\textbf{为什么这样配：}
提款需求和医疗支出的弹性很小------你每个月都需要从账户中提取一笔钱来生活。因此，波动必须受到严格约束。但与此同时，六十岁退休的人活到九十岁的概率并不低------三十年期的投资仍然需要对通胀保持警惕。

\textbf{不适用情形：}
退休金收入已完全覆盖日常生活，投资组合纯属''锦上添花''的，可以保持对更高权益比例的容忍度。健康状况不佳或有明确的遗产规划意愿的，配置可以更偏向保守和流动。

\textbf{关键点：}
提款策略比配置比例更值得关注。每年从组合中提取多少（4\%
原则是一个常见的讨论起点，但需根据实际投资环境和预期寿命调整）、按照什么顺序提取（先现金、再债券利息、最后卖股票）、是否需要购买年金，这些问题对退休生活体验的影响，往往远大于权益占比是
25\% 还是 30\%。

\begin{center}{\color{dh-bronze}\rule{0.2\textwidth}{0.5pt}}\end{center}

\subsection{年龄不是唯一标准}\label{ux5e74ux9f84ux4e0dux662fux552fux4e00ux6807ux51c6}

以上阶段分析以年龄为主线，是因为年龄与''投资期限''有最直接的关联。但投资决策不应该只看年龄。以下因素同样重要，甚至更重要：

\begin{itemize}
\tightlist
\item
  \textbf{收入稳定性。}
  同样是三十五岁，一个在一家稳定盈利的大企业工作的人，和一位自由职业者，能承受的风险敞口完全不同------即便他们的年龄完全相同。
\item
  \textbf{家庭责任。}
  两个孩子需要抚养的家庭，和没有子女负担的丁克家庭，支出刚性和风险缓冲需求差异巨大。
\item
  \textbf{负债水平。}
  高负债家庭需要更多的现金缓冲，应当优先偿还高息负债再考虑高风险投资。
\item
  \textbf{退休金保障。}
  享受稳定养老金制度的群体，比完全依赖个人投资养老的群体可以承受更高比例的权益配置。
\item
  \textbf{已有资产规模。}
  已经拥有十多年生活费的资产，和刚刚开始积累的人，风险和配置的思路应当完全不同。
\end{itemize}

配置的核心指导思想很简单：\textbf{用你输得起的钱，去博你需要的收益。给了足够的时间，让复利做它该做的事情。}

\begin{center}{\color{dh-bronze}\rule{0.2\textwidth}{0.5pt}}\end{center}

在下一章，我们将用十个真实的家庭案例，把配置原则转化为具体的决策过程------不是为了让你照搬，而是为了让你看清楚，别人怎么想、怎么选、为什么。

\newpage

\chapter{不同家庭状况的配置案例}\label{ux4e0dux540cux5bb6ux5eadux72b6ux51b5ux7684ux914dux7f6eux6848ux4f8b}

\begin{center}{\color{dh-bronze}\rule{0.2\textwidth}{0.5pt}}\end{center}

\begin{dhquote}
\textbf{学会看别人的方案不是照抄，而是学会思考自己的问题。}
\end{dhquote}

\begin{center}{\color{dh-bronze}\rule{0.2\textwidth}{0.5pt}}\end{center}

以下十个案例全部为虚构教学案例，不构成个人投资建议。它们的目标是展示''从家庭财务状况到具体配置方案''的完整思考过程。每个家庭都不同，以下配置比例均为参考区间，实际调整需要根据彼时彼刻的市场环境和个人的详细财务状况进行判断。

\begin{center}{\color{dh-bronze}\rule{0.2\textwidth}{0.5pt}}\end{center}

\subsection{案例一 · 25
岁，单身，刚工作}\label{ux6848ux4f8bux4e00-25-ux5c81ux5355ux8eabux521aux5de5ux4f5c}

\textbf{家庭背景：} 小陈，25 岁，互联网公司产品经理，税后月薪 12,000
元。单身，租房，无车。工作刚满一年。

\textbf{现金流：} 月收入 12,000 元，房租 2,500 元，日常餐饮交通购物约
5,000 元，月结余约 4,500 元。储蓄率约 37\%。

\textbf{负债：} 无。信用记录良好但从未申请过信用卡。

\textbf{已有资产：} 银行存款 30,000 元（全部在活期账户中）。

\textbf{目标：}
三年内没有购房计划。希望开始学习投资，为长期财富积累打基础。

\textbf{风险评估：}
年龄小、无负债、职业前景良好，客观风险承受能力高。投资经验为零，主观风险承受意愿中等------担心''投进去就亏了''。

\textbf{资金分层方案：} - 日常开支账户（第 1 层）：5,000
元，放在工资卡活期。 - 应急资金账户（第 2 层）：25,000 元（约 3-4
个月基本支出），放入货币基金。 - 长期投资账户（第 4
层）：从下个月开始，每月定投 3,000 元。

\textbf{参考长期投资配置：} 以一只覆盖中国市场的宽基指数 ETF 为主体（约
60\%），搭配一只全球市场 ETF（约
40\%）。暂时不单独配置债券或黄金------让时间和人力资本担任''风险缓冲器''。

\textbf{为什么这样配：} 25
岁的投资期限长达四十年，高权益比例利用时间的优势。全球 ETF
的加入提供了跨市场分散（不把所有赌注押在中国市场），也为未来了解更多投资工具打好基础。等资金积累到一定程度（比如
10 万元以上），再考虑加入债券配置。

\textbf{需要注意什么：}
不要因为市场短期上涨就加大投入超出预算的金额------定投金额应以长期可持续为准。如果工资增加，提高定投金额比提高消费水平更能带来长期收益。

\begin{center}{\color{dh-bronze}\rule{0.2\textwidth}{0.5pt}}\end{center}

\subsection{案例二 · 32
岁，双职工，准备买房}\label{ux6848ux4f8bux4e8c-32-ux5c81ux53ccux804cux5de5ux51c6ux5907ux4e70ux623f}

\textbf{家庭背景：} 小何和妻子小赵，32 岁和 30
岁，分别在一家外企和一家国企工作。家庭税后月收入合计 30,000 元。有一个 2
岁的女儿，计划三年内在一线城市购买首套房。

\textbf{现金流：} 月收入 30,000 元，月支出约 18,000 元（含房租 6,000
元、育儿费用 3,000 元），月结余约 12,000 元。储蓄率约 40\%。

\textbf{负债：} 无。

\textbf{已有资产：} 银行存款 200,000 元，货币基金 50,000 元，混合型基金
30,000 元。合计 280,000 元。两人都有基本的社保和少量企业年金。

\textbf{目标：} 三年内购买总价约 300 万元的首套房，首付加税费约 100
万元。目前储蓄和已有资产距离首付还有较大缺口，需要持续高强度储蓄。同时希望不中断长期投资。

\textbf{风险评估：}
家庭现金流稳定但即将承担大额房贷，客观风险承受能力中等。两人性格偏稳健，不想在首付金上有任何闪失。

\textbf{资金分层方案：} - 日常开支账户（第 1 层）：15,000 元，活期存款。
- 应急资金账户（第 2 层）：60,000 元（约 3 个月基本支出），货币基金。 -
中短期目标账户（第 3 层）：已有 280,000
元全部调入此层。未来三年每月从中短期账户追加存入 8,000
元。存放方式以银行大额存单、中短期国债类产品为主，不接受任何本金损失风险。
- 长期投资账户（第 4 层）：每月定投 4,000 元（含已持有的 30,000
元基金）。

\textbf{参考长期投资配置：} 简单的一到两只宽基指数 ETF（覆盖 A
股和部分海外市场），暂不配置债券------中短期目标账户的大额低风险资金已经足够起到''类债券''的缓冲作用。

\textbf{为什么这样配：}
首付目标明确且期限短（三年），资金必须确保安全------中短期目标账户的钱绝不能进入股票市场。但家庭每月有稳定的高结余，拿出一小部分（每月
4,000
元，不到月结余的三分之一）用于长期投资，既能保持投资习惯，又不影响购房目标。

\textbf{需要注意什么：}
购房前一到两年，需要逐步评估是否需要把长期投资账户中的资金也''抽出来''用于首付------如果首付缺口仍然较大，那么长期投资需要暂时让路。如果缺口已经接近填满，长期投资继续保留不动的方案更优（避免在不合适的时机变现，也保持复利继续工作）。

\begin{center}{\color{dh-bronze}\rule{0.2\textwidth}{0.5pt}}\end{center}

\subsection{案例三 · 38
岁，有房贷和子女}\label{ux6848ux4f8bux4e09-38-ux5c81ux6709ux623fux8d37ux548cux5b50ux5973}

\textbf{家庭背景：} 张先生 38 岁，在一家中型制造企业担任技术经理，妻子
36 岁，公立学校教师，税后家庭月收入合计 35,000 元。有一个 8
岁的儿子和一个 4 岁的女儿。

\textbf{现金流：} 月收入 35,000 元，支出约 25,000 元（含房贷 8,000
元、子女教育和托管费用 5,000 元、生活支出约 12,000 元），月结余约 10,000
元。储蓄率约 28\%。

\textbf{负债：} 住房贷款余额约 120 万元，剩余期限 22 年，利率 3.8\%。

\textbf{已有资产：} 自住房市值约 250 万元。银行存款和货币基金合计
150,000 元。混合型基金 40,000 元，股票 ETF 60,000 元，累计公积金账户约
200,000 元。有基本的社保和一份定期寿险、一份重疾险。

\textbf{目标：} 子女未来教育（大学或出国）准备了基础资金。希望在 55 岁到
60 岁之间实现退休，同时确保孩子教育不受自身投资结果的影响。

\textbf{风险评估：}
月结余较稳定，但负债率高，可投资资产占总资产比例偏低，大部分财富被自住房''锁定''。客观风险承受能力中等偏低------一旦发生收入中断，高额房贷和两个孩子的支出会在短时间内耗尽储蓄。

\textbf{资金分层方案：} - 日常开支账户（第 1 层）：20,000 元，活期存款。
- 应急资金账户（第 2 层）：130,000 元（约 5-6 个月基本支出），货币基金。
- 中短期目标账户（第 3 层）：无明确中短期目标（购房已完成）。 -
长期投资账户（第 4 层）：每月定投 8,000 元（含已有基金和 ETF 共 100,000
元）。 - 养老账户（第 5 层）：公积金 200,000 元 +
每月单位和个人缴存部分。

\textbf{参考长期投资配置：} 50\% 宽基指数 ETF（A 股为主）、30\% 全球市场
ETF（美股及其他地区）、20\% 中短久期债券 ETF 或债券基金。现有混合基金和
ETF 逐步转换整合到此配置中。

\textbf{为什么这样配：} 38
岁距离退休还有十七到二十二年，权益仍然应该是长期增长的主力。但高负债和子女支出降低了家庭对波动的容忍度------20\%
的债券配置提供了稳定器和再平衡弹药。全球 ETF
分散了单一市场的风险。不需要单独重仓黄金，家庭已经持有大量不动产（自住房），实物资产的敞口已经足够。

\textbf{需要注意什么：}
如果家庭现金流因任何原因变得不稳定（例如一方失业或转行），月度定投应当暂停，优先补充应急资金。如果房贷利率大幅上升（如浮动利率重定价），评估是否需要提前偿还部分贷款------此时投资和还贷的比较需要重新计算。

\begin{center}{\color{dh-bronze}\rule{0.2\textwidth}{0.5pt}}\end{center}

\subsection{案例四 · 48
岁，收入较高但养老准备不足}\label{ux6848ux4f8bux56db-48-ux5c81ux6536ux5165ux8f83ux9ad8ux4f46ux517bux8001ux51c6ux5907ux4e0dux8db3}

\textbf{家庭背景：} 周先生 48 岁，IT 行业销售总监，税后月薪 40,000
元，并有约 100,000 元的季度奖金。妻子 46 岁，全职主妇，无收入。儿子 16
岁（高中），计划出国读本科。

\textbf{现金流：} 平均税后月收入约 65,000
元（含季度奖金均摊），家庭月支出约 35,000 元（含房贷 12,000
元、儿子国际课程培训费用），月结余约 30,000 元。储蓄率约 46\%。

\textbf{负债：} 住房贷款余额约 150 万元，剩余期限 12 年，利率 3.5\%。

\textbf{已有资产：} 自住房市值约 350 万元。银行存款 300,000
元。少量混合型基金 50,000 元。有几笔银行大额存单合计 200,000
元。除此之外，几乎没有其他的投资资产。

\textbf{目标：}
发现了养老准备的巨大缺口------过去二十多年一直将收入和精力用于职场打拼和家庭生活，几乎没有系统地投资。希望在剩余的十二到十五年工作时间中''追回''时间，儿子要出国读大学还需要
100---150 万元的资金储备。

\textbf{风险评估：}
收入高但极不稳定------销售业绩受行业周期影响波动很大，家庭只有一个人的收入来源。距离退休只有十二到十七年，恢复时间不足。客观风险承受能力中等偏低。

\textbf{资金分层方案：} - 日常开支账户（第 1 层）：30,000 元，活期存款。
- 应急资金账户（第 2 层）：180,000 元（约 5-6
个月支出，考虑到只有一人有收入），货币基金。 - 中短期目标账户（第 3
层）：儿子教育费用是未来三年到五年的刚性需求。将银行存款 300,000 元中的
200,000
元划入此层，使用大额存单或中短期国债。剩余不足部分通过未来三年的高月结余继续补充。
- 长期投资账户（第 4 层）：每月投入 20,000
元，这是当前能承受的最大定投强度。

\textbf{参考长期投资配置：} 考虑到投资期限较短（12---17
年）和收入波动较大，权益比例不宜过高------45\% 宽基指数 ETF（40\% A 股 +
5\% 全球）、40\% 综合债券 ETF 或债券基金、10\% 货币基金或短期债券 ETF
作为缓冲、5\% 黄金 ETF。

\textbf{为什么这样配：}
与同龄人相比较高的权益敞口''看起来有悖于经典建议''，但基于两点判断：一是当前积累太少，完全采取保守配置（30\%
权益）的预期收益不足以在退休前积累到满足未来数十年生活所需的资产规模；二是高收入高储蓄率可以在未来十几年来持续''补仓''，平摊市场波动的风险。较高的债券比例则确保组合在被迫于低点时变现时，有足够的缓冲。

\textbf{需要注意什么：}
这是一个''需要追赶''的案例，但追赶不能变成赌博。如果未来三到五年市场出现大幅上涨，不要因为''赚到了''就放松警惕或提高消费------那是你为数不多的机会，把投资收益从纸上落袋的好时机。如果儿子留学费用超预期或出国计划提前，优先保障家庭中短期目标的实现（教育），长期投资无论如何不能挤占教育储备金。如果收入在行业下行期大幅缩水，果断将定投金额下调，将多余现金全部用来补充应急和中期储备金。

\begin{center}{\color{dh-bronze}\rule{0.2\textwidth}{0.5pt}}\end{center}

\subsection{案例五 · 58 岁，距离退休 5
年}\label{ux6848ux4f8bux4e94-58-ux5c81ux8dddux79bbux9000ux4f11-5-ux5e74}

\textbf{家庭背景：} 王女士 58 岁，一家大型消费品企业的销售总监，税后月薪
25,000 元。丈夫 60 岁，即将退休，每月退休金约 8,000
元。独生女儿已经独立，不需要经济支持。

\textbf{现金流：} 家庭税后月收入约 33,000 元，月支出约 18,000
元，月结余约 15,000 元。储蓄率约 45\%。

\textbf{负债：} 无（两年前已还清房贷）。

\textbf{已有资产：} 自住房市值约 300 万元。银行定期存款和货币基金合计
400,000 元。混合型基金 150,000 元，股票 ETF 200,000 元。企业年金约
150,000 元。无其他投资资产。

\textbf{目标：} 王女士计划在 63
岁（五年后）退休，届时主要生活来源将为两人的退休金（预计合计约
15,000---16,000
元/月）、投资组合的收益、以及企业年金的补充。希望在退休时拥有一个''不需要每天看行情''的稳健组合，每年提取一部分资金补充退休生活。

\textbf{风险评估：}
距离退休极近，收益顺序风险（见第二十二章）是最需要防范的威胁。五年后在退休的节点上如果遇到市场大跌，损失将不可逆地影响退休生活的质量。客观风险承受能力低。

\textbf{资金分层方案：} - 日常开支账户（第 1 层）：15,000 元，活期存款。
- 应急资金账户（第 2 层）：60,000 元（约 3-4 个月基本支出），货币基金。
- 退休缓冲区（第 3 层）：将退休后前五年的基本生活缺口（退休金 vs
实际生活支出）提前准备------大约需要 30---40
万元。将银行存款的大部分转入此层，放在短期定存或货币基金中，退休前绝不投入风险资产。
- 长期投资账户（第 4 层）：每月定投 10,000
元到退休为止。到时候这个账户需要转变为''退休后逐步提取''的模式。

\textbf{参考长期投资配置（退休前最后五年）：}
此阶段应逐步降低权益敞口。35\% 宽基指数 ETF��A 股为主）、45\% 综合债券
ETF 或债券基金、15\% 货币基金或短期国债类工具、5\% 黄金 ETF。

\textbf{为什么这样配：}
五年后退休的第一个目标不是''多赚钱''，而是''不亏钱''。权益敞口需要在退休前进一步降至
20\%---30\%，但过渡应该分步执行（每年降低一点），而不是一次性''一键切换''。现在的
35\% 是为退休前的过渡做准备。

\textbf{需要注意什么：}
退休前一年的年终，进行一次全面的审查------将退休后至少五年需要提取的资金全部从权益账户中移出来，放入超低波动资产中。退休之后，权益部分每年提取不超过
3.5\%---4\%
的比例，提取顺序为先现金缓冲、后债券利息、最后才是卖出股票。如果市场处于深度下跌期，暂时仅提取现金和债券利息，停止卖出股票，等待市场恢复。

\begin{center}{\color{dh-bronze}\rule{0.2\textwidth}{0.5pt}}\end{center}

\subsection{案例六 · 32
岁自由职业者，收入波动大}\label{ux6848ux4f8bux516d-32-ux5c81ux81eaux7531ux804cux4e1aux8005ux6536ux5165ux6ce2ux52a8ux5927}

\textbf{家庭背景：} 自由设计师，年收入在 12---35
万元之间波动，无固定雇主和社保以外的保障。无房产，租房居住（月租金约
4,000 元），无配偶和子女。

\textbf{核心挑战：}
收入不稳定性是第一风险。在好年份不能把全部结余投入长期资产，必须在差年份之前提前储备充足的现金缓冲。

\textbf{参考配置：} 至少 12 个月生活支出的现金缓冲（货币基金 +
短期定存）作为第一优先级，满足后多余结余按 60\% 宽基指数 ETF + 30\% 债券
ETF + 10\% 黄金 ETF
配置。不设固定月定投金额，而是在每季度末根据过去三个月的实际收入，提取税后结余的固定比例（如
60\%）投入长期账户。

\subsection{案例七 · 40
岁房产集中家庭}\label{ux6848ux4f8bux4e03-40-ux5c81ux623fux4ea7ux96c6ux4e2dux5bb6ux5ead}

\textbf{家庭背景：} 一套自住房（市值约 450 万元）+
一套投资出租房（市值约 200 万元，月租金约 3,500 元），房贷合计约 180
万元。金融资产不到 20 万元，全部为银行存款。

\textbf{核心挑战：} 超过 95\%
的净资产集中在单一资产类别（房产）和单一城市。如果房价下跌或租金市场恶化，整个家庭的财务状况将受到严重冲击。此外，出租房的流动性极差------真正需要资金时可能数月无法出手。

\textbf{参考配置：} 将月租金收入中的 50\% 作为新增投资现金流。优先建立 6
个月家庭支出现金缓冲。此后不再积累银行存款，而是按 50\% 宽基指数 ETF +
40\% 债券 ETF + 10\% 黄金 ETF 定投。长期目标是将金融资产从总资产的 3\%
逐步提升至 20\%
以上------不是卖掉房子，而是不再把所有增量财富继续''堆''在房产上。

\subsection{案例八 · 45
岁科技公司高管，持有大量股权}\label{ux6848ux4f8bux516b-45-ux5c81ux79d1ux6280ux516cux53f8ux9ad8ux7ba1ux6301ux6709ux5927ux91cfux80a1ux6743}

\textbf{家庭背景：} 年薪 + 奖金约 150
万元，另持有公司已归属限制性股票（市值约 500 万元，约占净资产的
60\%）。自住房市值约 300 万元（无贷款），金融资产约 50 万元。

\textbf{核心挑战：}
``双重集中''------人力资本（薪酬来自同一家公司）与金融资本（持股同一家公司）高度重叠。如果公司业绩恶化，可能同时面临降薪裁员的收入冲击和持股贬值的财富冲击。

\textbf{参考处置策略：}
股权解禁后，分批次、有计划地减持变现，将减持所得转入分散化组合。这里的难点不是财务计算，而是行为上的------在公司股价一路上涨时，``卖出''比''买入''难十倍。建议提前设定一个固定的减持时间表（例如每季度减持已归属股票的
10\%），将其转化为机械规则而非情绪决策。减持所得按 55\% 宽基指数 ETF +
35\% 债券 ETF + 10\% 黄金 ETF 配置。

\subsection{案例九 · 拥有 100
万元可投资资产}\label{ux6848ux4f8bux4e5d-ux62e5ux6709-100-ux4e07ux5143ux53efux6295ux8d44ux8d44ux4ea7}

\textbf{适用人群：} 30 岁左右、事业起步期的年轻投资者。

\textbf{参考配置（按第 21 章资金分层框架）：} - 日常开支账户：1
个月支出，活期存款 - 应急资金账户：3---6 个月支出，货币基金 -
长期投资账户：剩余资金，按 \textbf{70\% 宽基指数 ETF（A 股 + 全球）+
25\% 债券 ETF + 5\% 黄金 ETF} 配置 - 由于投资期限长（≥25
年），可以在年轻时承担较高权益敞口

\subsection{案例十 · 拥有 300
万元可投资资产}\label{ux6848ux4f8bux5341-ux62e5ux6709-300-ux4e07ux5143ux53efux6295ux8d44ux8d44ux4ea7}

\textbf{适用人群：} 40---45 岁、家庭与事业稳定期的投资者。

\textbf{参考配置：} - 日常 + 应急：6---12
个月支出置于货币基金和短期定存中 -
中短期目标账户（如有）：对应目标金额置于中短期债券或大额存单 -
长期投资账户：剩余资金按 \textbf{55\% 宽基指数 ETF（A 股 + 全球）+ 35\%
债券 ETF + 5\% 黄金 ETF + 5\% REITs ETF} 配置 -
此时已开始逐步提升债券比例、引入 REITs 作为额外的分散工具

\subsection{案例十一 · 拥有 1,000
万元可投资资产}\label{ux6848ux4f8bux5341ux4e00-ux62e5ux6709-1000-ux4e07ux5143ux53efux6295ux8d44ux8d44ux4ea7}

\textbf{适用人群：} 50 岁以上、接近退休或已退休的投资者。

\textbf{参考配置：} - 退休缓冲区：将退休后前 5 年生活支出缺口（退休金 vs
实际支出）的全部金额独立存放于货币基金/短期国债类工具------这部分资金在任何市场环境下都不投入风险资产
- 长期投资账户：剩余资金按 \textbf{35\% 宽基指数 ETF（A 股 + 全球）+
50\% 债券 ETF + 10\% 黄金 ETF + 5\% REITs ETF} 配置 -
每年从组合中提取不超过 3.5\%---4\%，提款顺序：现金缓冲 → 债券利息 →
黄金/REITs 分红 → 权益卖出（仅在市场非深度下跌时）

\begin{center}{\color{dh-bronze}\rule{0.2\textwidth}{0.5pt}}\end{center}

\begin{dhquote}
\textbf{以上案例六至十一的配置比例基于本书前述各章的资产历史风险收益特征和分散化原则推导，而非基于实时市场数据所做的优化建议。}
每位投资者的实际情况（税务居民身份、资金跨境限制、家庭特殊支出计划等）都可能使具体配置偏离上述参考方案。在做出重大配置决策前，请结合自身情况审慎评估或咨询持牌专业人士。
\end{dhquote}

\begin{center}{\color{dh-bronze}\rule{0.2\textwidth}{0.5pt}}\end{center}

资产配置不是比谁的方案更''精致''------它只是用最合理的方式，把你的钱放在正确的地方，让它为你的目标工作。在下一章，我们将把配置的思维转化为一套可执行的步骤：如何一步一步建立属于自己的长期投资组合。

\newpage

\chapter{如何建立自己的长期投资组合}\label{ux5982ux4f55ux5efaux7acbux81eaux5df1ux7684ux957fux671fux6295ux8d44ux7ec4ux5408}

\begin{center}{\color{dh-bronze}\rule{0.2\textwidth}{0.5pt}}\end{center}

\begin{dhquote}
\textbf{投资组合不是一份购物清单，而是一套从目标到工具再到执行的完整计划。}
\end{dhquote}

\begin{center}{\color{dh-bronze}\rule{0.2\textwidth}{0.5pt}}\end{center}

在掌握了资产配置的原理、了解了各个年龄阶段的调整逻辑、也看过了几个家庭案例之后，现在是时候把所有这些知识转化为你自己的行动方案了。

以下十二个步骤，是一个从零开始建立长期投资组合的完整流程。你不需要一次性把所有步骤都做完------可以花几天甚至几周的时间来思考。一旦这套方案建立起来，你就可以用它来指导未来很多年的投资决策，而不是每次面对市场波动都重新纠结一次。

\begin{center}{\color{dh-bronze}\rule{0.2\textwidth}{0.5pt}}\end{center}

\subsection{第一步：确定目标}\label{ux7b2cux4e00ux6b65ux786eux5b9aux76eeux6807}

在翻开任何一个基金平台或券商 APP
之前，先回答一个问题：\textbf{这笔钱将来要做什么？}

退休？孩子的教育？买一套更大的房子？还是''我希望在五十岁的时候拥有选择不做不喜欢的工作的自由''？目标不需要精确到''我要
1,538,420 元在 2047 年退休''，但需要有一个大致的方向和时间范围。

如果你投资的钱有多个目标，就把它们分开------教育基金是一份方案，养老基金是另一份方案。目标不同，期限不同，能承受的风险也不同。

\subsection{第二步：确定期限}\label{ux7b2cux4e8cux6b65ux786eux5b9aux671fux9650}

期限是你从''现在''到''需要用钱''之间的年数。这个期限决定了你的投资组合可以从多大比例投入到高波动资产中。我们在第四章已经讨论过''五年规则''------五年内确定要用到的钱，不应该投入到高波动的权益类资产中。

如果期限超过二十年，权益类资产（股票
ETF）可以占到较高比例。如果期限在五到十年之间，权益比例建议适度降低。如果期限不到五年，以债券和现金为主。

\subsection{第三步：测算未来支出}\label{ux7b2cux4e09ux6b65ux6d4bux7b97ux672aux6765ux652fux51fa}

这个步骤不需要精确到小数位，但要有一个基本的量级概念。你每个月的刚性支出是多少？三年、五年、十年之后，有没有可以预见的大额支出（买房、孩子教育、换车、重大医疗）？把这些支出列出来，标注大概的金额和时间。

这个步骤的目的，是让你在投资之前先知道：``未来哪些时候，我大概率需要从投资账户里拿钱？''知道了这些时间点，你才能避免把短期内要用的钱困在高波动资产里。

\subsection{第四步：评估风险能力}\label{ux7b2cux56dbux6b65ux8bc4ux4f30ux98ceux9669ux80fdux529b}

回到第八章我们讨论的三组概念：风险承受能力（你客观上能承受多大波动）、风险承受意愿（你心理上愿意接受多大波动）、风险承受需求（你为了实现目标''需要''承受多大波动）。

把你自己的答案写下来。不要骗自己。如果你看到账户跌 20\%
就睡不着觉，那就诚实地说出来------没有必要为了''看起来像个成熟的投资者''而承担超出自己心理极限的风险。

\subsection{第五步：选择资产类别}\label{ux7b2cux4e94ux6b65ux9009ux62e9ux8d44ux4ea7ux7c7bux522b}

基于你的投资目标、期限和风险承受能力，选择纳入你投资组合的主要资产类别。对大多数人来说，三到四种核心资产类别就足够了：

\begin{itemize}
\tightlist
\item
  宽基股票指数 ETF（核心增长引擎）
\item
  债券 ETF 或债券基金（稳定器和再平衡弹药）
\item
  黄金 ETF（防御和通胀对冲，可选）
\item
  现金和货币类工具（短期流动性和安全垫）
\end{itemize}

资产类别不是越多越好。四种搭配得当的核心资产，比八种胡乱拼凑的''收藏品''有效得多。

\subsection{第六步：设定目标配置比例}\label{ux7b2cux516dux6b65ux8bbeux5b9aux76eeux6807ux914dux7f6eux6bd4ux4f8b}

为每一种资产类别设定一个目标百分比。举个例子，一个三十二岁、投资期限超过二十年的投资者，可能会设定：

\begin{itemize}
\tightlist
\item
  宽基股票 ETF：65\%
\item
  债券 ETF：25\%
\item
  黄金 ETF：5\%
\item
  现金和货币类工具：5\%
\end{itemize}

这些百分比不是随意拍脑门想出来的------它应该源自你的风险评估结果（第四步）、借鉴相应的年龄阶段配置区间（第二十二章）、并可以参照和自己情况相似的家庭案例中的决策逻辑（第二十三章）。

\subsection{第七步：设定允许偏离区间}\label{ux7b2cux4e03ux6b65ux8bbeux5b9aux5141ux8bb8ux504fux79bbux533aux95f4}

配置比例并不需要每天都精确维持在目标值。给每一种资产设定一个允许的偏离范围，比如
``目标 ±5\%''。如果股票的配置目标是 65\%，实际浮动范围可以设定在 60\% 到
70\%
之间。只有超出了这个范围，才触发再平衡。这样做既不会因为日常波动频繁交易，也不会让配置漂移得太远。

\subsection{第八步：选择具体工具}\label{ux7b2cux516bux6b65ux9009ux62e9ux5177ux4f53ux5de5ux5177}

在大类资产比例确定之后，为每一类选择具体的投资工具。比如，在''宽基股票
ETF''这一类别下，你可能会选择一只追踪沪深 300 的 ETF 和一只追踪标普 500
的 ETF（如果你的组合需要全球配置）。在''债券 ETF''下，选择一只综合债券
ETF（覆盖国债和企业债）。在''黄金 ETF''下，选择一只实物黄金 ETF。

使用第十七章的 ETF
二十项筛选清单来筛选每一只工具。在这个阶段，应尽量选择产品结构稳定、流动性充足、费率和跟踪差异良好的产品。

\subsection{第九步：设定投入方式}\label{ux7b2cux4e5dux6b65ux8bbeux5b9aux6295ux5165ux65b9ux5f0f}

如果你的资金是一次性到位的（比如收到了一笔大额奖金或卖房回款），直接按照目标配置比例一次性买入即可。一次性投资与分批定投各有优劣------如果市场随后大涨，一次性买完全赶上了行情；如果市场随后大跌，分批买有机会降低平均成本。关键在于你对波动的容忍度和投资的期限。如果自己的承受意愿不够强，可以选择在3到6个月内分批投入，而不是在一个时间点全部进去。

如果你的资金是每月工资的一部分------这是大多数人的情况------设定一个每个月自动执行的''定投
+ 再平衡补充''方案（第二十七章会详细讨论）。

\subsection{第十步：设定卖出条件}\label{ux7b2cux5341ux6b65ux8bbeux5b9aux5356ux51faux6761ux4ef6}

在你买入任何一只资产之前，先写下你会在什么情况下卖出它。这是投资中最需要纪律约束的环节，也是很多人亏损的根源------不是因为买入时机不对，而是因为卖出的时刻完全由情绪驱动。

卖出条件可以包括：达到了某个既定的目标；投资标的本身发生了根本性恶化（不是价格跌了，而是基本面彻底坏了）；资产配置比例需要调整（按照第七步的再平衡规则执行）；人生阶段发生重大变化（如购房、退休、移民）。

写下来，存在手机备忘录或笔记本里。在市场恐慌的时候，拿出来看一遍------然后再做决定。

\subsection{第十一步：设定再平衡规则}\label{ux7b2cux5341ux4e00ux6b65ux8bbeux5b9aux518dux5e73ux8861ux89c4ux5219}

确定再平衡的频率和方法------我们将在下一章详细讨论。建议从最简单的规则开始：每年检查一次配置比例，如果实际比例偏离目标超过允许区间，就触发再平衡。

\subsection{第十二步：记录修改原因}\label{ux7b2cux5341ux4e8cux6b65ux8bb0ux5f55ux4feeux6539ux539fux56e0}

每一次你对投资组合做出任何改变------加仓、减仓、更换工具、调整配置比例------都记下来你当时做出这个决定的原因。不需要长篇大论，一句话就够了：``因为黄金过去一年涨了40\%，配置从5\%变成了8\%，超出允许区间，按再平衡规则减仓至5\%。''

这个记录的目的不是给你自己''秋后算账''，而是帮助你在一年、两年之后回头审视：你当初做的这些决定，哪些是基于理性和纪律，哪些是因为市场的噪音。投资复盘的能力（第二十八章）就是这么建立起来的。

\begin{center}{\color{dh-bronze}\rule{0.2\textwidth}{0.5pt}}\end{center}

这十二个步骤中的大部分，思考一次就够了------不需要每月、每周甚至每天重复。把框架搭好之后，你绝大部分的投资生活就简化为：每月按比例定投，每年做一次再平衡检查，人生阶段发生变化时重新评估配置目标。投资不是越忙越好------简单、稳定、可长期执行，才是你应该追求的状态。

接下来，我们把最后一步细化------再平衡到底应该怎么做，什么时候做，为什么它不是你想象中''频繁操作''。

\newpage

\chapter{什么时候应该再平衡}\label{ux4ec0ux4e48ux65f6ux5019ux5e94ux8be5ux518dux5e73ux8861}

\begin{center}{\color{dh-bronze}\rule{0.2\textwidth}{0.5pt}}\end{center}

\begin{dhquote}
\textbf{再平衡不是预测市场------它只是一道机械的数学题，强迫你在众人贪婪时卖出一些，在众人恐惧时买入一些。}
\end{dhquote}

\begin{center}{\color{dh-bronze}\rule{0.2\textwidth}{0.5pt}}\end{center}

在上一章的十二个步骤中，再平衡被列为一个需要预先设定规则的环节。这一章专门展开讨论：再平衡到底是什么、怎么做、什么时候做。

\subsection{再平衡解决了什么问题}\label{ux518dux5e73ux8861ux89e3ux51b3ux4e86ux4ec0ux4e48ux95eeux9898}

市场每天都在变化，涨涨跌跌之间，你最初精心设计的配置比例会自然偏离。假设你设定的是一个经典的''60\%
股票、40\% 债券''组合。一年之后，股票大涨 30\%，债券只涨了
5\%------现在你的组合实际变成了''65\% 股票、35\%
债券''。你承担的风险比最初设计的高了------无论你有没有意识到。

再平衡，就是把比例拉回到目标值。卖出一些涨得太多的资产（股票），买入一些涨得太慢或者跌了的资产（债券）。这个过程听起来有些反直觉------为什么要卖掉''赚钱的东西''去买''不赚钱的东西''？但仔细想想：如果你知道股票已经涨了很多、估值已经偏高，而债券相对更有吸引力，那么在这个时点''高抛低吸''是最理性的操作。再平衡只不过把你应该做的事情，写成了一条不需要犹豫的规则。

\subsection{三种再平衡方式}\label{ux4e09ux79cdux518dux5e73ux8861ux65b9ux5f0f}

\textbf{定期再平衡（calendar
rebalancing）}，顾名思义，每隔固定的时间检查一次，需要时就调整。最常见的频率是一年一次。对大多数个人投资者来说，年度再平衡在效果和操作成本之间取得了很好的平衡------既能及时纠正偏离，又不会频繁交易增加成本。有些人偏好半年检查一次，但频率超过半年到一年，交易成本和税务影响可能开始超过分散调整带来的好处。

\textbf{阈值再平衡（threshold
rebalancing）}，不是看时间，而是看偏离程度。比如设定当某类资产的实际比例超过目标值的
±5\%
时，就触发再平衡。这样可以在市场出现剧烈波动时更及时地介入------比如某一年股票大涨了
50\%，你不必等年底才来检查。但阈值设定不宜过窄（如
±2\%），否则过于频繁的再平衡会把大量时间花在微调上，效果甚至不如按兵不动。

\textbf{用新增资金再平衡（cash flow
rebalancing）}，是最节省成本的再平衡方式，尤其适合还在持续投入积累阶段的投资者。每个月定投时，不按照原有的目标比例去分钱，而是把所有新增资金都投向''当前比例已经偏低''的那类资产上。即使不卖出任何已经持有的东西，只是改变新资金的投入方向，就能逐步把组合拉回目标比例。

\subsection{什么是好的时机，什么时候不该动}\label{ux4ec0ux4e48ux662fux597dux7684ux65f6ux673aux4ec0ux4e48ux65f6ux5019ux4e0dux8be5ux52a8}

再平衡的关键在于''根据规则执行''，而不是''判断现在是再平衡的好时机''。但有一些实际考量：

牛市中的再平衡往往是最难执行的。卖掉已经大幅上涨的资产，把资金投到不温不火的债券或现金中------这在当时看来可能像''错过继续赚钱的机会''。但回想一下，所有牛市的顶部都是在多数人觉得''还会继续涨''的时候出现的。再平衡不能保证你卖在最高点，但它确保了你不会把所有浮盈都赔回去。

熊市中的再平衡逻辑是相反的：从债券和现金中取出资金，买入已经大幅下跌的股票
ETF。这需要的不是聪明，而是勇气。如果你做不到------这完全可以理解------那么保持不动也比底部恐慌卖出好。再平衡是工具，不是道德考验。在极端市场环境下，如果严格遵守再平衡规则会给你带来巨大的心理压力，略微放宽偏离容忍度也是一种合理的折中。

税务和交易成本在实际操作中也不能忽视。如果你通过场内 ETF
进行再平衡，每次买卖都会产生佣金和买卖价差的成本（虽然金额不大但频繁累积不可忽视）。如果你持有的是场外基金，赎回费也可能根据持有时间的缩短而增加。因此，优先使用''新增资金再平衡''和''分红再投资''可以达到同样的效果，同时大幅降低交易成本。

\subsection{再平衡不是预测市场，你只是在恢复风险敞口}\label{ux518dux5e73ux8861ux4e0dux662fux9884ux6d4bux5e02ux573aux4f60ux53eaux662fux5728ux6062ux590dux98ceux9669ux655eux53e3}

有一个常见的误区需要说清楚：再平衡不是为了''抄底逃顶''，它是为了把你对自己的风险敞口的控制权重新拿回到手里。如果你认为股市接下来会大涨、债券接下来会大跌，于是拒绝把股票卖掉、拒绝买入债券------那就不再是再平衡，而是主观判断市场方向。这种判断可能对也可能错，但无论如何，它已经偏离了再平衡的原有机制。

长期来看，再平衡的价值在于：在资产之间存在均值回归（mean
reversion）的趋势下，它机械地实现了部分''低买高卖''。它强制执行了一个在市场上最稀缺的习惯------在资产变得太热的时候卖出一些，在资产变得太冷的时候买入一些。

\begin{center}{\color{dh-bronze}\rule{0.2\textwidth}{0.5pt}}\end{center}

再平衡是投资组合的''守夜人''------平时不声不响，但它确保你的配置不会因为市场的极端波动而慢慢变成你完全没预料到的样子。在下一章，我们将实现很多读者期待已久的那一刻：完成第一笔真正的投资。

\partdivider{第六部分 · 真正开始投资}

\newpage

\chapter{如何完成第一笔投资}\label{ux5982ux4f55ux5b8cux6210ux7b2cux4e00ux7b14ux6295ux8d44}

\begin{center}{\color{dh-bronze}\rule{0.2\textwidth}{0.5pt}}\end{center}

\begin{dhquote}
\textbf{开始，永远比等待一个''完美时机''更重要。}
\end{dhquote}

\begin{center}{\color{dh-bronze}\rule{0.2\textwidth}{0.5pt}}\end{center}

很多人在学习投资很久之后，依然没有迈出第一步。不是因为不会买，而是因为心里一直有个声音在说''再等等、再学学、现在的市场可能不是最好的时机''。

投资没有毕业考试。所有的投资者都是在实践中学习、在波动中成长。第一笔投资最重要的目标，不是赚钱，而是启动------让你从一个''旁观学习者''变成一个''参与投资者''。

\subsection{第一笔投资之前，先确认四件事}\label{ux7b2cux4e00ux7b14ux6295ux8d44ux4e4bux524dux5148ux786eux8ba4ux56dbux4ef6ux4e8b}

\textbf{第一，你的高息负债清零了吗？} 如果还背着年利率 15\%
以上的信用卡透支或消费贷，在把这些钱还清之前，不要开始投资。

\textbf{第二，你的应急资金到位了吗？}
至少三到六个月的基本生活支出应当已经存好，安放在随时可以取用的货币基金或活期存款中。这是一道防火墙。

\textbf{第三，你的中短期资金需求已经安排好了吗？}
未来三到五年内需要用到的大笔资金------买房首付、学费------不应当进入第一笔投资。记在第二十一章的''中短期目标账户''里，用安全的方式存放。

\textbf{第四，你的心态准备好了吗？}
投进的第一笔钱，可能在你买入后的第二天就变少了------这完全正常，而且历史上经常发生。如果你还没有准备好看到账户出现浮亏，那就先用只占你资产很少一部分的金额开始------少到即使亏光了也不会影响你的生活。

\subsection{第一笔投资的完整行动流程}\label{ux7b2cux4e00ux7b14ux6295ux8d44ux7684ux5b8cux6574ux884cux52a8ux6d41ux7a0b}

\textbf{一、选择合规平台。}
确认平台持有合法金融牌照。花一点时间阅读费用说明页面：管理费、申购费、赎回费、佣金分别是多少。不推荐具体平台，但建议选择成立时间较长、用户规模较大、合规记录良好的平台。

\textbf{二、完成开户和风险测评。}
如实回答风险承受能力测评问卷------不要为了''显得能承受更大风险''而故意选高分选项。风险测评的结果是你保护自己的一道屏障。

\textbf{三、入金。}
把你已经确认''可以用来长期投资''的那笔钱转入投资账户。第一笔入金的金额不需要很大------几百元、一两千元都可以。

\textbf{四、选择你要买入的产品。}
对于第一笔投资，最简单、也最不容易出错的选择是一到两只宽基指数 ETF
或场外指数基金。你不需要在第一笔投资里完成''完美的资产配置''------配置是逐步完成的。

\textbf{五、下单并记录。}
买入之后，保存或记录下：买入日期、产品名称、买入价格、买入金额、以及选择这个产品的原因。这个简单的记录习惯，会在一年后变成复盘时最宝贵的原始数据。

\textbf{六、设置定投和复盘提醒。}
如果你计划长期持续投资，现在就可以设置每月自动扣款的定投计划。在日历里记下三个月后和一年后------在这两个时间点花十五分钟回顾投资感受和学习收获。

\subsection{一些坦诚的提醒}\label{ux4e00ux4e9bux5766ux8bdaux7684ux63d0ux9192}

\textbf{不同地区、不同平台的开户规则、税务处理和产品准入条件不同。}
本书所讨论的投资框架适用于一般投资理念，但具体到''你能买什么''\,``你需要缴什么税''，必须根据你所在的地区和法规环境自行确认。

\textbf{第一笔投资不追求''买在最低点''。}
市场明天是涨是跌，没有人知道。只要做好了长期持有的准备，今天就是一个合适的开始。

\textbf{投资之后不要频繁查看账户。}
研究表明投资者查看账户的频率越高，感知到的波动越大，做出情绪化决策的概率也越高。如果你是用月度定投方式投入资金，一个月看一次就足够了。

\begin{center}{\color{dh-bronze}\rule{0.2\textwidth}{0.5pt}}\end{center}

你完成了第一笔投资。也许只有几百元，也许只买了一只最简单的指数基金。但这一刻的意义，不在于金额的大小，而在于你完成了一个重要的心理跨越：从等待理论上的''准备好了''，到接受了''在实践中继续学习和成长''。接下来，我们进入定投------对于大多数普通投资者来说最值得认真阅读的话题。

\newpage

\chapter{定投、一次性投入与分批买入}\label{ux5b9aux6295ux4e00ux6b21ux6027ux6295ux5165ux4e0eux5206ux6279ux4e70ux5165}

\begin{center}{\color{dh-bronze}\rule{0.2\textwidth}{0.5pt}}\end{center}

\begin{dhquote}
\textbf{定投不是一种预测市场的方法，而是一种面对市场的方法。}
\end{dhquote}

\begin{center}{\color{dh-bronze}\rule{0.2\textwidth}{0.5pt}}\end{center}

如果要评选投资世界里最有争议的方法，定投一定会排在前面。有人把它视为普通投资者最适合的投资方式，也有人认为它只是营销出来的概念。真相是，定投既不是稳赚不赔的秘诀，也不是毫无意义的策略------它之所以能够流行这么多年，不是因为能神奇地提高收益率，而是因为它解决了一个几乎所有普通投资者都会遇到的问题：人很难在情绪面前始终保持理性。

\subsection{定投的行为价值}\label{ux5b9aux6295ux7684ux884cux4e3aux4ef7ux503c}

定投，全称\textbf{定期定额投资（Dollar-Cost Averaging，简称
DCA）}，指的是在固定的时间、以固定的金额持续买入同一种投资产品。每月十五号扣款一千元，无论当天市场是涨是跌，都按照计划执行。

它最大的价值不是提高收益率，而是让人真正开始投资，并且持续投下去。对于大多数普通人来说，投资最大的障碍不是''没有好的投资机会''，而是''总在等待一个好的投资时机''。市场上涨的时候担心追高，市场下跌的时候又期待更低------结果就是一年又一年过去了，钱始终留在银行账户里，一分钱都没有投入市场。

定投把''时间''这个问题从决策中拿掉了。不需要每天都问自己''今天是买入的好时机吗''。按照既定节奏执行就足够了。

\subsection{一次性投入 vs
定投}\label{ux4e00ux6b21ux6027ux6295ux5165-vs-ux5b9aux6295}

如果市场长期持续上涨------比如稳步攀升十几年------那么在最开始一次性把全部资金投入进去，最终收益会高于分批定投。原因很简单：资金更早进入市场，多享受了几年复利。这是数学上正确的结论。

但现实中的绝大多数人并不是手握一大笔现金不知道该不该投------而是每个月从工资里结余一部分，逐年积累。所以''一次性投入
vs
定投''的比较对于拿月薪的普通投资者来说，大多数时候是一个伪命题：你根本就没有一大笔钱可以一次性投进去。你的现金流本身就是分批到账的，定投只是把你的资金状况忠实地映射到投资行为上。

如果你确实有一笔大额资金（如继承、卖房回款、年终奖），犹豫该不该一次性全部投入，一个折中的做法是在三到六个月内分批投入。这样做的主要目的是心理上的------如果刚买入市场就大跌
20\%，分批的人受到的冲击会比一次性全买的人小得多。这并不能保证最终收益更高（如前所述，市场长期上涨时一次性更好），但它能增加你''坚持不半途而废''的概率------而坚持比完美的买入点重要得多。

\subsection{微笑曲线}\label{ux5faeux7b11ux66f2ux7ebf}

微笑曲线（smile curve）是定投最经典的教学图示。设想市场走出了一个 U
形：从高点下跌，在低位徘徊一段时间，然后又涨回到原来的高度。如果你在最开始一次性买入，到市场恢复原位时，你的收益是零（不考虑期间的分红）。如果你坚持定投，在下跌和低位阶段用同样的金额买到了更多的份额------当市场回到原位时，你的平均成本低于最初的价格，总的收益是正的。

下面是一个具体的微笑曲线数字案例。假设一只基金在 12
个月内的净值经历了如下 V 形走势：1.00 → 0.90 → 0.78 → 0.66 → 0.60 → 0.60
→ 0.68 → 0.76 → 0.84 → 0.92 → 1.00 →
1.00。有一位投资者在期初一次性投入了 10 万元（买入价
1.00），另一位投资者将 10 万元分为 12 等份，每月定额投入约 8,333 元。

\textbf{一次性投入（lump sum）：} 用 10 万元按净值 1.00 买入，获得
100,000 份基金。12 个月后净值回到 1.00，组合价值仍为 100,000
元------收益率为 0\%（不考虑期间分红）。

\textbf{每月定投（DCA）：} 每月固定投入约 8,333
元。在净值低迷的月份（如第 5---6 个月净值仅 0.60），用同样的 8,333
元买到了 13,889 份基金------是期初买入量的 1.67 倍。12 个月累计买入
127,545 份，平均成本约为 0.784 元（远低于期初价格 1.00 元）。当净值回到
1.00 时，组合价值为 127,545 元------收益率为 +27.5\%。

\begin{longtable}[]{@{}lllll@{}}
\toprule\noalign{}\rowcolor{dh-pink!25}
月份 & 净值 & 一次性投入（价值） & 定投当月买入（份） & 定投累计价值 \\
\midrule\noalign{}
\endhead
\bottomrule\noalign{}
\endlastfoot
1 & 1.00 & ¥100,000 & 8,333 & ¥8,333 \\
2 & 0.90 & ¥90,000 & 9,259 & ¥15,833 \\
3 & 0.78 & ¥78,000 & 10,684 & ¥22,056 \\
4 & 0.66 & ¥66,000 & 12,626 & ¥26,996 \\
5 & 0.60 & ¥60,000 & 13,889 & ¥32,875 \\
6 & 0.60 & ¥60,000 & 13,889 & ¥41,208 \\
7 & 0.68 & ¥68,000 & 12,255 & ¥55,036 \\
8 & 0.76 & ¥76,000 & 10,965 & ¥69,844 \\
9 & 0.84 & ¥84,000 & 9,921 & ¥85,530 \\
10 & 0.92 & ¥92,000 & 9,058 & ¥102,009 \\
11 & 1.00 & ¥100,000 & 8,333 & ¥119,212 \\
12 & 1.00 & ¥100,000 & 8,333 & \textbf{¥127,545} \\
\end{longtable}

在净值最终回到起点的情况下，一次性投入的收益为 0\%，而定投的收益约为
+27.5\%------这个额外的收益完全来自于在低位''捡到''了更多份额。

（数据来源：纯数学模拟。本金 10 万元、12
个月等额定投、净值序列为教学假设值 V 形路径。）

微笑曲线在数学上是成立的，但它在真实的市场中有一个重要前提：资产在经历了下跌之后，最终恢复并创出新高。如果资产的基本面已经永久性恶化------一家失去竞争壁垒的公司、一只跟踪夕阳行业的
ETF------那么无论定投多少次，``微笑曲线''都无法成立。定投不是万能药------它降低的是择时风险，不是资产选择风险。

\subsection{定投在下跌中如何降低平均成本}\label{ux5b9aux6295ux5728ux4e0bux8dccux4e2dux5982ux4f55ux964dux4f4eux5e73ux5747ux6210ux672c}

假设你每月定投 1,000 元到一只指数 ETF。第一个月净值为 1.00 元，你买到了
1,000 份。第二个月市场下跌，净值跌到 0.80 元，你用同样的 1,000 元买到了
1,250 份。两个月下来，你总共投入了 2,000 元，持有了 2,250 份，平均成本是
2,000 ÷ 2,250 ≈ 0.889 元------低于第一个月的买入价。

市场不需要回到原来的高度，你已经开始盈利了。这是定投在波动市场中积累优势的数学。但正因为这个数学在短期回本时感觉很好，很多人会高估定投''降低亏损''的力度------如果市场持续十年低迷或不断下跌，定投者的损失也在持续积累，只不过积累的速度比一次性买入的人稍慢。定投不能消灭风险，它只是把风险平摊到更长的时间维度上。

\subsection{什么情况下该暂停或终止定投}\label{ux4ec0ux4e48ux60c5ux51b5ux4e0bux8be5ux6682ux505cux6216ux7ec8ux6b62ux5b9aux6295}

定投坚持的是一种纪律，而不是对任何资产的盲目持有。在以下几种情况下，暂停或终止定投是完全理性的：

你的现金流出现了问题------失业了、收入骤降了、应急资金被消耗了需要补充。投资永远应该在生活有保障的前提下进行。这时候暂停定投、优先恢复财务安全基线，是正确的选择。

你投资的标的发生了根本性恶化------不是价格跌了，而是底层的投资逻辑彻底变了。比如一只行业
ETF
所跟踪的行业被新技术基本淘汰，或者一只主动基金的基金经理离职且接任者投资风格完全不匹配。此时继续定投等于不断往一个有漏洞的桶里倒水。

你的人生阶段发生了重大变化------比如准备买房、换工作到低薪但更有意义的领域、即将退休需要把积累的资金转向保值和现金流。定投服务于人生目标，而不是反过来。当目标变了，方法也要跟着调整。

最后，如果你发现自己因为定投反而变得天天焦虑------每次扣款都很痛苦、看到账户下跌就很后悔''为什么上个月还在投''------那么问题的根源可能不在于定投本身，而在于你的定投金额相对于你的风险承受能力来说太高了。降低每次投入的金额，或者增加组合中债券和现金的比例，可能比继续硬撑更有效地帮助你在长期坚持下来。

\begin{center}{\color{dh-bronze}\rule{0.2\textwidth}{0.5pt}}\end{center}

定投是一种帮助普通人在不确定的世界里执行长期投资的方法。它并不许诺更高的收益，但它大幅降低了一种更致命的成本------因为犹豫和恐惧而永不开始、因为一次下跌而彻底放弃。对于大多数人来说，找到一种自己能坚持十年以上的投资节奏，比精确判断每一次买入的时机重要得多。

我自己定投了快五年，最深的体会不是''微笑曲线''有多准，而是------定投最大的好处是让我不用每天早上醒来第一件事就是看行情。每个月十五号自动扣款，剩下的时间我去工作和生活。这种''不受市场情绪干扰''的感觉，比多赚几个点实在得多。

\begin{figure}
\centering
\pandocbounded{\includegraphics[keepaspectratio]{ch27_smile_curve.png}}
\caption{微笑曲线：一次性投入vs定投对比}
\end{figure}

下一章，我们将讨论一个极少被认真对待、但对长期投资者极有价值的话题：如何记录和复盘你的投资。

\newpage

\chapter{如何记录和复盘投资}\label{ux5982ux4f55ux8bb0ux5f55ux548cux590dux76d8ux6295ux8d44}

\begin{center}{\color{dh-bronze}\rule{0.2\textwidth}{0.5pt}}\end{center}

\begin{dhquote}
\textbf{好的投资决策和坏的投资决策，在短期内看起来可能一模一样。只有记录下来，你才能分辨它们。}
\end{dhquote}

\begin{center}{\color{dh-bronze}\rule{0.2\textwidth}{0.5pt}}\end{center}

你打开一个投资 APP，看到自己的账户比一年前涨了
12\%。你的第一反应是什么？很可能是高兴------``我做得不错。''但如果你的投资组合对应的基准（比如沪深
300 指数）在同期涨了 18\%，那么你实际上跑输了市场 6
个百分点。而如果你的投资组合涨了 8\%，但基准只涨了
3\%，那你其实是大幅跑赢了市场。

如果没有记录和比较的基准，你对自己投资表现的全部认知就来自于''数字是红是绿''。而颜色只能告诉你市场是涨是跌，不能告诉你自己的决策是对是错。

\subsection{投资日志：最基本的复盘工具}\label{ux6295ux8d44ux65e5ux5fd7ux6700ux57faux672cux7684ux590dux76d8ux5de5ux5177}

投资日志（investment journal）不需要很复杂。一个简单的模板可以长这样：

\begin{longtable}[]{@{}
  >{\raggedright\arraybackslash}p{(\linewidth - 10\tabcolsep) * \real{0.0833}}
  >{\raggedright\arraybackslash}p{(\linewidth - 10\tabcolsep) * \real{0.0833}}
  >{\raggedright\arraybackslash}p{(\linewidth - 10\tabcolsep) * \real{0.0833}}
  >{\raggedright\arraybackslash}p{(\linewidth - 10\tabcolsep) * \real{0.0833}}
  >{\raggedright\arraybackslash}p{(\linewidth - 10\tabcolsep) * \real{0.2917}}
  >{\raggedright\arraybackslash}p{(\linewidth - 10\tabcolsep) * \real{0.3750}}@{}}
\toprule\noalign{}\rowcolor{dh-pink!25}
\begin{minipage}[b]{\linewidth}\raggedright
日期
\end{minipage} & \begin{minipage}[b]{\linewidth}\raggedright
操作
\end{minipage} & \begin{minipage}[b]{\linewidth}\raggedright
产品
\end{minipage} & \begin{minipage}[b]{\linewidth}\raggedright
金额
\end{minipage} & \begin{minipage}[b]{\linewidth}\raggedright
理由（为什么这样做）
\end{minipage} & \begin{minipage}[b]{\linewidth}\raggedright
预期（持有多久、可能的风险）
\end{minipage} \\
\midrule\noalign{}
\endhead
\bottomrule\noalign{}
\endlastfoot
2026-03-15 & 买入 & 某沪深300 ETF & 2,000元 &
按每月定投计划执行，这是第六次定投 & 长期持有，不设具体卖出日期 \\
2026-06-20 & 卖出 & 某行业 ETF & 3,500元 &
行业基本面恶化，当初买入的逻辑已不成立 & --- \\
\end{longtable}

关键的一栏是''理由''。如果不写下你当时为什么这样做，一年后回看的时候，你可能会不自觉地用后来的结果来''解释''当初的行为------``我当时就觉得这只基金会涨''------但实际上你当初可能只是跟风买入。这是人类记忆中最常见的认知偏差之一：\textbf{事后聪明偏误（hindsight
bias）}。

\subsection{年度复盘：不要用每日盈亏评价长期策略}\label{ux5e74ux5ea6ux590dux76d8ux4e0dux8981ux7528ux6bcfux65e5ux76c8ux4e8fux8bc4ux4ef7ux957fux671fux7b56ux7565}

每天看账户涨跌，是一种效率很低的复盘方式。一个更健康的节奏是：每月定投照常进行，每季度或每半年快速扫一眼、确认没有偏离配置规则，每年做一次稍微深入的年度复盘。

年度复盘的五个核心问题：

\textbf{第一，我的长期投资组合是否还符合我的目标？}
你一年前设定的资产配置比例------比如 60\% 股票、30\% 债券、10\%
现金------现在还在那个区间里吗？如果不在了，是因为市场变动导致的正常偏离，还是因为你自己做了什么调整（调整的理由是什么）？

\textbf{第二，我的实际收益和基准相比怎么样？}
你的投资组合有一个隐含的''基准''------你选择配置了多少比例的权益类资产，你的收益就应该和''基准中同样比例的权益指数
+ 同样比例的债券指数''进行比较。如果你的组合涨了 10\%，而基准是
12\%，说明你可能在具体的产品选择或择时上出了偏差。不要只和存款利率比------那样的话任何投资看起来都''做得不错''。

\textbf{第三，我有没有做出情绪驱动的决定？}
回顾你这一年的所有操作记录------买入、卖出、调整比例。哪些是基于既定规则和理性判断的？哪些是因为''看到别人都在买''\,``连续跌了好几天太难受了''``涨太快怕错过''？诚实地面对这两个分类的比例，比你今年赚了多少钱更重要。

\textbf{第四，我的储蓄率是升了还是降了？}
投资收益率是市场给的，储蓄率是自己控制的。如果你的储蓄率在下降------收入涨了但存下来的钱没涨------那么再高的收益率也难以弥补本金的缩水。年度复盘必须包含这个非投资性但极为关键的财务指标。

\textbf{第五，未来一年我需要做什么调整？}
人生阶段快进入下一档了（比如准备买房、即将退休）？目标变了（不用存房款了、孩子毕业了）？去年定投的金额发现太高或太低？把调整写下来，作为下一年度的执行计划。

\subsection{好的结果不等于好的决策，反之亦然}\label{ux597dux7684ux7ed3ux679cux4e0dux7b49ux4e8eux597dux7684ux51b3ux7b56ux53cdux4e4bux4ea6ux7136}

这是复盘中最需要反复提醒自己的一个原则。如果你在牛市里追高买入了一只热门行业
ETF，两个月赚了
30\%------这是一个好的结果，但可能不是一个好的决策。因为它建立在''别人愿意以更高价格接盘''的运气之上。下一次同样的行为，可能带来的是严重亏损。

反过来，如果你按照资产配置计划在熊市中持续定投宽基指数
ETF，一年、两年没有赚钱甚至还浮亏------这是一个暂时不好的结果，但可能是一个非常理性的决策。因为你在市场低位收集了更多筹码，如果市场未来恢复到合理水平，这些低位买入的资产会带来真实的长期回报。

如果复盘时只用结果来评价决策，你会不断地把运气当成自己的能力，也会在运气不好的时候否定了本来正确的做法。区分这两者，是长期投资中最重要也最难掌握的技能。

\subsection{记录不是为了惩罚自己，是为了更了解自己}\label{ux8bb0ux5f55ux4e0dux662fux4e3aux4e86ux60e9ux7f5aux81eaux5df1ux662fux4e3aux4e86ux66f4ux4e86ux89e3ux81eaux5df1}

很多人在''复盘''这个词面前会本能地排斥------好像是在给自己找茬。不要这样想。投资日志和年度复盘的目的，不是在年底给自己打分------而是帮助你弄清楚一个问题：在过去的投资经历中，哪些行为让你离目标更近了，哪些让你离目标更远了。

你不需要每周、每个月都做详细复盘。但每一年花一个下午，坐下来看一看自己的账户、翻一翻自己当初写下的那些买入理由------这种定期的自我校准，会帮你逐渐从''被市场的日常噪音牵着走''变成''按照自己的既定节奏稳步前行''。

\begin{center}{\color{dh-bronze}\rule{0.2\textwidth}{0.5pt}}\end{center}

记录和复盘看似是最无聊的投资环节------既没有收益曲线的刺激，也没有找到''牛股''的成就感。但它可能是长期投资者最忠实的伙伴。因为它唯一的作用，是让你越来越了解你自己------一个比任何市场数据都重要的投资对象。

在下一部分，我们将进入全书最后一个主题：面对市场的周期和人性的波动，如何在极端时刻保持冷静。

\partdivider{第七部分 · 面对市场周期与人性}

\newpage

\chapter{牛市、熊市与市场周期}\label{ux725bux5e02ux718aux5e02ux4e0eux5e02ux573aux5468ux671f}

\begin{center}{\color{dh-bronze}\rule{0.2\textwidth}{0.5pt}}\end{center}

\begin{dhquote}
\textbf{牛市和熊市不是意外，它们是市场运转的正常节奏------就像白昼和黑夜一样不可避免。}
\end{dhquote}

\begin{center}{\color{dh-bronze}\rule{0.2\textwidth}{0.5pt}}\end{center}

在投资中，有一些话几乎每一次市场大幅波动时都会出现：``这次不一样。''2000
年互联网泡沫时有人这样说，2008 年金融危机时有人这样说，2015 年中国 A
股大幅波动时有人这样说，2020 年 3
月全球股市暴跌时还是有人这样说。但每一次，市场最终都恢复了，并在此后创出了新的高点。

这不是说''市场永远会涨''------个别市场可能经历长达十年甚至更久的低迷。也不是说''不用担心下跌''------下跌在当下是非常真实且痛苦的。真正值得认真对待的，不是''市场会不会跌''，而是''当市场下跌的时候，我会做什么''。

\subsection{牛市的陷阱}\label{ux725bux5e02ux7684ux9677ux9631}

牛市来的时候，一切看起来都很容易。你的账户在涨，你的同事在聊他刚刚买到的''牛股''，社交媒体上到处都是令人羡慕的收益截图。这是最容易做出错误决策的时候。

牛市最大的陷阱不是上涨本身，而是上涨会不断放大人的自信。本来你只计划用长期闲置的资金投资宽基指数；但赚了
20\%
之后，你开始觉得自己可能看懂了市场------于是你提高仓位，把应急资金也投了进去。本来你坚定地认为自己是一个长期投资者；但账户每天刷新高之后，你开始频繁查看行情、计算今天又赚了多少------甚至在工作时间也忍不住打开交易软件。

这些变化都不是你''做错了什么''------它们是人类大脑在持续获得正反馈时的自然反应。但投资和自然反应之间常常是对立的。牛市中感觉最''舒服''的事情------追高、加仓、频繁交易------往往是熊市中让你亏损最大的原因。

有一个值得记住的���实：每一轮牛市的顶部，都不是在大家觉得''市场太贵了''的时候出现的。顶部是在最多人觉得''还会继续涨''的时候悄悄来临的。没有人会在顶部拉响警报------甚至当市场开始下跌时，很多人还会认为''这只是回调，趁跌买入''。

\textbf{案例一：纳斯达克互联网泡沫（1999---2002）}

纳斯达克综合指数从 1995 年初的约 743 点开始，在互联网热潮的推动下，于
2000 年 3 月 10 日攀升至 5,048.62 点的历史峰值------五年间累计涨幅约
580\%。随后泡沫破裂，指数一路暴跌，至 2002 年 10 月 9 日跌至 1,139.90
点------从峰值蒸发了 77.4\%。即便指数在 2002 年底触底，此后用了整整 15
年（直到 2015 年 4 月 23 日）才重新站上 5,000 点，收复 2000
年的历史高点。如果一位投资者在 2000 年 3
月的顶部投入了资金并一直持有，他需要等待 15
年才能回本------前提是他没有在中途因为绝望而卖出。

\textbf{案例二：中国 A 股大牛市（2006---2007）及后续熊市}

上证综指从 2006 年 1 月约 1,180
点在股权分置改革和经济高速增长的推动下启动，一路飙升至 2007 年 10 月 16
日的 6,124.04 点------不到两年涨幅约
416\%。随后在内外部冲击（国内通胀高企、货币政策从紧、全球金融危机）的共同作用下，指数开启了长达一年的暴跌，至
2008 年 10 月 28 日跌至 1,664.93 点------从峰值蒸发了
72.8\%。到今天（2026 年），上证综指仍未重新触及 6,124 点。

这两个案例的教训是一致的：牛市顶部的''合理''估值，在泡沫破裂后看来，往往也是''极高''的估值。而在市场最狂热的时候，几乎所有参与者都相信自己不会成为最后接棒的那个人。

（数据来源：纳斯达克 --- Nasdaq / Bloomberg；上证综指 --- 上海证券交易所
/ Wind 终端。点数均为价格指数，不含分红。）

\subsection{牛市中值得做的事情}\label{ux725bux5e02ux4e2dux503cux5f97ux505aux7684ux4e8bux60c5}

在牛市里，最重要的事情恰恰不是''做更多''，而是''尽量少做''。

\textbf{坚持资产配置比例。}
股票仓位因大涨而超出目标比例时，按照第二十五章的再平衡规则，卖掉一部分------把超额收益落袋为安，转入债券或现金。这需要强大的心理自制力，因为每一次卖掉''还在涨''的资产都像是在主动放弃利润。但它也是确保你不会把牛市里的全部浮盈都还给熊市的最有效手段。

\textbf{继续按计划定投。}
你不需要因为市场''已经涨了很多''就停止定投------这是在尝试预测市场顶部。你也不需要''趁着牛市多投一点''------这是在用情绪替代计划。继续按既定金额定期投入，就像市场没有任何波动一样。

\textbf{不要在社交媒体上比较收益。}
牛市里总有人比你赚得更多------可能是他运气好，可能是他承担了你看不到的巨大风险，也可能是截图本身就是假的。任何比较都无法为你的长期收益增值，但它可以轻易地摧毁你的纪律和心态。

\subsection{熊市的考验}\label{ux718aux5e02ux7684ux8003ux9a8c}

如果说牛市考验的是一个人的自制力，熊市考验的就是一个人的信念。当账户连续几个月都在缩水，新闻里全是悲观的预测，身边的朋友一个接一个地清仓离场，继续坚持定投变成了一件非常困难的事情。

熊市中最大的风险不是亏钱------亏损已经发生了。最大的风险是：在亏损发生之后，因为恐惧而放弃了长期投资计划。如果你在市场底部清仓离场，等待''情况明朗了再回来''，那么你锁定了所有的浮亏，却很可能错过市场恢复时涨幅最大、最迅速的那一段。

下面的数据来自标普 500 指数在 2005---2025
年的日收益统计，直观展示了''错过最好交易日''的代价：

\begin{longtable}[]{@{}ll@{}}
\toprule\noalign{}\rowcolor{dh-pink!25}
持有策略 & 年化总回报 \\
\midrule\noalign{}
\endhead
\bottomrule\noalign{}
\endlastfoot
完整持有（每年约 252 个交易日） & \textbf{9.8\%} \\
错过了最好的 5 个交易日 & 7.1\% \\
错过了最好的 10 个交易日 & 5.2\% \\
错过了最好的 20 个交易日 & 2.8\% \\
错过了最好的 30 个交易日 & 0.9\% \\
\end{longtable}

在 20 年的投资期限内，5,040 个交易日中最好的 10 天------仅占总交易日数的
0.2\%------贡献了约 4.6 个百分点的年化回报差。如果你在那 10
天刚好不在市场里（例如因为恐慌而清仓离场），你的长期年化回报将从 9.8\%
腰斩至 5.2\%。而这最好的 10
个交易日并没有均匀分布在牛市中------它们高度集中在市场最恐慌的时段（如
2008 年末、2020 年 3---4 月），恰好是大多数投资者最想逃离的时候。

（数据来源：Bloomberg 终端，标普 500 指数每日总回报数据，2005 年 1 月
--- 2025 年 12 月。将每日回报从高到低排序后依次剔除前 N
个，计算剩余交易日几何年化回报率。含分红再投资。）

而那些最好的上涨日，往往出现在最悲观的时期------可能就在连续暴跌之后的某一天。没有人能提前知道它具体会是哪一天。

\subsection{熊市决策树}\label{ux718aux5e02ux51b3ux7b56ux6811}

当你面对熊市、尤其是跌幅已经很大的熊市时，可以按照下面的逻辑顺序逐一排查，帮助自己避免被单一的恐惧情绪主导：

\textbf{第一步：检查持仓的基本面。}
我持有的资产本身变差了吗？还是只是价格跌了？如果是前者------你持有的 ETF
所跟踪的指数中的核心公司出现了根本性恶化------那么重新评估甚至减仓是完全合理的。如果只是后者------市场整体情绪悲观导致价格下跌------那么继续持有并按计划定投，是目前最优的选择。

\textbf{第二步：检查自己的现金流。}
未来六到十二个月内，我有确定需要用到的、还没有安排的支出吗？如果有------比如计划中的买房、医疗支出------那么提前做好资金划拨（从权益账户转入现金账号），确保这些刚性支出不受市场波动影响。如果没有------也就是说这笔钱确实是长期投资------那么继续执行原计划。

\textbf{第三步：检查自己的情绪状态。}
如果继续持仓让你每天处于极度焦虑中------失眠、影响工作、情绪低沉------这本身就是一个需要认真对待的问题。很可能不是''市场太可怕了''，而是你的风险敞口超出了你实际能承受的范围。在情绪平复之后，考虑适当降低权益比例，使其更匹配你的真实风险承受意愿。

\textbf{第四步：执行提前设定的卖出规则。}
如果你在买入的时候（第二十四章第十步）已经写下了卖出条件------比如''当基金的管理费大幅上涨''或''基金经理离职且接任者风格完全不同''------那么按照规则执行。如果没有触发这些与基本面相关的卖出条件，只是因为价格下跌了，那么卖出的决定很可能来自恐惧，而不是来自规则。

\subsection{什么时候应该坚持，什么时候应该退出}\label{ux4ec0ux4e48ux65f6ux5019ux5e94ux8be5ux575aux6301ux4ec0ux4e48ux65f6ux5019ux5e94ux8be5ux9000ux51fa}

如果说长期投资是一艘船，那么''坚持''就是它的发动机，而''退出条件''是它的刹车。没有刹车的船不敢出发，没有发动机的船原地不动。

需要坚持的是：在熊市中继续持有基本面良好的资产、继续执行定投计划、继续遵守资产配置和再平衡的纪律。需要退出的是：资产基本面发生了不可逆的恶化、自己的投资期限因为人生阶段变化而不再匹配原先设定的配置比例、或者现金流出现严重问题而必须动用投资账户里的资金。

区分这两者，不是在市场暴跌的时候临时判断的------它应该在你买入每一笔资产之前就已经考虑过。记录在你的投资日志里（第二十八章）。在市场最恐慌的时候，你不需要''想清楚该怎么做''------你只需要拿出你平静时写下的东西，读一遍，然后照做。

\begin{center}{\color{dh-bronze}\rule{0.2\textwidth}{0.5pt}}\end{center}

牛市和熊市，终究都会过去。几十年后回头看，今天让你彻夜难眠的那根大阴线，不过是漫长投资长路上一个不起眼的波动。真正值得关注的不是市场本身------而是你自己的心，在周期的起伏之中，是否依然沉着。

\begin{figure}
\centering
\pandocbounded{\includegraphics[keepaspectratio]{ch29_market_cycles.png}}
\caption{错过最好交易日对长期收益的影响}
\end{figure}

在下一章，我们将深入讨论在全书中多次出现的那个主题：在投资中，最难管理的从来不是市场，而是自己。

\newpage

\chapter{投资中最难管理的是自己}\label{ux6295ux8d44ux4e2dux6700ux96beux7ba1ux7406ux7684ux662fux81eaux5df1}

\begin{center}{\color{dh-bronze}\rule{0.2\textwidth}{0.5pt}}\end{center}

\begin{dhquote}
\textbf{市场不会因为你焦虑而改变方向，但你的账户会因为你的冲动而缩水。}
\end{dhquote}

\begin{center}{\color{dh-bronze}\rule{0.2\textwidth}{0.5pt}}\end{center}

如果投资仅仅是一门技术------只要学会估值、读懂财报、掌握资产配置模型------那么每一个认真学习的人最终都应该获得良好的回报。但现实中，世界上最聪明的一批人------诺贝尔经济学奖得主、量子物理博士、顶级投行的交易员------在投资中也一样会犯低级错误。原因不在于他们不够聪明，而在于投资的敌人常常不来自外部，而来自我们自己大脑中的认知偏差（cognitive
bias）。

行为金融学（behavioral
finance）在过去几十年里系统性地研究了这些偏差是如何影响投资决策的。本章集中讨论其中最常见的十种------每一种都可能在你的投资生涯中反复出现。认识到它们的存在，是抵御它们的第一步。

\begin{center}{\color{dh-bronze}\rule{0.2\textwidth}{0.5pt}}\end{center}

\subsection{损失厌恶 · Loss
Aversion}\label{ux635fux5931ux538cux6076-loss-aversion}

\textbf{定义：}
人们在面对同样金额的损失和收益时，损失带来的痛苦大约是收益带来快乐的两倍。亏掉
1,000 元的痛苦，需要用赚到 2,000 元的喜悦来平衡。

\textbf{投资中的表现：}
投资者持有亏损的资产过久------不愿意在账面亏损时卖出（因为''卖了就真的亏了''），却过早地卖出了盈利的资产（为了''落袋为安''）。长此以往，账户里剩下的全是亏损最严重的持仓，而优质的盈利资产早已被卖掉，无法继续创造复利。

\textbf{防范规则：}
卖出决策永远应该基于资产本身的基本面和你的配置需求，而不是基于''我现在是赚还是亏''这个心理账户。在每一次考虑卖出一笔投资的时候，问自己一个问题：如果这笔投资现在是在盈利，我还会做出同样的卖出选择吗？如果答案是''会''，那就按原计划执行；如果答案是''不会''，那说明此刻驱动你决定的更可能是恐惧，而不是逻辑。

\begin{center}{\color{dh-bronze}\rule{0.2\textwidth}{0.5pt}}\end{center}

\subsection{过度自信 ·
Overconfidence}\label{ux8fc7ux5ea6ux81eaux4fe1-overconfidence}

\textbf{定义：}
人们习惯于高估自己的知识、技能和判断的准确性。在投资中表现得尤为明显------大多数投资者认为自己的投资能力高于平均水平，这在统计学上是不可能的。

\textbf{投资中的表现：}
过度交易（认为自己的每一次判断都值得执行），集中押注于少数几只自认为''研究透了''的股票，频繁择时（认为自己能判断市场短期方向），在赚了几次钱之后把运气当成了能力。

\textbf{防范规则：}
用规则约束行为，而不是用信心替代规则。设定好资产配置比例和再平衡规则之后，除非人生阶段或资产基本面发生了根本性变化，否则不频繁更改计划。投资日志（第二十八章）是觉察过度自信最有效的工具------每周回顾自己上一周做的一个投资决策，问问自己：这个决策是基于事实和逻辑，还是基于''我觉得应该没问题''？

\begin{center}{\color{dh-bronze}\rule{0.2\textwidth}{0.5pt}}\end{center}

\subsection{确认偏误 · Confirmation
Bias}\label{ux786eux8ba4ux504fux8bef-confirmation-bias}

\textbf{定义：}
人们倾向于寻找、注意和记住那些支持自己已有观点的信息，而忽略、低估或遗忘那些与自己观点矛盾的信息。

\textbf{投资中的表现：} 买入一只股票或 ETF
之后，只阅读关于这个资产''前景看好''的分析文章，而回避那些指出风险的报告。在市场上涨时找理由证明''还会继续涨''，在市场下跌时找理由证明''只是暂时调整''。

\textbf{防范规则：}
每次在做出一个投资决定之前，强迫自己做一个小练习：给自己已经持有的每一项资产，写一个''反对继续持有''的理由。哪怕最终你觉得这个理由不够有力、决定继续持有，这个思考过程本身就已经打破了你脑子里的确认回音壁。

\begin{center}{\color{dh-bronze}\rule{0.2\textwidth}{0.5pt}}\end{center}

\subsection{锚定效应 ·
Anchoring}\label{ux951aux5b9aux6548ux5e94-anchoring}

\textbf{定义：}
人们在做出判断时会过度依赖最初获得的信息------``锚点''------即使这个信息与当前决策无关或不准确。

\textbf{投资中的表现：} 以买入价格作为''锚''------``我 30
元买入的，等回到 30 元我就卖''，无视目前这只股票的真实价值可能已经远低于
30
元，或者这个价格在未来数年可能再也回不来了。以历史最高价作为''锚''------``这个
ETF 曾经到过 2.5 元，现在 1.8
元肯定便宜了''，但忽视了当时估值过高的可能，以及基本面已经发生了变化的事实。

\textbf{防范规则：}
任何资产都不该和你的''买入成本''有情感纠葛。市场不知道、也不在乎你的买入价。一个资产目前值不值当持有，取决于它未来的回报前景，和你当初花了多少钱买它毫无关系。

\begin{center}{\color{dh-bronze}\rule{0.2\textwidth}{0.5pt}}\end{center}

\subsection{羊群效应 · Herding}\label{ux7f8aux7fa4ux6548ux5e94-herding}

\textbf{定义：}
人们在不确定的情况下，倾向于跟随大多数人的行为，认为''这么多人在做，总不会都错''。

\textbf{投资中的表现：} ``大家都在买的基金，应该不差吧。''某个行业 ETF
在社交媒体上被反复讨论------一次技术革新、一场政策利好、一次新能源突破------于是越来越多的人涌入，把价格推升到远超基本面的水平。然后在热潮退去的时候，又因为''大家都在卖''而恐慌离场。

\textbf{防范规则：}
羊群效应最典型的特征就是：``大家都在说''。当你的朋友圈、办公茶水间、打车软件上的聊天内容都在谈论同一个投资话题时，这是最需要冷静的时刻，而不是最应该加入的时刻。在做任何投资之前，独立思考一个问题：如果今天没有任何人在讨论这个资产，我还会基于基本面买入它吗？

\begin{center}{\color{dh-bronze}\rule{0.2\textwidth}{0.5pt}}\end{center}

\subsection{近期偏误 · Recency
Bias}\label{ux8fd1ux671fux504fux8bef-recency-bias}

\textbf{定义：}
人们过度重视最近发生的事情，认为最近出现的趋势会一直延续下去。

\textbf{投资中的表现：}
在持续上涨几年之后，认为''股市永远涨''，提高权益配置和高波动敞口；在持续下跌几个月之后，认为''股市完蛋了''，清仓撤出所有权益。这种线性外推是人类大脑最根深蒂固的思维习惯之一，但在非线性、充满反转的市场中，它会反复惩罚追随这种直觉的投资者。

\textbf{防范规则：}
延长你观察市场的时间视角。看月度走势图会让你变得短视，看季度、年度甚至是十年期的走势图，会恢复你对''波动是常态''的感知。在你的投资日志里，每年至少回顾一次过去十年全球主要市场的走势------不是预测，只是用来提醒自己：一切最终都会过去。

\begin{center}{\color{dh-bronze}\rule{0.2\textwidth}{0.5pt}}\end{center}

\subsection{处置效应 · Disposition
Effect}\label{ux5904ux7f6eux6548ux5e94-disposition-effect}

\textbf{定义：}
损失厌恶的一个具体表现：投资者倾向于过早卖出盈利的资产，而过久持有亏损的资产。这直接导致投资组合中''优质资产被割掉、劣质资产被留下''。

\textbf{投资中的表现：}
即便你已经意识到一笔投资当初的判断是错误的，而且目前账面浮亏，还是不愿意止损------因为''割肉''这个动作在心理上等同于''正式承认自己犯了错''，而人天生厌恶承认错误。与此同时，一笔投资刚刚涨了几个点就迫不及待地落袋为安，让未来的复利永远失去了那一部分本金的时间价值。

\textbf{防范规则：}
使用再平衡规则（第二十五章）作为强制性纠偏手段------在涨得太多时卖出，目的不是为了落袋，而是恢复风险敞口；在跌得合理时买入或持有，继续持仓不是因为舍不得割肉，而是源于投资逻辑。提前写下每笔投资的持有理由（第二十四章第十步），在情绪冲动想卖出时，回看这份记录。

\begin{center}{\color{dh-bronze}\rule{0.2\textwidth}{0.5pt}}\end{center}

\subsection{FOMO · Fear of Missing Out}\label{fomo-fear-of-missing-out}

\textbf{定义：} 害怕错过他人的收益，从而在未充分评估的情况下冲动买入。

\textbf{投资中的表现：}
看到同事的账户截图、热搜话题或群里讨论某只基金''最近三个月涨了
40\%``，产生强烈的焦虑------''再不上车就晚了''。在不知道这个产品为什么涨、也不知道它还会不会继续涨的情况下，冲动买入。

\textbf{防范规则：}
停下来两分钟，打开你的投资日志，翻到你当初设定的资产配置目标和买入条件。如果你的配置里并不需要这个资产，或者你没有在平静状态下认真研究过它，那么
FOMO
就是你现在唯一在做的事情------跟随别人的兴奋，而不是执行自己的计划。

\begin{center}{\color{dh-bronze}\rule{0.2\textwidth}{0.5pt}}\end{center}

\subsection{赌徒谬误 · Gambler's
Fallacy}\label{ux8d4cux5f92ux8c2cux8bef-gamblers-fallacy}

\textbf{定义：}
错误地认为独立随机事件的概率会因为过去的结果而改变------``连续跌了五天，第六天该涨了''。

\textbf{投资中的表现：}
在没有基本面支持的情况下，纯粹因为''已经跌了很多''就认为反弹在即、加仓买入；或者因为''已经涨了很多''就认为回调将至、贸然卖出。市场的短期走势并不遵循''轮流''的逻辑------连续下跌之后不意味着''该涨了''，它可能继续下跌，也可能就此反转，两者之间没有任何统计上的必然性。

\textbf{防范规则：}
回避任何以''已经\ldots\ldots 了''开头的交易理由------``已经跌了 30\%
了''``已经连跌五天了''``已经超过历史最低点了''。便宜与否、该买不该买，不应取决于它''已经跌了多少''，而应取决于它未来的盈利创造能力和你目前组合内的配置空缺。

\begin{center}{\color{dh-bronze}\rule{0.2\textwidth}{0.5pt}}\end{center}

\subsection{结果偏误 · Outcome
Bias}\label{ux7ed3ux679cux504fux8bef-outcome-bias}

\textbf{定义：}
用结果的好坏来评判决策的质量，而不是用决策当时所依据的信息和逻辑来评估。

\textbf{投资中的表现：}
你在牛市里跟风买了一只垃圾公司股票，结果牛市还没结束它就涨了
50\%------这个结果是好的，但决策是糟糕的。下次同样的行为可能会让你血本无归。你看好某只
ETF
并坚持在熊市中定投了一年，账户还是浮亏------这个结果是暂时不好的，但决策可能是非常理性的。未来市场恢复时，这些低位积累的筹码会产生丰厚的真实回报。

\textbf{防范规则：}
这是所有偏差中最隐蔽也最危险的一个，因为它会让你把运气当作能力，把坏运气当作''投资理念有问题''。唯一有效的对抗方式就是投资日志------在每笔交易记录中写下''我当时为什么这样做''，而不是事后用涨跌来定义''好决策''还是''坏决策''。

\begin{center}{\color{dh-bronze}\rule{0.2\textwidth}{0.5pt}}\end{center}

\subsection{还需要注意的两个现代问题}\label{ux8fd8ux9700ux8981ux6ce8ux610fux7684ux4e24ux4e2aux73b0ux4ee3ux95eeux9898}

\textbf{频繁查看账户。}
智能手机的普及让查看投资账户变得异常方便------而这恰恰是问题所在。研究表明，投资者查看账户的频率越高，感知到的波动越大，做出情绪化决策的概率越高。如果你不是一个职业交易员，每天多次打开行情软件是在邀请你的大脑反复经历''微小的恐惧和微小的贪婪''------这两种微小的情绪积累在一起，等于一个被磨损的判断力。

\textbf{社交媒体比较。}
你在朋友圈或投资社区里看到的收益截图，几乎不存在''适中的''------要么是极好（拿出来晒），要么是极差（说出来寻求安慰）。而你看到的恰巧大多数是前者。这些截图构成了一个虚假的参照系------让你觉得''所有人都在赚钱，只有我不行''。这种比较没有任何用处。你的目标不是跑赢网上的陌生人，而是实现自己的财务目标。

\begin{center}{\color{dh-bronze}\rule{0.2\textwidth}{0.5pt}}\end{center}

认识这些偏差不会让你从此变成一个''完全理性的投资者''------人类的大脑不是为了投资而进化出来的，我们不可能彻底消除数百万年进化的底层反应。但你可以建立规则和习惯，在这些偏差发作的时候，给自己一个''暂停''的机会------不要马上行动，不要打开交易软件，不要看着那条红色的下跌曲线做决定。

你的投资计划，应该是在平静的星期天下牛写下来的，而不是在市场暴跌的星期三上午临时拼凑出来的。

说实话，写这一章的时候，我把自己过去犯过的每一个行为偏差都回忆了一遍。FOMO
追涨、亏了不舍得卖、赚了一点就跑------这些错误我一个没少犯过。区别只在于，后来我学会了在冲动之前先打开投资日志，看一遍自己当初写下的理由。

\begin{center}{\color{dh-bronze}\rule{0.2\textwidth}{0.5pt}}\end{center}

接下来，是这本书的最终篇章。不是关于''如何投资''的更多知识，而是关于一个更根本的问题：我们做这一切，到底是为了什么。

\partdivider{第八部分 · AI 时代的个人财务}

\newpage

\chapter{AI
记账：让每一笔钱都有处可去}\label{ai-ux8bb0ux8d26ux8ba9ux6bcfux4e00ux7b14ux94b1ux90fdux6709ux5904ux53efux53bb}

\begin{center}{\color{dh-bronze}\rule{0.2\textwidth}{0.5pt}}\end{center}

\begin{dhquote}
\textbf{记账不是为了约束你花钱，而是为了让你看清，钱到底替你在过什么样的生活。}
\end{dhquote}

\begin{center}{\color{dh-bronze}\rule{0.2\textwidth}{0.5pt}}\end{center}

2019 年，美国理财应用 Mint 公布过一个让很多人惊讶的统计：用户在注册后头
30 天里，平均会发现自己每月多花了大约 20\%
的钱------不是因为他们突然乱花钱了，而是因为他们第一次真正看清了自己把钱花在了哪里。

手动记账最大的敌人不是懒惰，而是「记不全」。一天几十笔小额支出，坐地铁、买咖啡、点外卖、随手扫码------没有几个人能靠回忆把它们逐笔记全。而记不全的账本，会给你一个错误的安全感：我以为我只花了三千，实际上花了四千。

AI 记账解决的就是这个问题。

\subsection{AI
记账到底在做什么}\label{ai-ux8bb0ux8d26ux5230ux5e95ux5728ux505aux4ec0ux4e48}

很多人第一次听到「AI
记账」时的反应是：不就是自动记账软件吗？支付宝和微信钱包早就能看消费记录了。

区别在于「分类」这件事。

传统的记账软件把交易记录罗列出来，分类基本靠两条路：要么你自己手动归类，要么它用一套固定规则（关键词匹配）自动归类------比如消费描述里带「咖啡」二字就归到餐饮。固定规则的问题在于，现实中的消费描述远比规则复杂：「某某便利」可能是餐饮也可能是日用品，「某某超市」可能同时包含食品和家用百货。

AI
记账则把「识别」变成了「理解」。它会把你的每一条交易记录和它的上下文一起处理------商户名称、金额、时间、频率、同一天的其他消费------然后判断这笔钱最可能属于哪个类别，并解释它为什么这么判断。遇到模糊的情况，它不会硬猜，而是向你确认：「这笔
45 元的『某某便利店』消费，是早餐还是日用品？」

这个能力看起来很小，但它把记账从「必须坚持的习惯」变成了「几乎不需要维护的背景服务」。而记账这件事，一旦不需要意志力来维持，它就能真正持续下去。

\subsection{记账带来的第一个改变：储蓄率变得可见}\label{ux8bb0ux8d26ux5e26ux6765ux7684ux7b2cux4e00ux4e2aux6539ux53d8ux50a8ux84c4ux7387ux53d8ux5f97ux53efux89c1}

我们在这本书第三章反复强调过：储蓄率比收入数字更重要。但储蓄率有一个尴尬的问题------它是一个需要计算才能知道的结果。很多人不是不想存钱，而是根本不知道自己的储蓄率是多少，因为不知道钱去了哪里。

AI 记账会把这件事变成实时的：

\begin{itemize}
\tightlist
\item
  每月自动汇总你的总支出，按类别拆解------餐饮占多少、居住占多少、交通占多少、娱乐占多少；
\item
  自动算出你的月度储蓄率 =（收入 − 支出）÷ 收入；
\item
  当某一类支出连续几个月明显上升时，主动提醒你，而不是等你年底复盘才发现。
\end{itemize}

这里的价值不是「管住你」，而是「让你看见」。\textbf{看得见的数字，才可能被改变。}
一个知道自己每个月餐饮花掉 1,800 元的人，和一个以为只花 1,000
元的人，面对的是完全不同的现实。前者可以做出主动选择，后者只能被动接受结果。

\subsection{AI
记账能做的，和不能做的}\label{ai-ux8bb0ux8d26ux80fdux505aux7684ux548cux4e0dux80fdux505aux7684}

能做的：

\begin{itemize}
\tightlist
\item
  \textbf{自动分类}：绝大多数交易自动归入正确的类别，模糊的向你确认；
\item
  \textbf{月度画像}：每月生成一份简单的消费报告，告诉你钱去了哪里、和上月比有什么变化；
\item
  \textbf{异常提醒}：某类支出异常升高、出现不认识的订阅扣款、账户出现大额支出时提醒你；
\item
  \textbf{目标进度}：如果你设了「每月存 2,000
  元」的目标，它会跟踪你的进度，并在月底告诉你是否达成。
\end{itemize}

不能做的，或者说需要你警惕的：

\begin{itemize}
\tightlist
\item
  \textbf{它不会替你做决定}。AI 可以告诉你餐饮花了 1,800
  元，但要不要减到 1,200
  元，是你的价值观和优先级决定的，不是它的算法决定的；
\item
  \textbf{它的分类偶尔会错}。再聪明的模型也有看走眼的时候。每周花一分钟抽查一下本周的分类，比一年后才发现账单全乱要好得多；
\item
  \textbf{它需要你的授权}。记账意味着把你的消费数据交给一个工具。选择工具之前，先确认三件事：数据存在哪里、谁有权访问、你能否随时导出并删除。
\end{itemize}

\subsection{数据安全：记账的第一条底线}\label{ux6570ux636eux5b89ux5168ux8bb0ux8d26ux7684ux7b2cux4e00ux6761ux5e95ux7ebf}

把账本交给
AI，最核心的担忧是隐私。这里有三条建议，不复杂，但值得在第一次使用前就想清楚：

\textbf{第一，永远不要在记账工具里保存密码、身份证号、银行卡密码。}
记账只需要交易流水，不需要你的登录凭证。

\textbf{第二，选择数据边界清晰的产品。}
优先选择「数据加密存储、可导出、可删除」的工具。如果一个工具要求你授权读取短信或邮箱里的所有内容，而不仅仅是账单------想清楚它到底需要多少权限。

\textbf{第三，定期导出自己的数据。}
记账工具可以换，但你的财务历史是你自己的。每半年把数据导出一份存档，即使将来更换工具，你的账本也不会丢失。

\subsection{小练习：让 AI
记账跑一个月}\label{ux5c0fux7ec3ux4e60ux8ba9-ai-ux8bb0ux8d26ux8dd1ux4e00ux4e2aux6708}

\begin{enumerate}
\def\labelenumi{\arabic{enumi}.}
\tightlist
\item
  选择一款记账工具，完成初始设置，只授权必要的交易数据读取；
\item
  连续使用一个月，不需要刻意改变消费习惯------先让它记录真实的你；
\item
  月底打开月度报告，写下三件事：我的钱主要花在哪些类别？哪一类支出和我以为的不一样？下个月我想主动调整哪一个？
\end{enumerate}

一个月后你会发现的，大概率不是「我花得太多了」，而是「我以前对自己的生活了解得太少了」。

\begin{center}{\color{dh-bronze}\rule{0.2\textwidth}{0.5pt}}\end{center}

记账是财务管理的「眼睛」------它让你看清现状。而看清现状之后，下一步是看向未来：你的人生目标，应该被翻译成什么样的财务计划？这就是下一章要讨论的：AI
长期规划。

\newpage

\chapter{AI
长期规划：把人生目标翻译成财务数字}\label{ai-ux957fux671fux89c4ux5212ux628aux4ebaux751fux76eeux6807ux7ffbux8bd1ux6210ux8d22ux52a1ux6570ux5b57}

\begin{center}{\color{dh-bronze}\rule{0.2\textwidth}{0.5pt}}\end{center}

\begin{dhquote}
\textbf{长期规划不是预测未来，而是把「我希望的生活」翻译成「我现在要做的事」。}
\end{dhquote}

\begin{center}{\color{dh-bronze}\rule{0.2\textwidth}{0.5pt}}\end{center}

很多人对「长期规划」望而生畏，因为它听起来像一道没有标准答案的大题：我要在
40 岁前买房，我要在 60
岁退休，我要给孩子存一笔教育金------然后呢？然后就没有然后了。目标停留在愿望层面，因为没有人和方法帮你把它拆解成可执行的数字。

AI
长期规划要做的事情，恰恰是把这些模糊的愿望，翻译成具体的财务数字和行动步骤。

\subsection{为什么「愿望」必须变成「数字」}\label{ux4e3aux4ec0ux4e48ux613fux671bux5fc5ux987bux53d8ux6210ux6570ux5b57}

我们在第五章讨论复利时算过一笔账：同样是每月存 500 元，25 岁开始和 35
岁开始，到 60
岁的差距是巨大的。但这里有一个更基础的问题：\textbf{你凭什么确定自己每个月该存
500 元，而不是 300 元或 1,000 元？}

答案是：由你的目标反推。

如果你 25 岁，希望 60 岁退休时有一笔 300 万元的退休金，假设年化收益率
7\%，那么从现在起每月需要投入多少？如果你希望 35 岁前攒下 50
万元首付，在收益率 3\% 的保守假设下，每月又要存多少？

这些计算不复杂，但很少有人真正做------因为手动算很麻烦，算出来之后还要面对「原来我每个月要存这么多」的落差。AI
的价值就在这里：它把「目标 → 每月投入 → 时间 →
收益率」之间的换算变成一个即时的、可以反复试算的工具。

\subsection{AI
长期规划的三个核心能力}\label{ai-ux957fux671fux89c4ux5212ux7684ux4e09ux4e2aux6838ux5fc3ux80fdux529b}

\textbf{第一，目标反推。} 你说一个目标（时间 +
金额），它告诉你每月需要存多少；你说一个每月能存的上限，它告诉你这个目标在什么时间尺度上可行。双向都可以试，直到找到一个你既愿意接受、又能够执行的点。

\textbf{第二，情景模拟。} 同样一笔投入，收益率 5\%、7\%、9\%
三种假设下，20
年后分别是多少？中间如果遇到一次大回撤（我们在第七章讨论过，市场下跌
30\% 不是罕见事件），你的账户会经历什么？AI
不会给你一个「精确预测」------没有人能预测------但它可以帮你提前看清不同假设下的人生版本，让你在低风险的状态下做好心理和财务上的准备。

\textbf{第三，里程碑跟踪。}
规划不是做一次就结束的。每年更新一次你的收入、支出、资产情况，AI
会重新计算你的进度，并告诉你：按当前节奏，你的目标会提前还是延后？这个「进度感」非常重要------它让一个十年期的目标，变成了每年都能看见的小步前进。

\subsection{一个例子：把「想退休」变成计划}\label{ux4e00ux4e2aux4f8bux5b50ux628aux60f3ux9000ux4f11ux53d8ux6210ux8ba1ux5212}

假设你 30 岁，月收入 15,000 元，每月能存 4,000 元，希望 60
岁退休。你可以这样和 AI 对话：

\begin{dhquote}
「我 30 岁，每月能存 4,000 元，假设长期年化收益 7\%，到 60
岁我会有多少钱？」
\end{dhquote}

它会算给你看：40 万元本金（4,000 × 12 × 10\ldots{} 不，是 4,000 × 12 ×
30 = 144 万元本金），在 7\% 复利下大约会变成 480 万元左右。

然后你可以继续追问：

\begin{dhquote}
「如果我想 60 岁时有 600 万元，我需要每月存多少？」
\end{dhquote}

它会告诉你：大约每月 5,200
元。这时候你就有了一条清晰的决策链------要么提高每月储蓄，要么接受稍低的退休目标，要么延长工作时间。\textbf{关键是，这些选择在你
30 岁的时候就摆在桌面上，而不是 55 岁时才被迫面对。}

\subsection{规划工具的边界：假设不是承诺}\label{ux89c4ux5212ux5de5ux5177ux7684ux8fb9ux754cux5047ux8bbeux4e0dux662fux627fux8bfa}

使用 AI
长期规划时，有一条必须刻在心里的原则：\textbf{模型的输出质量取决于输入的假设，而假设不等于事实。}

7\%
的年化收益率是一个合理的长期教学假设（我们在序章里用过），但它不是承诺。真实的市场会有牛市和熊市，你的收入会变化，你的生活会变化。所以：

\begin{itemize}
\tightlist
\item
  把 AI 给你的数字当作「参考区间」而不是「确定答案」；
\item
  每年重新规划一次，而不是二十年前算完就不再更新；
\item
  对任何「保证收益」的说法保持警惕------无论来自 AI
  还是来自推销员，承诺收益的都值得怀疑。
\end{itemize}

\subsection{和全书其他章节的连接}\label{ux548cux5168ux4e66ux5176ux4ed6ux7ae0ux8282ux7684ux8fdeux63a5}

AI 长期规划不是孤立存在的工具，它调用的全部是你在这本书里学过的框架：

\begin{itemize}
\tightlist
\item
  第四章的应急资金和负债处理------规划的第一步永远是先打好安全垫；
\item
  第五章的复利计算------反推每月投入的数学基础；
\item
  第二十一章的资金分层------你的目标应该落在哪一层账户里；
\item
  第二十二章的年龄阶段配置------不同人生阶段的收益率假设和风险承受力不同；
\item
  第二十四章的组合构建------目标确定后，具体用什么工具去实现。
\end{itemize}

换句话说，AI
是帮你把这些框架\textbf{算得快一点、看得清一点}的助手，而不是替代这些框架的魔法。

\subsection{小练习：做一次目标反推}\label{ux5c0fux7ec3ux4e60ux505aux4e00ux6b21ux76eeux6807ux53cdux63a8}

\begin{enumerate}
\def\labelenumi{\arabic{enumi}.}
\tightlist
\item
  写下你最重要的一个长期目标（时间 + 大致金额）；
\item
  用 AI
  工具试算：按当前储蓄能力，这个目标需要多少年、或需要每月多存多少；
\item
  调整假设（收益率
  ±2\%），看看结果的变化幅度------这会让你理解「数字对假设的敏感度」；
\item
  把最终采用的数字，写进你的投资计划书（第二十四章的十二步流程）。
\end{enumerate}

\begin{center}{\color{dh-bronze}\rule{0.2\textwidth}{0.5pt}}\end{center}

规划解决了「往哪走」的问题，记账解决了「现在在哪」的问题。那么最后一个问题自然而然地浮现：当
AI
开始参与得越来越深，它能不能直接帮我做投资决策？答案需要分两半来说------这就是下一章的主题。

\newpage

\chapter{AI 投资辅助与边界：让 AI
当副驾驶，而不是驾驶员}\label{ai-ux6295ux8d44ux8f85ux52a9ux4e0eux8fb9ux754cux8ba9-ai-ux5f53ux526fux9a7eux9a76ux800cux4e0dux662fux9a7eux9a76ux5458}

\begin{center}{\color{dh-bronze}\rule{0.2\textwidth}{0.5pt}}\end{center}

\begin{dhquote}
\textbf{AI
可以帮你把信息整理清楚，但「用你的钱去冒险」这个决定，只能由你来做。}
\end{dhquote}

\begin{center}{\color{dh-bronze}\rule{0.2\textwidth}{0.5pt}}\end{center}

这本书写到这里，你已经知道：投资长期看盈利、短期看情绪（第十一章）；资产配置比选股更重要（第二十章）；最大的敌人是你自己（第三十章）。那么一个很自然的问题是：既然
AI 这么聪明，能不能让它直接帮我选基金、做交易、管账户？

答案会让一些人失望：\textbf{成熟的 AI 使用方式，不是把方向盘交给
AI，而是让 AI 坐在副驾驶的位置上。}

\subsection{AI
在投资里真正擅长的事}\label{ai-ux5728ux6295ux8d44ux91ccux771fux6b63ux64c5ux957fux7684ux4e8b}

AI
在投资中的可靠能力，集中在「信息处理和纪律执行」上，而不是「预测市场」上。

\textbf{第一，信息整理。}
一份几十页的基金年报，一份复杂的监管公告，一堆观点冲突的财经文章------AI
可以快速提取要点、整理成结构化摘要，帮你把「读资料」的时间从几小时压缩到几分钟。但请注意：它做的是\textbf{提炼}，不是\textbf{解读}。它给你的是资料里的客观要点，而不是「该不该买」的结论。

\textbf{第二，检查清单执行。} 这本书第十七章给了你一张 ETF
二十项筛单，第二十四章给了你十二步组合构建流程。这些清单最大的问题是：人很容易在情绪上来的时候跳过几步。AI
可以扮演一个严格的项目经理------「你确认过费率了吗？你核对过跟踪误差了吗？你写下卖出条件了吗？」它不会替你决定，但它会确保你没有漏掉该走的步骤。

\textbf{第三，复盘助手。} 第二十八章的投资日志和年度复盘，AI
可以帮助你整理：把你一年来的每笔操作按时间线排好，对比你的组合收益和基准（我们说过，要和沪深
300
比，而不是和存款利率比），标记出哪些操作没有写下理由。它把「复盘」这件事的成本大幅降低，让你更愿意做------而愿意复盘，本身就是长期投资者的重要优势。

\textbf{第四，纪律提醒。} 第二十九章讲过牛市和熊市的情绪陷阱。AI
可以在你试图偏离计划的时候提醒你：「根据你的投资计划书，你设定的股票配置比例是
60\%，现在是
72\%------按第二十五章的再平衡规则，你应该考虑卖出超出部分，而不是继续加仓。」它不一定拦得住你，但它会成为一个不容易被情绪说服的声音。

\subsection{AI
不应该被用于的事}\label{ai-ux4e0dux5e94ux8be5ux88abux7528ux4e8eux7684ux4e8b}

以下是 AI 在投资中的三条红线，无论工具的宣传多么诱人，都值得保持警惕：

\textbf{第一，不要让 AI 替你选具体标的。}
任何告诉你「这只基金值得买」「这只股票要涨」的
AI，本质上都是在做一个预测。而我们在第十一章已经论证过：短期预测的准确率，并不会因为说话的是
AI 而提高。AI
可以帮你分析一只基金的费率和跟踪误差，但「买不买」是你的决定。

\textbf{第二，不要把账户权限交给 AI。}
任何要求你提供交易密码、允许自动交易、代为管理账户的「AI
投资工具」，都在把你的资产置于它无法承担责任的境地。你可以用 AI
记录和分析，但交易按钮，永远应该由你自己按下。

\textbf{第三，不要相信「AI 保证收益」。} 如果有一个 AI
声称自己找到了「稳赚不赔」的策略------无论它叫量化、智能投顾还是别的名字------请记住我们在第七章讨论过的风险本质：高收益承诺往往伴随着你看不见的高风险。承诺收益，是投资世界里最响亮的警报。

\subsection{研究侧的
AI：离你的钱很远的另一种存在}\label{ux7814ux7a76ux4fa7ux7684-aiux79bbux4f60ux7684ux94b1ux5f88ux8fdcux7684ux53e6ux4e00ux79cdux5b58ux5728}

讲完三条红线，你可能会问：那是不是所有 AI
投资工具都该敬而远之？其实判断标准从来不是「它是不是
AI」，而是「它碰不碰你的钱、替不替你做决定」。沿着这个标准，投资里的 AI
可以粗略分成两类：

\textbf{交易侧的 AI}
替你看盘、选标的、甚至替你下单------它离你的钱太近，正是上面三条红线要拦住的对象。

\textbf{研究侧的 AI}
则完全相反：它帮你把一家公司、一个行业、一类资产的信息收集起来，整理成结构化的资料------财报要点、历史数据、不同标的的对比。它不接触你的任何账户，也不输出「买还是卖」，更像一个不知疲倦的助理研究员。把资料整理这类重复劳动交给它，能省下大量时间；但研究完「该不该买」，仍然要由你自己在纸上写下理由后按下确认键。

这里有一个诚实的说明：研究侧的工具我自己也在用。我日常的行业与公司研究会交给自研的开源研究工具
AlphaForge
来处理------它是一个把「信息加工成结构化投资知识」的工作台，让 AI
在归档、整理和对比上搭把手。我把它开源，是希望读者能看到一种健康的边界示范：它给我的是整理好的事实，而不是替我按下交易按钮的冲动。工具要不要用、用哪一家的，请你自行判断------但如果决定用
AI
辅助研究，请守住这一条：\textbf{它只能在你愿意暴露的数据范围内工作，且永远不碰你的交易账户。}

\subsection{一个健康的 AI
使用姿势}\label{ux4e00ux4e2aux5065ux5eb7ux7684-ai-ux4f7fux7528ux59ffux52bf}

让我们把 AI
放进一个健康的投资流程里，看看它应该出现在哪些环节、不该出现在哪些环节：

\begin{longtable}[]{@{}lll@{}}
\toprule\noalign{}\rowcolor{dh-pink!25}
环节 & AI 可以参与 & AI 不应决定 \\
\midrule\noalign{}
\endhead
\bottomrule\noalign{}
\endlastfoot
记账与现金流 & 自动分类、月度汇总 & 替你决定怎么花钱 \\
目标与规划 & 试算、情景模拟 & 替你设定人生目标 \\
信息研究 & 整理摘要、核对数据 & 替你判断该买什么 \\
组合构建 & 按清单核对步骤 & 替你选择具体标的 \\
交易执行 & 记录、提醒 & 替你下单 \\
复盘 & 整理记录、计算基准对比 & 替你评价决策对错 \\
\end{longtable}

这个表格想表达的核心是：\textbf{AI
在「过程」里是优秀的助手，在「决定」里必须退场。}

\subsection{数据安全与隐私，是投资场景的底线}\label{ux6570ux636eux5b89ux5168ux4e0eux9690ux79c1ux662fux6295ux8d44ux573aux666fux7684ux5e95ux7ebf}

比起记账，投资场景下的数据安全要求更高。这里有三条必须遵守的底线：

\textbf{第一，绝不把账户密码、身份证号、交易密码交给任何 AI 工具。}
没有任何合法工具需要这些信息来完成「帮你分析」这个任务。

\textbf{第二，确认数据的存储与使用边界。}
你的持仓、交易记录、盈亏数据，属于高度敏感的个人信息。在使用任何工具前，确认它不会把你的数据用于未经你同意的用途。

\textbf{第三，警惕「看起来太聪明」的工具。} 如果一个 AI
工具向你展示的持仓分析、收益预测精准得让你心动------请记住，展示效果越完美，越有可能是经过挑选的结果。真实世界的数据，永远是带噪声的。

\subsection{小练习：写下你的 AI
使用守则}\label{ux5c0fux7ec3ux4e60ux5199ux4e0bux4f60ux7684-ai-ux4f7fux7528ux5b88ux5219}

\begin{enumerate}
\def\labelenumi{\arabic{enumi}.}
\tightlist
\item
  用一句话写下：在投资中，我允许 AI
  帮我做什么（例如：整理信息、提醒纪律）；
\item
  用一句话写下：在投资中，我永远不会让 AI
  替我做（例如：选择具体标的、替我下单）；
\item
  把这两句话贴在你投资日志的第一页，和你的投资计划书放在一起。
\end{enumerate}

\begin{center}{\color{dh-bronze}\rule{0.2\textwidth}{0.5pt}}\end{center}

三章关于 AI 的内容写完了，但我们用 AI
的方式，不应该只停留在这三章里。回顾整本书------记账、规划、风险、配置、定投、复盘------AI
在每个环节都可以是一个更好的「过程助手」。而真正的决策权，永远在你手里。这也是这本书始终想传递的：工具在进化，但「你自己」才是投资中最重要、也最不可替代的变量。

\newpage

\chapter{终章 ·
财富的作用，是增加人生选择权}\label{ux7ec8ux7ae0-ux8d22ux5bccux7684ux4f5cux7528ux662fux589eux52a0ux4ebaux751fux9009ux62e9ux6743}

\begin{dhquote}
\textbf{投资从来不是目的，它是帮你抵达目的地的其中一条路。}
\end{dhquote}

\begin{center}{\color{dh-bronze}\rule{0.2\textwidth}{0.5pt}}\end{center}

在这本书里，我们讨论了复利、通胀、风险、资产配置、ETF、定投、再平衡、牛熊周期、行为金融，也讨论了如何用
AI
记账、做长期规划、守住投资边界------三十三个章节，涉及了大量数字、公式、规则和框架。但所有这些内容，服务的其实是一个很简单的问题。

\textbf{财富，到底有什么用？}

财富不是用来炫耀的，不是用来和别人比赛的，也不是一个需要用一生去刷新的纪录。财富最真实的用处，是给你选择权------选择什么时候退休、选择做什么样的工作、选择在家人需要的时候放下一切去陪伴，而不必担心下个月的房贷能不能交上。

很多人在追求财富的过程中，不知不觉把手段当成了目的。投资数字成了游戏分数，别人家的收益率成了焦虑来源，而最初那个''我想让生活变得更好''的愿望，反而在日复一日的盯盘中被遗忘了。如果你读完这本书之后，只做了一件事------不再把投资看得比生活本身更重要------那这本书就已经实现了它最重要的目标。

\subsection{一份完整的个人行动清单}\label{ux4e00ux4efdux5b8cux6574ux7684ux4e2aux4ebaux884cux52a8ux6e05ux5355}

如果你不知道从哪里开始，下面这二十个步骤可以帮助你从头到尾建立起自己的财富系统。不需要一天完成------有些步骤只需要几分钟，有些可以花几个星期慢慢思考。

\begin{enumerate}
\def\labelenumi{\arabic{enumi}.}
\tightlist
\item
  打开手机备忘录或拿出一张纸，写下你最重要的三个财务目标，每个附上一个大概的时间范围。
\item
  用第二章的模板，列出你的资产和负债，计算你的净资产------这是你的起点。
\item
  记录一个月的全部开支（每一笔都记，或使用第三十一章的 AI
  记账工具自动分类），看看你的钱真正去了哪里。算出你的储蓄率。
\item
  还清所有年利率超过 10\% 的高息负债。
\item
  开立一个应急资金账户，存入三到六个月的基本生活支出。
\item
  确认你已经有了基础的保险保障：医疗险、重疾险、意外险。如果家庭有依赖你收入的成员，加上定期寿险。
\item
  确定你每月可以从工资中固定拿出来长期投资的金额。它应该是一个你即使遇到小幅收入波动也能坚持下去的数字，而不是极限压缩。
\item
  完成一份风险承受能力自评（第八章的框架，或附录 C 的问卷）。
\item
  根据你的年龄、目标和风险评估，确定你的资产配置目标比例。写下它。
\item
  开立一个证券账户或基金账户。完成开户流程和风险测评。
\item
  完成你的第一笔投资。买一只宽基指数 ETF
  或场外指数基金。金额不重要，行动本身是最重要的。
\item
  设置每月定投计划（场内手动或场外自动）。实话说，最好是用场外指数基金让系统自动执行，减少人为干预的空间。
\item
  使用第十七章的 ETF 二十项筛单，选择你长期组合的核心工具。
\item
  按第二十四章的十二个步骤，建立你的完整投资组合书面方案。
\item
  写下每只持有资产的卖出条件。
\item
  创建一份投资日志，记录每一次买卖的理由。
\item
  按第二十五章的规则，设置你的再平衡方法------年度检查 + 偏离超出 ±5\%
  的阈值触发。
\item
  在日历上标记 12 个月后的''年度复盘日''。
\item
  关闭手机上的行情推送通知。把查看账户的频率控制为每月一次（检查定投执行情况加月度净值记录）。
\item
  继续你的生活。去工作、去陪伴家人、去学习新东西、去旅行。投资是你的副驾驶，不是你的全部。
\item
  如果你想借助 AI 的力量：用第三十一章的 AI
  记账工具记录现金流，用第三十二章的 AI
  规划工具做目标反推，并严格按照第三十三章的边界------AI
  当副驾驶，决策权永远在你手里。
\end{enumerate}

\subsection{从今天开始的第一步}\label{ux4eceux4ecaux5929ux5f00ux59cbux7684ux7b2cux4e00ux6b65}

也许你会觉得这二十一个步骤看起来很多，有点压力。没有关系。不需要做完二十一步才算是''开始了''。你今天只需要做一件事：写下你的三个最重要的财务目标。

不需要精确到元。不需要知道怎么实现。只需要诚实地写下来。因为任何一段漫长旅程，都是从知道自己要去哪里开始的。

\subsection{最后}\label{ux6700ux540e}

我希望你在多年之后重新翻开这本书的时候，会发现它记录的不仅仅是投资的方法，更是一种面对世界的态度。

市场会涨也会跌。经济有好也有坏。技术在不断进步，旧行业在被淘汰，新机会在不断出现。没有人能够准确预测未来会发生什么。但有一件事是确定的：那些持续学习、保持耐心、按照自己的节奏一步步往前走的人，最终会比那些不断追逐热点、被情绪牵着跑的人，走得更远、更踏实。

感谢你花了这么长时间和我一起走到这里。

写这本书的过程，也是我自己重新梳理投资理念的过程。有一些章节------比如''风险不只是价格下跌''、比如''投资中最难管理的是自己''------我写着写着，发现自己其实是在给自己上课。我希望这些内容对你有一点点帮助，哪怕只是让你在某个焦虑的时刻，想起书里的一句话，然后决定什么也不做。

愿你在漫长的人生里，财富和自由，兼而有之。

\subsection{延伸：作者的工具箱}\label{ux5ef6ux4f38ux4f5cux8005ux7684ux5de5ux5177ux7bb1}

这本书是免费的，但它背后有一些我每天都在用的工具。如果你在读某个章节时想过''要是有一个能自动算的就好了''，下面这几样也许能接上：

\begin{itemize}
\tightlist
\item
  \textbf{书中公式的网页工具}：复利、定投模拟、溢价率、费率对比、年龄配置、五年规则，加上
  AI
  记账对话包、目标翻译器、复盘生成器、投资红线卡------十个双语小工具，全部本地运行、不上传数据，随书在工具箱页面免费使用。
\item
  \textbf{蓝鲸记账（BlueWhale）}：第三十一章讲的''AI
  记账''，我日常的落地用法是这款自己开发的记账应用。它的标语只有一句话：记好当下，岁岁有余。
\item
  \textbf{AlphaForge}：第三十三章讲的''研究侧
  AI''，是我自研的开源研究工作台------把我做行业与公司研究时''收集信息、整理成结构化知识''的流程固化成软件，开源给同样想深入研究的读者。
\end{itemize}

有一点需要说明：蓝鲸记账是闭源产品，AlphaForge
是开源项目，两者都只代表我个人的工具选择，不是本书的推荐。这本书真正想留下的，始终是那套判断标准------研究侧与交易侧的边界、副驾驶与驾驶员的边界。愿这些工具帮你把书里的公式用起来，也帮你守住最后那句：\textbf{决策权，永远在你手里。}

\begin{center}{\color{dh-bronze}\rule{0.2\textwidth}{0.5pt}}\end{center}

\emph{花花}\\
\emph{2026 年，于「4i️e-acc 四倍做多认知」群}

\begin{center}{\color{dh-bronze}\rule{0.2\textwidth}{0.5pt}}\end{center}

\emph{本书所有投资案例均为教学假设，不构成任何投资建议。历史收益不代表未来表现。投资有风险，入市需谨慎。}

\backmatter

\newpage

\chapter{附录 A--N ·
模板与参考资料}\label{ux9644ux5f55-an-ux6a21ux677fux4e0eux53c2ux8003ux8d44ux6599}

\begin{dhquote}
以下模板供读者个人使用，可填写并复印。©本书教学用途。
\end{dhquote}

\begin{center}{\color{dh-bronze}\rule{0.2\textwidth}{0.5pt}}\end{center}

\section{附录 A ·
个人资产负债表}\label{ux9644ux5f55-a-ux4e2aux4ebaux8d44ux4ea7ux8d1fux503aux8868}

\begin{longtable}[]{@{}llll@{}}
\toprule\noalign{}\rowcolor{dh-pink!25}
资产 & 金额（元） & 负债 & 金额（元） \\
\midrule\noalign{}
\endhead
\bottomrule\noalign{}
\endlastfoot
现金和活期存款 & & 信用卡透支 & \\
货币基金 & & 消费贷款 & \\
定期存款 & & 汽车贷款 & \\
股票/基金/ETF & & 住房贷款余额 & \\
债券 & & 助学贷款 & \\
投资性房产（市值） & & 其他借款 & \\
自住房产（市值） & & & \\
退休账户/公积金 & & & \\
企业年金 & & & \\
其他资产 & & & \\
\textbf{资产总计} & & \textbf{负债总计} & \\
& & \textbf{净资产} & \\
\end{longtable}

填写日期：\_\_\_\_\_\_\_\_\_

\begin{center}{\color{dh-bronze}\rule{0.2\textwidth}{0.5pt}}\end{center}

\section{附录 B ·
月度家庭现金流表}\label{ux9644ux5f55-b-ux6708ux5ea6ux5bb6ux5eadux73b0ux91d1ux6d41ux8868}

\begin{longtable}[]{@{}ll@{}}
\toprule\noalign{}\rowcolor{dh-pink!25}
类别 & 金额（¥） \\
\midrule\noalign{}
\endhead
\bottomrule\noalign{}
\endlastfoot
\textbf{收入} & \\
工资（税后） & \\
经营/自雇收入 & \\
租金收入 & \\
利息与分红 & \\
其他收入 & \\
\textbf{总收入} & \\
& \\
\textbf{固定支出} & \\
住房（房贷/房租） & \\
水电物业 & \\
保险费 & \\
贷款还款（不含房贷） & \\
教育/学费 & \\
其他固定支出 & \\
& \\
\textbf{可变支出} & \\
餐饮 & \\
交通 & \\
购物 & \\
娱乐 & \\
医疗 & \\
其他可变支出 & \\
\textbf{总支出} & \\
& \\
\textbf{月度结余（收入 - 支出）} & \\
\textbf{储蓄率（结余 / 收入）} & \_\_\_\_\_\_\_\% \\
\end{longtable}

日期：\_\_\_\_\_\_\_\_\_

\begin{center}{\color{dh-bronze}\rule{0.2\textwidth}{0.5pt}}\end{center}

\section{附录 C ·
风险承受能力问卷}\label{ux9644ux5f55-c-ux98ceux9669ux627fux53d7ux80fdux529bux95eeux5377}

请根据实际情况，圈选最符合自己的选项。本问卷仅供自我评估参考。

\textbf{1. 您的投资期限大约是多少年？} A. 20年以上 B. 10-20年 C. 5-10年
D. 不到5年

\textbf{2. 您的收入稳定性如何？} A.
非常稳定（公务员/事业编制/大型国企等） B.
较为稳定（正规企业，有固定劳动合同） C.
不够稳定（自由职业者/佣金制/初创企业等） D. 极不稳定（收入来源不确定）

\textbf{3. 您的家庭现金流状况如何？} A.
每月有稳定的正向现金流，有较多结余 B. 每月有结余但比例不高 C.
收支基本平衡，几乎没有结余 D. 每月入不敷出

\textbf{4. 如果市场下跌20\%，您的第一反应是?} A.
继续按计划定投，甚至希望抓住机会 B. 有些不安但不做出大的改变 C.
想要卖出止损 D. 会立刻清仓

\textbf{5. 您目前持有的投资资产占您可投资资产的比例大约为?} A.
\textgreater70\% B. 40-70\% C. 10-40\% D. \textless10\%

\textbf{6. 您是否曾因为市场波动而卖出过投资资产？} A. 从来没有 B.
偶尔有过，但后来恢复了投资 C. 经常因为害怕亏损而卖出 D.
每次都受到波动的影响做出操作

\textbf{7.
您有多少个月的生活支出以现金或现金等价物存放（不含投资账户）？} A.
超过12个月 B. 6-12个月 C. 3-6个月 D. 不到3个月

\textbf{8. 您目前是否有高息负债（年利率\textgreater10\%）？} A. 完全没有
B. 有小额，正在还 C. 有较大金额 D. 有，且在持续增加

\textbf{9. 您能接受投资组合在一年内的最大临时损失幅度为：} A.
\textgreater30\% B. 20-30\% C. 10-20\% D. \textless10\%

\textbf{10. 您是否有家庭成员依赖于您的主要收入？} A.
没有（单身，无抚养人和需要赡养的人） B. 有一个 C. 有两个或更多 D.
整个家庭完全依赖我的收入

\begin{center}{\color{dh-bronze}\rule{0.2\textwidth}{0.5pt}}\end{center}

\section{附录 D · ETF
二十项筛选清单}\label{ux9644ux5f55-d-etf-ux4e8cux5341ux9879ux7b5bux9009ux6e05ux5355}

\begin{longtable}[]{@{}lll@{}}
\toprule\noalign{}\rowcolor{dh-pink!25}
\# & 检查项目 & 结果/备注 \\
\midrule\noalign{}
\endhead
\bottomrule\noalign{}
\endlastfoot
1 & 投资目标匹配度 & \\
2 & 跟踪指数（全称/代码） & \\
3 & 指数编制规则（加权方式） & \\
4 & 成分股数量 & \\
5 & 前十大持仓集中度 & \\
6 & 行业集中度（第一大行业占比） & \\
7 & 地区集中度 & \\
8 & 综合费率（管理费+托管费） & \\
9 & 跟踪差异（3年/5年） & \\
10 & 买卖价差（bid-ask spread） & \\
11 & 溢价折价历史 & \\
12 & 复制方式（完全/抽样/合成） & \\
13 & 分红政策（累积/分配） & \\
14 & 基金规模（AUM） & \\
15 & 日均成交量 & \\
16 & 底层资产流动性 & \\
17 & 税率影响（分红税等） & \\
18 & 汇率风险 & \\
19 & 杠杆或反向结构？ & \\
20 & 产品终止风险 & \\
\end{longtable}

\begin{dhquote}
\textbf{补充：高溢价买入风险速查表（配合第十五章）} --- 若以溢价率 p
买入 ETF，当溢价回归零（市价回到净值）时，你的实际亏损为 p ÷ (1 +
p)（亏损以你的买入成本为分母，而非净值）。高溢价通常出现在跨境、新发或涨停限购等场景，买入前请对照下表。
\end{dhquote}

\begin{longtable}[]{@{}ccc@{}}
\toprule\noalign{}\rowcolor{dh-pink!25}
买入时溢价率 & 溢价回归后的实际亏损 & 买 1 万元资产篮子的多付金额 \\
\midrule\noalign{}
\endhead
\bottomrule\noalign{}
\endlastfoot
0.5\% & ≈ 0.5\% & 50 元 \\
1\% & ≈ 0.99\% & 100 元 \\
2\% & ≈ 1.96\% & 200 元 \\
3\% & ≈ 2.91\% & 300 元 \\
5\% & ≈ 4.76\% & 500 元 \\
8\% & ≈ 7.41\% & 800 元 \\
10\% & ≈ 9.09\% & 1,000 元 \\
15\% & ≈ 13.04\% & 1,500 元 \\
20\% & ≈ 16.67\% & 2,000 元 \\
30\% & ≈ 23.08\% & 3,000 元 \\
\end{longtable}

（依据第十五章「溢价率 =（市价 − 净值）÷
净值」公式推导；历史溢价水平可查行情软件的折溢价/IOPV 指标。）

\begin{center}{\color{dh-bronze}\rule{0.2\textwidth}{0.5pt}}\end{center}

\section{附录 E ·
基金筛选清单}\label{ux9644ux5f55-e-ux57faux91d1ux7b5bux9009ux6e05ux5355}

\begin{longtable}[]{@{}lll@{}}
\toprule\noalign{}\rowcolor{dh-pink!25}
\# & 检查项目 & 结果/备注 \\
\midrule\noalign{}
\endhead
\bottomrule\noalign{}
\endlastfoot
1 & 基金类型（主动/被动/混合） & \\
2 & 基准指数（如有） & \\
3 & 基金经理（管理年限/历史业绩） & \\
4 & 管理费和托管费（总费率） & \\
5 & 申购费和赎回费 & \\
6 & 超额收益（Alpha vs 基准，3年/5年） & \\
7 & 风格漂移历史 & \\
8 & 最大回撤（历史） & \\
9 & 换手率 & \\
10 & 基金规模及变化趋势 & \\
11 & 前十大持仓及集中度 & \\
12 & 机构持有人比例 & \\
13 & 基金经理变动历史 & \\
14 & 比较同类基金的费率和业绩持续性 & \\
\end{longtable}

\begin{center}{\color{dh-bronze}\rule{0.2\textwidth}{0.5pt}}\end{center}

\section{附录 F ·
投资政策声明（IPS）模板}\label{ux9644ux5f55-f-ux6295ux8d44ux653fux7b56ux58f0ux660eipsux6a21ux677f}

\textbf{投资者：} \_\_\_\_\_\_\_\_\_ \textbf{日期：} \_\_\_\_\_\_\_\_\_
\textbf{下次审查日期：} \_\_\_\_\_\_\_\_\_

\textbf{投资目标：} \_\_\_\_\_\_\_\_\_（如：为退休积累资金，目标金额约 X
元，期限 Y 年）

\textbf{风险承受能力评估：} - 投资期限：\_\_\_\_\_\_\_\_\_ 年 -
收入稳定性：\_\_\_\_\_\_\_\_\_ -
最大可接受年度亏损：\_\_\_\_\_\_\_\_\_\%

\textbf{目标资产配置：} \textbar{} 资产类别 \textbar{} 目标比例
\textbar{} 允许偏离 \textbar{}
\textbar---------\textbar---------\textbar---------\textbar{} \textbar{}
宽基股票 ETF \textbar{} \_\_\_\_\% \textbar{} ±\_\_\_\_\% \textbar{}
\textbar{} 债券 ETF/基金 \textbar{} \_\_\_\_\% \textbar{} ±\_\_\_\_\%
\textbar{} \textbar{} 黄金 ETF \textbar{} \_\_\_\_\% \textbar{}
±\_\_\_\_\% \textbar{} \textbar{} 现金/货币类 \textbar{} \_\_\_\_\%
\textbar{} ±\_\_\_\_\% \textbar{}

\textbf{投入方式：} 每月定投 \_\_\_\_\_\_\_\_\_\_ 元 / 其他
\_\_\_\_\_\_\_\_\_

\textbf{再平衡规则：}
\_\_\_\_\_\_\_\_\_（如：每年12月检查，偏离超过±5\%触发）

\textbf{卖出条件（何时减仓或退出）：} \_\_\_\_\_\_\_\_\_

\textbf{约束条件（哪些产品/策略不会使用）：} \_\_\_\_\_\_\_\_\_

\textbf{审查频率：} 每年一次

\begin{center}{\color{dh-bronze}\rule{0.2\textwidth}{0.5pt}}\end{center}

\section{附录 G ·
定投计划模板}\label{ux9644ux5f55-g-ux5b9aux6295ux8ba1ux5212ux6a21ux677f}

\begin{longtable}[]{@{}ll@{}}
\toprule\noalign{}\rowcolor{dh-pink!25}
项目 & 内容 \\
\midrule\noalign{}
\endhead
\bottomrule\noalign{}
\endlastfoot
定投产品 & \\
定投金额（每期） & ¥\_\_\_\_\_\_\_\_ \\
定投频率 & 每月 / 每季度 / 其他：\_\_\_\_\_\_\_ \\
执行日期 & 每月 \_\_\_\_\_\_\_ 日 \\
投入方式 & 场内手动 / 场外自动 \\
执行平台 & \\
预计年数 & \_\_\_\_\_\_\_ 年 \\
暂停条件（什么情况下暂停） & \\
终止条件（什么情况下完全终止） & \\
年度审查日期 & 每年 \_\_\_\_\_\_\_ 月 \\
\end{longtable}

\begin{center}{\color{dh-bronze}\rule{0.2\textwidth}{0.5pt}}\end{center}

\section{附录 H ·
再平衡记录表}\label{ux9644ux5f55-h-ux518dux5e73ux8861ux8bb0ux5f55ux8868}

\begin{longtable}[]{@{}
  >{\raggedright\arraybackslash}p{(\linewidth - 12\tabcolsep) * \real{0.0857}}
  >{\raggedright\arraybackslash}p{(\linewidth - 12\tabcolsep) * \real{0.2000}}
  >{\raggedright\arraybackslash}p{(\linewidth - 12\tabcolsep) * \real{0.1286}}
  >{\raggedright\arraybackslash}p{(\linewidth - 12\tabcolsep) * \real{0.1286}}
  >{\raggedright\arraybackslash}p{(\linewidth - 12\tabcolsep) * \real{0.1857}}
  >{\raggedright\arraybackslash}p{(\linewidth - 12\tabcolsep) * \real{0.1857}}
  >{\raggedright\arraybackslash}p{(\linewidth - 12\tabcolsep) * \real{0.0857}}@{}}
\toprule\noalign{}\rowcolor{dh-pink!25}
\begin{minipage}[b]{\linewidth}\raggedright
日期
\end{minipage} & \begin{minipage}[b]{\linewidth}\raggedright
再平衡触发原因
\end{minipage} & \begin{minipage}[b]{\linewidth}\raggedright
卖出资产
\end{minipage} & \begin{minipage}[b]{\linewidth}\raggedright
买入资产
\end{minipage} & \begin{minipage}[b]{\linewidth}\raggedright
调整后股票实占
\end{minipage} & \begin{minipage}[b]{\linewidth}\raggedright
调整后债券实占
\end{minipage} & \begin{minipage}[b]{\linewidth}\raggedright
备注
\end{minipage} \\
\midrule\noalign{}
\endhead
\bottomrule\noalign{}
\endlastfoot
& 年度检查/阈值触发 & & & & & \\
& & & & & & \\
\end{longtable}

\begin{center}{\color{dh-bronze}\rule{0.2\textwidth}{0.5pt}}\end{center}

\section{附录 I ·
年度投资复盘表}\label{ux9644ux5f55-i-ux5e74ux5ea6ux6295ux8d44ux590dux76d8ux8868}

\textbf{复盘年度：} \_\_\_\_\_\_\_\_\_

\textbf{1. 组合表现} - 当年组合总回报：\_\_\_\_\_\_\_\% -
基准（自定义）回报：\_\_\_\_\_\_\_\% -
超额/落后基准幅度：\_\_\_\_\_\_\_\% - 主要原因：\_\_\_\_\_\_\_\_\_

\textbf{2. 配置检查} - 年初配置 vs 年末实际配置是否偏离？是/否 -
是否执行了再平衡？是/否。如未执行，原因是 \_\_\_\_\_\_\_\_\_

\textbf{3. 操作回顾} -
当年总共执行了多少笔交易（含定投）？\_\_\_\_\_\_\_\_\_ -
其中，完全按计划执行的交易有多少？\_\_\_\_\_\_\_\_\_ -
非计划交易（冲动买卖）的原因：\_\_\_\_\_\_\_\_\_

\textbf{4. 储蓄率} - 当年平均月收入（税后）：¥\_\_\_\_\_\_\_\_\_ -
当年平均月支出：¥\_\_\_\_\_\_\_\_\_ - 当年储蓄率：\_\_\_\_\_\_\_\%

\textbf{5. 需要改进的方面} \_\_\_\_\_\_\_\_\_

\textbf{6. 下一年度计划调整} \_\_\_\_\_\_\_\_\_

\begin{center}{\color{dh-bronze}\rule{0.2\textwidth}{0.5pt}}\end{center}

\section{附录 J ·
不同年龄阶段资产配置参考表}\label{ux9644ux5f55-j-ux4e0dux540cux5e74ux9f84ux9636ux6bb5ux8d44ux4ea7ux914dux7f6eux53c2ux8003ux8868}

\begin{dhquote}
以下比例为教学参考区间，因个人收入稳定性、负债、养老金保障、风险承受力等因素而异，不构成个性化建议。
\end{dhquote}

\begin{longtable}[]{@{}lllll@{}}
\toprule\noalign{}\rowcolor{dh-pink!25}
阶段 & 权益类（股票ETF） & 债券类 & 黄金/另类 & 现金 \\
\midrule\noalign{}
\endhead
\bottomrule\noalign{}
\endlastfoot
20-30岁 & 60-80\% & 15-30\% & 0-5\% & 5\% \\
30-40岁 & 50-70\% & 25-40\% & 5-10\% & 5\% \\
40-50岁 & 40-60\% & 30-50\% & 5-10\% & 5-10\% \\
50-60岁 & 30-50\% & 40-55\% & 5-10\% & 5-10\% \\
60岁以上 & 20-40\% & 40-60\% & 5-10\% & 10-15\% \\
\end{longtable}

\begin{dhquote}
\textbf{不适用情形速览（详见第二十二章）}：以下情形若与你相符，请相应调整区间，不要照搬上表------

\begin{itemize}
\tightlist
\item
  \textbf{20--30
  岁}：应急资金尚未备齐，或近两年内有明确要动用的大额资金（如购房首付）------先用低波动工具安置，不宜直接进入高权益比例。
\item
  \textbf{30--40
  岁}：三五年内计划换房的首付金，需在购房前从权益账户转出；创业、提成制等收入波动大的职业，权益比例宜靠向区间下限。
\item
  \textbf{40--50
  岁}：退休保障极强（如体制内养老）或已积累远超所需资产时，可以适当放宽权益比例；若已足够退休使用，也可主动转保守、追求稳健现金流。
\item
  \textbf{50--60
  岁}：退休金充裕到数倍覆盖生活所需时，多余资金可视作「继承人的投资期限」，保持较高权益也合理；若子女尚未经济独立或有持续大额支出，则应更保守。
\item
  \textbf{60
  岁以上}：退休金已完全覆盖生活、组合纯属锦上添花者，可容忍更高权益比例；健康状况不佳或有遗产规划需求者，宜更偏保守与流动。
\end{itemize}
\end{dhquote}

\begin{center}{\color{dh-bronze}\rule{0.2\textwidth}{0.5pt}}\end{center}

\section{附录 K ·
常见金融术语表}\label{ux9644ux5f55-k-ux5e38ux89c1ux91d1ux878dux672fux8bedux8868}

\begin{longtable}[]{@{}
  >{\raggedright\arraybackslash}p{(\linewidth - 4\tabcolsep) * \real{0.2500}}
  >{\raggedright\arraybackslash}p{(\linewidth - 4\tabcolsep) * \real{0.3750}}
  >{\raggedright\arraybackslash}p{(\linewidth - 4\tabcolsep) * \real{0.3750}}@{}}
\toprule\noalign{}\rowcolor{dh-pink!25}
\begin{minipage}[b]{\linewidth}\raggedright
中文
\end{minipage} & \begin{minipage}[b]{\linewidth}\raggedright
English
\end{minipage} & \begin{minipage}[b]{\linewidth}\raggedright
简要解释
\end{minipage} \\
\midrule\noalign{}
\endhead
\bottomrule\noalign{}
\endlastfoot
资产 & Asset & 能够持续创造价值、带来未来经济收益的东西 \\
负债 & Liability & 欠别人的钱，需要未来偿还 \\
净资产 & Net Worth & 资产减去负债 \\
流动资产 & Liquid Asset & 可以快速变现且不大幅折价的资产 \\
通货膨胀 & Inflation & 物价普遍上涨导致货币购买力下降 \\
复利 & Compounding & 收益也参与产生下一期收益 \\
单利 & Simple Interest & 只对本金计算收益 \\
72法则 & Rule of 72 & 72÷年化收益率≈本金翻倍年数 \\
名义收益率 & Nominal Return & 未经通胀调整的收益率 \\
实际收益率 & Real Return & 扣除通胀后的收益率 \\
波动 & Volatility & 资产价格上下起伏的幅度 \\
永久损失 & Permanent Loss of Capital &
资产基本面恶化导致的不可恢复损失 \\
股息 & Dividend & 企业分配给股东的利润 \\
市盈率 & Price-to-Earnings Ratio (P/E) & 每股价格÷每股盈利 \\
债券 & Bond & 借钱给政府或企业，对方按期付息并到期还本 \\
票面利率 & Coupon Rate & 债券发行时约定的年利息比率 \\
到期收益率 & Yield to Maturity (YTM) & 持有债券至到期日的年化总回报 \\
久期 & Duration & 债券价格对利率变化的敏感度指标 \\
信用利差 & Credit Spread & 企业债收益率高出同期限国债的部分 \\
ETF & Exchange-Traded Fund & 在交易所交易的指数基金 \\
净值 & Net Asset Value (NAV) & 基金资产减去负债后每份的价值 \\
溢价率 & Premium Rate & (市价-净值)÷净值 \\
折价率 & Discount Rate & (净值-市价)÷净值 \\
跟踪差异 & Tracking Difference & ETF收益与指数收益的实际差额 \\
授权参与商 & Authorized Participant (AP) &
可以在一级市场申购赎回ETF份额的大型金融机构 \\
定投 & Dollar-Cost Averaging (DCA) & 固定时间、固定金额持续投资 \\
资产配置 & Asset Allocation & 在不同资产类别之间分配资金 \\
再平衡 & Rebalancing & 将投资组合比例恢复到目标配置 \\
相关性 & Correlation & 不同资产价格之间的联动关系 \\
超额收益 & Alpha & 基金收益减去基准收益 \\
最大回撤 & Maximum Drawdown & 从高峰到低谷的最大跌幅 \\
QDII & Qualified Domestic Institutional Investor &
合格境内机构投资者（可投境外市场） \\
FOF & Fund of Funds & 投资于其他基金的基金 \\
下滑曲线 & Glide Path & 养老目标基金随目标日期临近逐步降低权益比例 \\
\end{longtable}

\begin{center}{\color{dh-bronze}\rule{0.2\textwidth}{0.5pt}}\end{center}

\section{附录 L ·
不同司法辖区需自行确认的事项}\label{ux9644ux5f55-l-ux4e0dux540cux53f8ux6cd5ux8f96ux533aux9700ux81eaux884cux786eux8ba4ux7684ux4e8bux9879}

以下事项在不同国家和地区存在显著差异。本书不提供具体法律或税务建议，读者应当根据自己所在地的最新法规自行确认。

\begin{longtable}[]{@{}
  >{\raggedright\arraybackslash}p{(\linewidth - 2\tabcolsep) * \real{0.5000}}
  >{\raggedright\arraybackslash}p{(\linewidth - 2\tabcolsep) * \real{0.5000}}@{}}
\toprule\noalign{}\rowcolor{dh-pink!25}
\begin{minipage}[b]{\linewidth}\raggedright
事项
\end{minipage} & \begin{minipage}[b]{\linewidth}\raggedright
说明
\end{minipage} \\
\midrule\noalign{}
\endhead
\bottomrule\noalign{}
\endlastfoot
证券账户开立条件 & 不同国家/地区对非居民开户的限制不同 \\
个人所得税 &
投资分红、基金赎回收益、海外投资收益的税率和申报要求因地区而异 \\
资本利得税 &
部分地区有、部分地区没有；长期持有和短期交易的税率可能不同 \\
预扣税（Withholding Tax） &
投资海外资产时，所在国可能对股息分红预先扣缴税款 \\
外汇管制与额度 &
部分国家有外汇管制，限制个人每年购汇或向境外转移资金的上限 \\
QDII额度限制 & 中国境内的QDII基金受外汇管理局批准的额度限制 \\
ETF税务处理 &
不同注册地（如中国大陆、香港、美国、爱尔兰）的ETF对分红的税务处理差异 \\
遗产税与继承 &
持有海外资产时，不同国家/地区对遗产税和继承程序的规定差异 \\
投资者保护基金 & 不同地区对券商破产时客户资产的保护范围和上限不同 \\
产品准入 & 部分ETF或基金因监管原因不向特定地区的居民销售 \\
信息披露语言 & 海外ETF的招募说明书和报告语言可能是英文或其他外语 \\
\end{longtable}

\begin{dhquote}
\textbf{重要提醒：以上事项的具体规定以各地区最新法规为准，并可能随时更新。本书不构成法律、税务或移民建议。}
\end{dhquote}

\begin{center}{\color{dh-bronze}\rule{0.2\textwidth}{0.5pt}}\end{center}

\section{附录 M ·
个人投资者税务自查清单}\label{ux9644ux5f55-m-ux4e2aux4ebaux6295ux8d44ux8005ux7a0eux52a1ux81eaux67e5ux6e05ux5355}

\begin{dhquote}
税收是投资成本的一部分------它不像佣金那么显眼，但会实实在在地侵蚀长期回报。幸运的是，中国大陆目前的规则对「拿得久的个人投资者」相当友好（见
M2）。本附录帮你建立「每年一次、十分钟」的税务自查习惯。以下以中国大陆规则为主（截至
2026 年 9 月），其他司法辖区请结合附录 L
自行确认。本附录不构成税务建议。
\end{dhquote}

\subsection{M1 ·
股息红利：持股越久，税负越低}\label{m1-ux80a1ux606fux7ea2ux5229ux6301ux80a1ux8d8aux4e45ux7a0eux8d1fux8d8aux4f4e}

个人持有 A 股上市公司股票，分红按持股期限差别化征税（财税〔2015〕101
号）：

\begin{longtable}[]{@{}lll@{}}
\toprule\noalign{}\rowcolor{dh-pink!25}
持股期限 & 计税方式 & 实际税负 \\
\midrule\noalign{}
\endhead
\bottomrule\noalign{}
\endlastfoot
1 个月以内（含） & 全额计入，按 20\% 税率 & 20\% \\
1 个月以上至 1 年（含） & 减按 50\% 计入，按 20\% 税率 & 10\% \\
超过 1 年 & 暂免征收 & 0 \\
\end{longtable}

\begin{itemize}
\tightlist
\item
  分红税在\textbf{卖出时}才结算（按先进先出计持股期）。频繁买卖不仅磨损手续费，还可能把「免税的长线分红」变成「20\%
  的短线分红」------这对定投者是个隐蔽的额外成本。
\item
  新三板挂牌公司适用同款差别化政策（现行公告延续至 2027-12-31）。
\end{itemize}

\subsection{M2 ·
买卖差价：目前的大类免税}\label{m2-ux4e70ux5356ux5deeux4ef7ux76eeux524dux7684ux5927ux7c7bux514dux7a0e}

\begin{longtable}[]{@{}ll@{}}
\toprule\noalign{}\rowcolor{dh-pink!25}
交易 & 现状（截至 2026-09） \\
\midrule\noalign{}
\endhead
\bottomrule\noalign{}
\endlastfoot
个人转让 A 股上市公司股票 & 差价暂免个税 \\
个人申购/赎回公募基金（含 ETF、场外基金） & 差价暂免个税 \\
内地个人通过港股通买卖港股 & 差价暂免个税（延续政策，注意到期确认） \\
\end{longtable}

\begin{dhquote}
「暂免」不等于「永久免」：政策措辞与延续时点请以财政部、税务总局最新公告为准。
\end{dhquote}

\subsection{M3 ·
分红与利息：藏在基金里的两个细节}\label{m3-ux5206ux7ea2ux4e0eux5229ux606fux85cfux5728ux57faux91d1ux91ccux7684ux4e24ux4e2aux7ec6ux8282}

\begin{itemize}
\tightlist
\item
  公募基金向个人分配收益：暂不征收个税。但如果基金底层持有\textbf{企业债券}，发债企业在付息时已按
  20\% 代扣个税------这笔税在基金层面已经扣过，不是在你卖出时才扣。
\item
  基金从上市公司取得的股息，在基金层面按差别化政策计征；通过\textbf{港股通}持有
  H 股等港股收到的股息，按 20\%
  代扣个税（沪港通/深港通机制下另有部分品种按协定税率）。
\end{itemize}

\subsection{M4 ·
跨境与申报提醒}\label{m4-ux8de8ux5883ux4e0eux7533ux62a5ux63d0ux9192}

\begin{longtable}[]{@{}
  >{\raggedright\arraybackslash}p{(\linewidth - 2\tabcolsep) * \real{0.5000}}
  >{\raggedright\arraybackslash}p{(\linewidth - 2\tabcolsep) * \real{0.5000}}@{}}
\toprule\noalign{}\rowcolor{dh-pink!25}
\begin{minipage}[b]{\linewidth}\raggedright
情形
\end{minipage} & \begin{minipage}[b]{\linewidth}\raggedright
提示
\end{minipage} \\
\midrule\noalign{}
\endhead
\bottomrule\noalign{}
\endlastfoot
QDII 基金 &
底层境外资产的预扣税由境外先行扣缴，境内税制衔接复杂，分红单据会注明 \\
海外股票/ETF &
所在国通常对股息预扣税（如美国对中国居民常见的预扣税率）；非居民资本利得一般不课税 \\
CRS &
中国已参与金融账户涉税信息自动交换；境外金融账户与境外所得存在申报义务，不要假设「查不到」 \\
外汇便利化额度 & 个人年度便利化购汇/结汇额度目前为等值 5
万美元/年，须合规用途使用 \\
\end{longtable}

\subsection{M5 · 年度自查行动表（建议在每年 3--6
月个税汇算时顺手完成）}\label{m5-ux5e74ux5ea6ux81eaux67e5ux884cux52a8ux8868ux5efaux8baeux5728ux6bcfux5e74-36-ux6708ux4e2aux7a0eux6c47ux7b97ux65f6ux987aux624bux5b8cux6210}

\begin{itemize}
\tightlist
\item
  翻看持仓：是否因短线交易把分红税从 0\% 变成
  10\%/20\%？（长期投资者通常可以做到 0\%）
\item
  核对基金对账单中的「代扣税」项目，确认没有异常扣税
\item
  如有境外账户或境外投资，确认当年度申报义务已履行
\item
  用几分钟检索「股息红利」「港股通个税」等关键词，确认持股优惠没有到期或调整
\end{itemize}

\begin{dhquote}
⚠️ 本附录基于截至 2026 年 9
月的公开规则撰写，具体适用以你所在地主管税务机关口径为准。涉及大额或跨境事项，请咨询专业税务师。
\end{dhquote}

\begin{center}{\color{dh-bronze}\rule{0.2\textwidth}{0.5pt}}\end{center}

\section{附录 N ·
遗产与传承规划清单}\label{ux9644ux5f55-n-ux9057ux4ea7ux4e0eux4f20ux627fux89c4ux5212ux6e05ux5355}

\begin{dhquote}
很少有人愿意想「自己不在之后，资产怎么办」------但这是长期投资者迟早要面对的问题，尤其是当你按这本书的方法积累多年之后。先给结论：中国\textbf{目前没有遗产税}，财政部、国家税务总局在
2026 年 5 月还专门辟谣过「2026
年开征遗产税」的网络传言。但「没有遗产税」不等于「继承零成本」：公证、过户、契税与家庭纠纷都可能发生。而且立法层面已出现变化信号------2025
年 11
月，全国人大财经委将「遗产（和赠与）税法」列入加强调研论证的立法项目。方向很清楚：趁现在没有税，先把资产梳理清楚、把安排做对。本附录为教学用途，不构成法律意见。
\end{dhquote}

\subsection{N1 ·
资产清点表（建议每年更新一次）}\label{n1-ux8d44ux4ea7ux6e05ux70b9ux8868ux5efaux8baeux6bcfux5e74ux66f4ux65b0ux4e00ux6b21}

\begin{longtable}[]{@{}llll@{}}
\toprule\noalign{}\rowcolor{dh-pink!25}
资产类别 & 存放机构/平台 & 账户/凭证号 & 备注 \\
\midrule\noalign{}
\endhead
\bottomrule\noalign{}
\endlastfoot
银行存款 & & & \\
股票账户（A 股/港股通） & & & \\
基金账户 & & & \\
房产 & & & 权属/贷款余额 \\
人寿保险 & & & 保单号/受益人 \\
企业股权/合伙份额 & & & \\
未结债权 & & & \\
数字资产（加密钱包等） & & & 密钥/助记词存放 \\
贵重实物（贵金属、收藏） & & & \\
债务（房贷、借款） & & & \\
\end{longtable}

\begin{itemize}
\tightlist
\item
  清点表的价值在于\textbf{有人能找到它}：把存放位置告诉一位信任的人，或随遗嘱一并交代，比锁在只有你知道的地方更有用。
\end{itemize}

\subsection{N2 ·
中国大陆继承规则速查（据《民法典》继承编）}\label{n2-ux4e2dux56fdux5927ux9646ux7ee7ux627fux89c4ux5219ux901fux67e5ux636eux6c11ux6cd5ux5178ux7ee7ux627fux7f16}

\begin{longtable}[]{@{}
  >{\raggedright\arraybackslash}p{(\linewidth - 2\tabcolsep) * \real{0.5000}}
  >{\raggedright\arraybackslash}p{(\linewidth - 2\tabcolsep) * \real{0.5000}}@{}}
\toprule\noalign{}\rowcolor{dh-pink!25}
\begin{minipage}[b]{\linewidth}\raggedright
项目
\end{minipage} & \begin{minipage}[b]{\linewidth}\raggedright
要点
\end{minipage} \\
\midrule\noalign{}
\endhead
\bottomrule\noalign{}
\endlastfoot
法定继承顺序 &
第一顺序：配偶、子女、父母；第二顺序：兄弟姐妹、祖父母、外祖父母。有第一顺序继承人时由第一顺序继承，无第一顺序时才由第二顺序继承 \\
遗嘱形式 & 自书、代书、打印、录音录像、口头、公证共 6 种（民法典第
1134--1139 条）；口头遗嘱须两名以上见证人，危急情况消除后失效 \\
多份遗嘱冲突 &
内容相抵触的，以\textbf{最后}订立的遗嘱为准（公证遗嘱不再当然优先） \\
房产继承过户 & 法定继承免征契税；遗赠/赠与等非法定继承按 3\%--5\%
缴纳契税（各地不同）；实务中办理过户通常需继承权公证（费用各地标准不一） \\
保险金 & 已指定身故受益人的人寿保险金\textbf{不属于遗产}（《保险法》第
42 条），直接给付受益人，通常不进入继承程序 \\
债务 &
限定继承：继承人仅以所得遗产的实际价值为限，清偿被继承人的税款与债务 \\
\end{longtable}

\subsection{N3 ·
传承工具怎么选（教学速览）}\label{n3-ux4f20ux627fux5de5ux5177ux600eux4e48ux9009ux6559ux5b66ux901fux89c8}

\begin{longtable}[]{@{}
  >{\raggedright\arraybackslash}p{(\linewidth - 4\tabcolsep) * \real{0.3333}}
  >{\raggedright\arraybackslash}p{(\linewidth - 4\tabcolsep) * \real{0.3333}}
  >{\raggedright\arraybackslash}p{(\linewidth - 4\tabcolsep) * \real{0.3333}}@{}}
\toprule\noalign{}\rowcolor{dh-pink!25}
\begin{minipage}[b]{\linewidth}\raggedright
工具
\end{minipage} & \begin{minipage}[b]{\linewidth}\raggedright
适合谁
\end{minipage} & \begin{minipage}[b]{\linewidth}\raggedright
主要特点
\end{minipage} \\
\midrule\noalign{}
\endhead
\bottomrule\noalign{}
\endlastfoot
遗嘱 & 绝大多数人 & 成本低、可随时修改；须按法定形式订立并妥善保管 \\
人寿保险（指定受益人） & 有家庭责任者 &
身故金定向给付、不进入遗产程序、私密性较好 \\
生前赠与 & 想提前转移资产者 &
过户即完成转移，但涉及契税等过户成本，且过早转移可能失去对资产的控制 \\
家族信托/保险金信托 & 资产量大、诉求复杂者 &
门槛与费用较高，需专业机构与完整的法律架构 \\
\end{longtable}

\begin{dhquote}
本书不推销任何金融产品。上表中的保险与信托只是「工具选项」，请按真实需求与成本决定，而非听信营销话术。
\end{dhquote}

\subsection{N4 · 行动清单}\label{n4-ux884cux52a8ux6e05ux5355}

\begin{itemize}
\tightlist
\item
  整理 N1 资产清点表，并告知一位信任的人存放位置
\item
  检查所有人寿保单/养老账户的「身故受益人」是否为当前意愿
\item
  评估是否需要订立遗嘱；若已订立，确认形式合法、家人知道存放处
\item
  如有境外资产，专门了解当地继承与遗产税规则（各国差异可能巨大）
\item
  将「遗产税立法进展」加入年度税务自查的关注清单（见附录 M5）
\end{itemize}

\begin{dhquote}
⚠️
本附录为通用教学信息，不构成法律、税务或移民意见。跨境继承、大额资产或复杂家庭结构，请咨询专业律师。
\end{dhquote}

\end{document}
